% informe-final-v1.0.0.tex
%
% Documento maestro del Informe Academico Final de la Entrega Final (v1.0.0)
% de BIOPET. Clase: report. Idioma: espanol (babel). Codificacion: UTF-8.
%
% Compilacion (desde docs/informe/):
%   pdflatex -interaction=nonstopmode informe-final-v1.0.0.tex
%   bibtex informe-final-v1.0.0
%   pdflatex -interaction=nonstopmode informe-final-v1.0.0.tex
%   pdflatex -interaction=nonstopmode informe-final-v1.0.0.tex
% o, si esta disponible:
%   latexmk -pdf -interaction=nonstopmode -halt-on-error informe-final-v1.0.0.tex
%
% Este documento es la evolucion del informe de la Tercera Entrega
% (informe-entrega-3.tex, que se conserva sin modificar como registro
% historico de esa entrega). No reescribe la historia: donde una cifra
% cambio entre la Tercera Entrega y la Entrega Final (numero de pruebas,
% umbral de cobertura, SUS, Lighthouse, seguridad dinamica/estatica,
% despliegue, DOI), este documento reporta el valor REAL vigente a la
% fecha de cierre de la Entrega Final, citando su fuente en el propio
% repositorio.
%
% Las figuras de evidencia siguen protegidas con \IfFileExists (comando
% \evidencia, definido abajo): la ausencia de una captura no deja huecos
% visibles ni texto de marcador en el PDF.

\documentclass[12pt,a4paper]{report}

% --- Idioma y codificacion --------------------------------------------
\usepackage[utf8]{inputenc}
\usepackage[T1]{fontenc}
\usepackage{lmodern}
\usepackage[spanish]{babel}

% --- Geometria de pagina -------------------------------------------------
\usepackage[a4paper,margin=2.5cm]{geometry}

% --- Manejo seguro de guion bajo en modo texto --------------------------
\usepackage{underscore}

% --- Tablas, figuras, listados --------------------------------------------
\usepackage{graphicx}
\usepackage{float}
\usepackage{booktabs}
\usepackage{longtable}
\usepackage{array}
\usepackage{caption}
\usepackage{enumitem}
\usepackage{textcomp}
\usepackage{amsmath}
\usepackage{amssymb}
\usepackage{pdflscape}
\usepackage{xcolor}
% microtype: mejora la justificacion (expansion/compresion de fuente,
% protrusion de margen) y reduce sustancialmente los "Overfull \hbox"
% causados por rutas de archivo largas en \texttt{} dentro de parrafos y
% captions, sin cambiar el contenido visible del texto.
\usepackage{microtype}
% emergencystretch: colchon de estiramiento de ultimo recurso que TeX solo
% usa cuando un parrafo no tiene ningun corte de linea que cumpla el
% \tolerance normal (tipico con \texttt{} largo: rutas, timestamps ISO,
% identificadores). No es \sloppy (que relaja \tolerance/\hyphenpenalty en
% TODO el documento y empeora el espaciado en parrafos que ya estaban
% bien); aqui solo se activa como respaldo puntual, sin alterar la
% justificacion normal de los parrafos que ya encajan.
\setlength{\emergencystretch}{3em}
% listings: para los listados de codigo cortos exigidos por la Guia Final
% (seccion Implementacion), con indice de listados en la materia
% preliminar (\lstlistoflistings).
\usepackage{listings}
\lstset{
  basicstyle=\ttfamily\footnotesize,
  breaklines=true,
  breakatwhitespace=false,
  columns=fullflexible,
  frame=single,
  numbers=left,
  numberstyle=\tiny,
  captionpos=b,
  language=Java,
  showstringspaces=false,
  tabsize=2
}
% \flushbottom (por defecto en la clase report) fuerza que cada pagina
% llene exactamente el alto de texto, estirando el espacio entre
% parrafos/tablas; con contenido denso (tablas, listas) eso puede producir
% un "Overfull \vbox" puntual y, como efecto colateral del algoritmo de
% salida, una pagina practicamente en blanco justo antes del capitulo
% siguiente. \raggedbottom permite que cada pagina termine a una altura
% natural en vez de estirarse, eliminando ese efecto.
\raggedbottom
% \rutalarga{a/b/c.ext}: version de \texttt{} para rutas de archivo largas,
% que inserta puntos de corte de linea (\allowbreak) despues de cada "/",
% en vez de dejar que "/" (que TeX no trata como punto de corte por
% defecto) produzca una sola palabra indivisible que desborda el margen.
\newcommand{\rutalarga}[1]{\texttt{\rutalargaimpl#1/\relax}}
\def\rutalargaimpl#1/#2\relax{%
  \ifx\relax#2\relax #1%
  \else #1/\allowbreak\rutalargaimpl#2\relax\fi
}
% El capitulo de Anexos tiene mas de nueve \section (Glosario, Endpoints,
% Documentos de referencia, Clases de prueba, Anexo A al G, Evidencias
% pendientes), asi que su numeracion en el indice llega a dos digitos
% (p. ej. "17.10."); el ancho reservado por defecto para el numero en
% \l@section (2.3em) alcanza para un digito y desborda por un par de
% puntos con dos. Se ensancha ligeramente esa caja, sin tocar ningun otro
% aspecto del indice.
\makeatletter
\renewcommand*\l@section{\@dottedtocline{1}{1.5em}{2.8em}}
\makeatother

% --- Encabezados/pies de pagina -------------------------------------------
\usepackage{fancyhdr}
\setlength{\headheight}{15.4pt}
% Encabezado de capitulo mas corto (solo numero + titulo, sin la palabra
% "Capitulo" en mayusculas ni el nombre del proyecto repetido), para que
% quepa con margen frente al bloque izquierdo sin que ambos campos se
% peguen visualmente en una sola linea de texto corrido.
\renewcommand{\chaptermark}[1]{\markboth{\thechapter.\ #1}{}}
\pagestyle{fancy}
\fancyhf{}
\fancyhead[L]{\small BIOPET v1.0.0}
\fancyhead[R]{\small\leftmark}
\fancyfoot[C]{\thepage}
\renewcommand{\headrulewidth}{0.4pt}

% --- Enlaces internos ------------------------------------------------------
\usepackage[hidelinks]{hyperref}
\hypersetup{
  pdftitle={BIOPET -- Sistema web de gestion veterinaria: Informe academico final v1.0.0},
  pdfauthor={Mariscal Cabrera Jaime Josue; Beltran Montiel Fred Adrian; Taipe Mora Zaida Melissa},
  pdfsubject={Informe academico final -- Entrega Final v1.0.0},
  bookmarksnumbered=true,
  breaklinks=true,
  colorlinks=false
}

% --- Bibliografia estilo IEEE (numerico), via BibTeX clasico -------------
% (usa la misma referencias.bib de docs/informe/, sin duplicarla)

% --- Comando auxiliar para evidencias condicionales -----------------------
\newcommand{\evidencia}[3]{%
  \IfFileExists{#1}{%
    \begin{figure}[H]
      \centering
      \includegraphics[width=0.82\textwidth]{#1}
      \caption{#2}
      \label{#3}
    \end{figure}
  }{}%
}

% --- Metadatos institucionales (usados en la portada) ---------------------
\newcommand{\universidad}{Universidad T\'ecnica Estatal de Quevedo}
\newcommand{\facultad}{Facultad de Ciencias de la Computaci\'on}
\newcommand{\carrera}{Carrera de Software}
\newcommand{\asignatura}{Aplicaciones Web}
\newcommand{\docente}{Dr.\ Gleiston Cicer\'on Guerrero Ulloa, Ph.D.}
\newcommand{\correodocente}{gguerrero@uteq.edu.ec}
\newcommand{\periodoacademico}{Presencial 2026 -- 2027}
\newcommand{\cursoparalelo}{A}
\newcommand{\equipo}{BMT}
\newcommand{\tituloinforme}{BIOPET --- Sistema web de gesti\'on veterinaria:\\ Informe acad\'emico final --- Entrega Final v1.0.0}
\newcommand{\ciudadfecha}{Quevedo, 18 de agosto de 2026}
% Los \allowbreak despues de cada "/" son puntos de corte de linea
% invisibles (no imprimen guion ni cambian el valor); evitan que estos
% identificadores largos, usados dentro de \texttt{}, produzcan un
% "Overfull \hbox" cuando aparecen en una columna estrecha o junto a mas
% texto en la misma linea.
\newcommand{\doisoftware}{10.5281/\allowbreak zenodo.21988746}
\newcommand{\doidataset}{10.5281/\allowbreak zenodo.21988785}
\newcommand{\urlfrontend}{https://\allowbreak biopet-frontend.onrender.com}
\newcommand{\urlbackend}{https://\allowbreak biopet-backend-dh5e.onrender.com}
\newcommand{\urlrepositorio}{https://\allowbreak github.com/\allowbreak Grinjoww/\allowbreak ENTREGA-FINAL-BIOPET}
% ORCID verificados de los tres integrantes (digito de control valido).
\newcommand{\orcidjaime}{0009-0007-2206-3941}
\newcommand{\orcidfred}{0009-0007-8303-2137}
\newcommand{\orcidzaida}{0009-0000-1227-3258}
% \commithashcorto: hash corto del commit AL QUE APUNTA el tag historico
% v1.0.0 (verificado con "git rev-list -n 1 v1.0.0"), NO el HEAD actual
% de esta rama de recalificacion. El tag v1.0.0 ya existe y es inmutable;
% este valor identifica ese estado historico de cierre de la Entrega
% Final, distinto de los commits posteriores de correccion/recalificacion
% que se agregan encima sin mover el tag. La fila de la portada rotula
% explicitamente el valor como "Tag v1.0.0 / commit" para que no se
% confunda con el HEAD vigente.
\newcommand{\commithashcorto}{0d5cd52}

\begin{document}
% Desactiva el atajo "-shorthand" del modulo spanish de babel (usa """
% como caracter activo para digrafos tipograficos): el informe usa comillas
% dobles literales dentro de \texttt{} (JSON, atributos HTML) que de otro
% modo se interpretarian como ese atajo y romperian la compilacion.
\shorthandoff{"}

% ===========================================================================
% PORTADA
% ===========================================================================
\begin{titlepage}
  \centering
  \vspace*{0.3cm}

  {\large \universidad \par}
  {\large \facultad \par}
  \vspace{0.3cm}
  {\normalsize \carrera \par}
  {\normalsize Asignatura: \asignatura \par}

  \vspace{1.6cm}
  {\Huge \bfseries BIOPET \par}
  \vspace{0.4cm}
  {\Large Sistema web de gesti\'on veterinaria \par}
  \vspace{0.6cm}
  {\Large\itshape Informe acad\'emico final --- Entrega Final v1.0.0 \par}

  \vspace{1.6cm}
  {\large Equipo \textbf{\equipo} \par}
  \vspace{0.6cm}

  {\normalsize
    \begin{tabular}{>{\centering\arraybackslash}p{15cm}}
      Mariscal Cabrera Jaime Josu\'e \quad --- \quad jmariscalc@\allowbreak uteq.edu.ec \quad --- \quad ORCID: \orcidjaime \\
      Beltr\'an Montiel Fred Adri\'an \quad --- \quad fbeltranm@\allowbreak uteq.edu.ec \quad --- \quad ORCID: \orcidfred \\
      Taipe Mora Zaida Melissa \quad --- \quad ztaipem@\allowbreak uteq.edu.ec \quad --- \quad ORCID: \orcidzaida \\
    \end{tabular}
    \par
  }

  \vspace{1.4cm}
  {\normalsize Docente: \docente \par}
  {\normalsize \texttt{\correodocente} \par}
  \vspace{0.3cm}
  {\normalsize Per\'iodo acad\'emico: \periodoacademico \quad --- \quad Curso/paralelo: \cursoparalelo \par}

  \vspace{1.1cm}
  {\small
    \begin{tabular}{@{}>{\raggedright\arraybackslash}p{6.9cm}>{\raggedright\arraybackslash}p{8.4cm}@{}}
      DOI del software (Zenodo) & \texttt{\doisoftware} \\
      DOI del \emph{dataset} de evidencias (Zenodo) & \texttt{\doidataset} \\
      Repositorio & \texttt{\urlrepositorio} \\
      Imagen de contenedor (GHCR, digest inmutable) & \texttt{sha256:ef1e857a\ldots6bef5163} \\
      Frontend (producci\'on) & \texttt{\urlfrontend} \\
      Backend (producci\'on) & \texttt{\urlbackend} \\
      \ifx\commithashcorto\empty\else
      Tag v1.0.0 / commit & \texttt{\commithashcorto} \\
      \fi
    \end{tabular}
    \par
  }

  \vfill

  {\normalsize \ciudadfecha \par}
  \vspace{0.5cm}

\end{titlepage}

% ===========================================================================
% MATERIA PRELIMINAR
% ===========================================================================
\pagenumbering{roman}

\tableofcontents
\clearpage
\listoffigures
\clearpage
\listoftables
\clearpage
\renewcommand{\lstlistlistingname}{\'Indice de listados}
\lstlistoflistings
\clearpage

\pagenumbering{arabic}
\setcounter{page}{1}

% ===========================================================================
% CAPITULOS
% ===========================================================================
\chapter*{Resumen}
\addcontentsline{toc}{chapter}{Resumen}

\textbf{Problema.} La mayor\'ia de cl\'inicas veterinarias en Ecuador
operan con registros manuales o en papel: p\'erdida de expedientes,
dif\'icil consulta del historial cl\'inico y falta de indicadores de
gesti\'on~\cite{cedeno2021veterinaria}. \textbf{Soluci\'on.} BIOPET es un
sistema web de gesti\'on veterinaria (Spring~Boot, Angular, PostgreSQL,
Redis/Valkey) con autenticaci\'on por rol/propiedad, CRUD de mascotas,
citas, consultas, vacunas y acceso a rutinas almacenadas.
\textbf{Metodolog\'ia.} \emph{Engineering Research}/\emph{Design
Science} (Peffers et al.~\cite{peffers2007dsrm}) con marco
Goal--Question--Metric~\cite{vansolingen2002gqm}, sin dise\~no
experimental controlado. \textbf{Principales validaciones.} Corte
hist\'orico v1.0.0: 205 pruebas automatizadas (0 fallos, 0 errores) y
cobertura JaCoCo de 91{,}80\,\% en l\'ineas y 79{,}39\,\% en ramas
(umbral 70\,\%); rendimiento con k6
(0{,}0\,\% de error); usabilidad SUS sobre 18 participantes (media
74{,}44/100, IC~95\,\% $[63{,}33, 85{,}56]$); seguridad (OWASP
Top~10 manual; SpotBugs est\'atico y auditor\'ia de SQL din\'amico sin
hallazgos; ZAP: \emph{baseline} hist\'orico y escaneo autenticado sin
alertas altas); y despliegue en producci\'on (Render,
HTTPS activo). Software, evidencias y backend archivados con DOI
(\texttt{\doisoftware}/\texttt{\doidataset}, Zenodo) y \emph{digest}
inmutable (GHCR). \textbf{Conclusi\'on
principal.} BIOPET demuestra, con evidencia reproducible
(ISO/IEC~25010~\cite{iso25010-2011}), que es posible construir y evaluar
emp\'iricamente, en un Proyecto Fin de Curso,
un artefacto funcional, seguro seg\'un los controles auditados y
reproducible; las limitaciones (entorno acad\'emico, tama\~no muestral)
se documentan en el cap\'itulo de amenazas a la validez.

\textbf{Palabras clave:} gesti\'on veterinaria; aplicaciones web;
Spring~Boot; Angular; JWT; PostgreSQL; ISO/IEC~25010; usabilidad (SUS);
seguridad OWASP; reproducibilidad.

\bigskip
\hrule
\bigskip

\chapter*{Abstract}
\addcontentsline{toc}{chapter}{Abstract}

\textbf{Problem.} Most veterinary clinics in Ecuador and similar contexts
still rely on manual or paper-based records: lost files, difficult
clinical-history retrieval, and a lack of management
indicators~\cite{cedeno2021veterinaria}. \textbf{Solution.} BIOPET is a
web-based veterinary management system (Spring~Boot, Angular,
PostgreSQL, Redis/Valkey) with role/ownership-based authentication, pet
CRUD, appointments, consultations, vaccines, and access to stored
routines. \textbf{Methodology.} Engineering Research/Design Science
(Peffers et al.~\cite{peffers2007dsrm}) with a Goal--Question--Metric
framework~\cite{vansolingen2002gqm}, no controlled experimental design.
\textbf{Main validations.} In the historical v1.0.0 snapshot, 205
automated tests were recorded (0 failures, 0 errors), with JaCoCo line
coverage of 91.80\,\% and branch coverage of 79.39\,\% (70\,\% gate); k6
performance testing (0.0\,\% error rate); usability measured
with SUS on 18 participants (mean 74.44/100, 95\,\% CI $[63.33, 85.56]$);
security audited (manual OWASP Top~10; SpotBugs static analysis and a
dynamic-SQL audit of the stored routines, both with no findings; dynamic
OWASP~ZAP: a historical baseline scan and an authenticated scan on the
local TLS environment, both with no high-severity alerts); and
production deployment (Render, HTTPS enabled). Software, evidence and
backend image remain archived with a DOI
(\texttt{\doisoftware}/\texttt{\doidataset}, Zenodo) and an immutable
\emph{digest} (GHCR). \textbf{Main conclusion.}
BIOPET shows, with reproducible evidence (ISO/IEC~25010~\cite{iso25010-2011}),
that it is possible to build and empirically evaluate, within a capstone
course project, a functional artifact that is secure with respect to the
audited controls and reproducible; limitations (academic environment,
sample size) are documented in the threats-to-validity chapter.

\textbf{Keywords:} veterinary management; web applications;
Spring~Boot; Angular; JWT; PostgreSQL; ISO/IEC~25010; usability (SUS);
OWASP security; reproducibility.

\chapter{Introducci\'on}
\label{cap:introduccion}

\section{Contexto: gesti\'on veterinaria y digitalizaci\'on incompleta}

La administraci\'on de una cl\'inica veterinaria involucra procesos que, en
ausencia de una herramienta digital centralizada, dependen de registros
manuales o en hojas de c\'alculo dispersas: identificaci\'on de mascotas y
due\~nos, historial de atenciones, control de vacunas y programaci\'on de
citas. La literatura reciente sobre el contexto ecuatoriano confirma que
esta carencia no es an\'ecdota: Cede\~no Ochoa, Catuto Murillo \& Rodas-Silva
documentan expl\'icitamente c\'omo la mayor\'ia de cl\'inicas veterinarias
locales carecen de herramientas tecnol\'ogicas de gesti\'on, lo que provoca
p\'erdida y duplicidad de expedientes por manejo manual~\cite{cedeno2021veterinaria}.
Rodr\'iguez et al. documentan un patr\'on similar en otro contexto
geogr\'afico, motivando un sistema de c\'odigo abierto para el mismo
dominio~\cite{rodriguez2024oscrum}. BIOPET nace de ese mismo diagn\'ostico:
la necesidad de un sistema web que centralice identificaci\'on,
autenticaci\'on, y la gesti\'on de mascotas, citas, consultas y vacunas de
una cl\'inica veterinaria, con evidencia emp\'irica de que efectivamente
funciona y es razonablemente seguro, no solo de que fue programado.

\section{Problema abordado}

El problema que BIOPET aborda tiene dos capas. La primera es funcional:
proveer una plataforma web que permita registrar usuarios con distintos
roles (administrador, veterinario, auxiliar, due\~no), administrar mascotas
asociadas a un due\~no, y gestionar el ciclo de atenci\'on veterinaria
(citas, consultas, vacunas) con control de acceso adecuado a cada rol. La
segunda capa, menos visible pero igual de cr\'itica, es de \textbf{evidencia
emp\'irica}: un sistema acad\'emico puede afirmar que ``funciona'' sin que
esa afirmaci\'on est\'e respaldada por pruebas automatizadas, medici\'on de
rendimiento bajo carga, auditor\'ia de seguridad real o percepci\'on de
usabilidad medida con un instrumento estandarizado. BIOPET trata esa
segunda capa como parte del problema mismo a resolver, no como un anexo
opcional: cada afirmaci\'on cuantitativa de este informe cita su fuente de
evidencia real en el repositorio.

\section{Motivaci\'on}

Tres entregas previas (Entrega~1A, Entrega~1B y Tercera Entrega) generaron
retroalimentaci\'on docente real, registrada \'integramente en
\texttt{docs/observaciones/OBSERVACIONES.md}: 15 observaciones formales,
todas cerradas a la fecha de este informe, cubriendo desde la ausencia de un
repositorio accesible hasta la falta de evidencia Lighthouse o de un DOI de
archivado permanente. Esa bitácora de correcciones es, en s\'i misma, parte
de la motivaci\'on de este documento: la Entrega Final no reescribe el
proyecto desde cero, sino que consolida el trabajo ya evaluado, corrigiendo
expl\'icitamente lo se\~nalado y documentando con honestidad lo que sigue
pendiente.

\section{Preguntas de investigaci\'on}
\label{sec:preguntas-investigacion}

Este informe organiza su evaluaci\'on emp\'irica alrededor de tres
preguntas de investigaci\'on, cada una medible con evidencia real ya
recolectada y respondida expl\'icitamente en el
cap\'itulo~\ref{cap:discusion} (secci\'on~\ref{sec:respuestas-rq}):

\begin{description}[nosep]
  \item[RQ1.] \textquestiondown En qu\'e medida BIOPET satisface los
    requisitos funcionales y no funcionales definidos para la Entrega
    Final, y cu\'ales permanecen expl\'icitamente pendientes?
  \item[RQ2.] \textquestiondown Qu\'e resultados presenta BIOPET en
    cobertura de pruebas, rendimiento, usabilidad, calidad web y seguridad
    bajo las evaluaciones emp\'iricas realmente ejecutadas?
  \item[RQ3.] \textquestiondown En qu\'e medida el proceso de despliegue,
    integraci\'on continua, publicaci\'on de la imagen de contenedor y
    archivado permanente permiten reproducir y verificar el artefacto
    final de forma independiente?
\end{description}

\section{Objetivo general}

Consolidar, evaluar emp\'iricamente y documentar acad\'emicamente el sistema
BIOPET en su versi\'on v1.0.0, integrando evidencia real de calidad
funcional, seguridad, rendimiento, usabilidad, mantenibilidad y
reproducibilidad, enmarcada en el modelo de calidad ISO/IEC~25010 y en una
metodolog\'ia de investigaci\'on de ingenier\'ia expl\'icita, sin inventar
resultados ni ocultar limitaciones.

\subsection*{Objetivos espec\'ificos}
\begin{enumerate}[nosep]
  \item Verificar el cumplimiento de los requisitos funcionales y no
    funcionales declarados en el SRS y documentar la arquitectura del
    sistema (incluidos los m\'odulos de Unidad~IV), distinguiendo
    expl\'icitamente lo verificado de lo pendiente. \emph{(RQ1)}
  \item Reportar la evidencia emp\'irica real de pruebas automatizadas,
    cobertura, rendimiento, usabilidad y calidad web recolectada hasta el
    cierre de esta entrega. \emph{(RQ2)}
  \item Auditar la seguridad del sistema desde tres \'angulos
    complementarios (revisi\'on manual OWASP, an\'alisis din\'amico ZAP,
    an\'alisis est\'atico SpotBugs/Find Security Bugs), incluyendo el
    tratamiento expl\'icito del \'unico hallazgo de severidad alta
    detectado. \emph{(RQ2)}
  \item Documentar el despliegue real en producci\'on, la publicaci\'on de
    la imagen de contenedor en GHCR, el archivado permanente en Zenodo y
    la reproducibilidad de punta a punta del artefacto. \emph{(RQ3)}
  \item Discutir los resultados frente a la literatura relacionada
    disponible y declarar honestamente las amenazas a la validez y las
    limitaciones del estudio, incluyendo condiciones a\'un no verificables
    (por ejemplo, operaci\'on sostenida durante per\'iodos prolongados).
    \emph{(RQ1--RQ3)}
\end{enumerate}

\section{Alcance}

El alcance funcional implementado y verificado cubre: registro y
autenticaci\'on de usuarios por cookies seguras; autorizaci\'on por rol y
por propiedad; CRUD completo de mascotas; y los m\'odulos de citas,
consultas y vacunas incorporados en la Unidad~IV, respaldados por seis
rutinas almacenadas de PostgreSQL. Quedan expl\'icitamente \textbf{fuera}
del alcance implementado, y as\'i se documentan en el cap\'itulo de
requisitos (cap\'itulo~\ref{cap:requisitos}): historial cl\'inico
cronol\'ogico completo como m\'odulo independiente, integraci\'on con
dispositivos IoT, facturaci\'on digital, reportes exportables en PDF/Excel,
y recomendaciones cl\'inicas generadas por IA externa. El alcance de
evaluaci\'on emp\'irica cubre pruebas automatizadas, cobertura, rendimiento,
usabilidad, calidad web, y seguridad est\'atica/din\'amica, todas ejecutadas
en un entorno acad\'emico (desarrollo local y despliegue en el plan gratuito
de Render), no en un entorno industrial de alta disponibilidad.

\section{Contribuci\'on del proyecto}

La contribuci\'on principal de BIOPET, frente al panorama de trabajos
relacionados revisado en el cap\'itulo~\ref{cap:trabajos-relacionados}, no
es una t\'ecnica nueva de ingenier\'ia de software, sino la
\textbf{amplitud verificable de la evidencia emp\'irica} recolectada sobre
un mismo artefacto real: cobertura de pruebas, rendimiento bajo carga,
usabilidad medida con instrumento estandarizado, calidad web automatizada, y
seguridad auditada por tres m\'etodos distintos, todo trazable a
requisitos formales y reproducible con comandos documentados
(\texttt{make up}, \texttt{mvn clean verify}, scripts de
\texttt{scripts/}). Ning\'un trabajo relacionado revisado reporta las seis
dimensiones a la vez sobre el mismo sistema (secci\'on~\ref{sec:sintesis-trabajos-relacionados}).

\section{Organizaci\'on de este informe}

El resto del documento se organiza as\'i: el cap\'itulo~\ref{cap:marco-teorico}
resume el marco te\'orico y normativo usado en todo el informe; el
cap\'itulo~\ref{cap:trabajos-relacionados}
revisa los trabajos relacionados bajo un protocolo inspirado en
PRISMA~2020; el cap\'itulo~\ref{cap:metodologia} describe la metodolog\'ia
de investigaci\'on de ingenier\'ia adoptada; el cap\'itulo~\ref{cap:requisitos}
documenta la ingenier\'ia de requisitos; el cap\'itulo~\ref{cap:arquitectura}
la arquitectura y las decisiones de dise\~no; el cap\'itulo~\ref{cap:implementacion}
los m\'odulos implementados; los cap\'itulos~\ref{cap:pruebas-calidad} y
\ref{cap:seguridad} la evidencia de pruebas/calidad y de seguridad; el
cap\'itulo~\ref{cap:despliegue} el despliegue y la reproducibilidad; el
cap\'itulo~\ref{cap:trazabilidad-final} la matriz de trazabilidad; los
cap\'itulos~\ref{cap:resultados-final} y \ref{cap:discusion} integran y
discuten los resultados; el cap\'itulo~\ref{cap:amenazas-limitaciones-final}
documenta las amenazas a la validez y las limitaciones; el
cap\'itulo~\ref{cap:conclusiones-final} presenta las conclusiones; el
cap\'itulo~\ref{cap:declaraciones} recoge las declaraciones formales
(CRediT, \'etica, disponibilidad); y los anexos
(cap\'itulo~\ref{cap:anexos-final}) complementan con glosario, endpoints y
documentos de referencia.

\chapter{Marco te\'orico}
\label{cap:marco-teorico}

Este cap\'itulo re\'une, de forma compacta, los conceptos y normas que
sustentan las decisiones t\'ecnicas y metodol\'ogicas descritas en el resto
del informe. No repite el detalle de aplicaci\'on de cada concepto a
BIOPET (eso corresponde a los cap\'itulos~\ref{cap:requisitos} a
\ref{cap:despliegue}); se limita a definir el marco conceptual usado y a
se\~nalar d\'onde se aplica.

\section{Ingenier\'ia de requisitos}

La especificaci\'on de requisitos de BIOPET sigue la estructura de
ISO/IEC/IEEE~29148:2018, est\'andar internacional que define el contenido
m\'inimo de un documento de Especificaci\'on de Requisitos de Software
(SRS): descripci\'on global del producto, requisitos funcionales y no
funcionales enunciados de forma verificable, y trazabilidad expl\'icita
hacia artefactos derivados (historias de usuario, casos de uso, pruebas).
La redacci\'on individual de cada requisito sigue adem\'as las
caracter\'isticas de calidad de la \emph{INCOSE Guide to Writing
Requirements}~\cite{incose2023requirements} --- en particular
\emph{Unambiguous} (un requisito admite una \'unica interpretaci\'on) y
\emph{Verifiable} (existe un m\'etodo objetivo para confirmar su
cumplimiento) --- aplicadas en el cap\'itulo~\ref{cap:requisitos}. Las
historias de usuario y los casos de uso son las dos t\'ecnicas
complementarias de elicitaci\'on/especificaci\'on usadas para descomponer
cada requisito en un flujo de interacci\'on concreto, seg\'un la pr\'actica
comparada por Pacheco, Garc\'ia \& Reyes~\cite{pacheco2018requirements}
(cap\'itulo~\ref{cap:trabajos-relacionados}).

\section{Arquitectura de software: el modelo C4}

El modelo C4 de Brown~\cite{c4-model} describe la arquitectura de un
sistema en niveles de abstracci\'on progresivos (Contexto, Contenedores,
Componentes, C\'odigo), cada uno dirigido a una audiencia distinta. BIOPET
documenta formalmente los niveles de Contenedores y Componentes
(cap\'itulo~\ref{cap:arquitectura}, secciones~\ref{sec:c4-nivel-2} y
\ref{sec:c4-nivel-3}), eligiendo ese nivel de detalle porque es el que
permite razonar sobre las decisiones de despliegue y de dise\~no interno
del backend sin caer en el detalle de clases individuales, que cambia con
mayor frecuencia que la arquitectura.

\section{Calidad de software: ISO/IEC~25010}

ISO/IEC~25010:2011~\cite{iso25010-2011} define un modelo de calidad de
producto de software organizado en ocho caracter\'isticas (Funcionalidad,
Eficiencia de desempe\~no, Compatibilidad, Usabilidad, Confiabilidad,
Seguridad, Mantenibilidad, Portabilidad), cada una con subcaracter\'isticas
espec\'ificas. Se usa en este informe como vocabulario com\'un para
organizar la evidencia emp\'irica recolectada (cap\'itulo~\ref{cap:pruebas-calidad}),
siguiendo el enfoque de aplicaci\'on pr\'actica descrito por Estdale \&
Georgiadou~\cite{estdale2018applying25010}: mapear cada caracter\'istica
contra evidencia real, declarando expl\'icitamente cu\'ando una
subcaracter\'istica no fue evaluada, en vez de omitirla silenciosamente
(secci\'on~\ref{sec:iso25010}).

\section{Arquitectura REST y JWT}

La API de BIOPET sigue el estilo arquitect\'onico REST (\emph{Representational
State Transfer}) sobre HTTP: recursos identificados por URL, verbos HTTP
con sem\'antica definida, y respuestas sin estado del lado del servidor
entre solicitudes. La autenticaci\'on usa JSON Web Token (JWT), formato
compacto y autocontenido definido en RFC~7519~\cite{rfc7519}: un token
firmado digitalmente que transporta un conjunto de \emph{claims}
(declaraciones) verificables sin necesidad de una consulta adicional al
servidor. BIOPET usa JWT firmado con HMAC-SHA256, transportado en cookies
en vez de en el cuerpo de la respuesta o en \texttt{localStorage}
(cap\'itulo~\ref{cap:seguridad}), decisi\'on de arquitectura registrada en
ADR-006. Los errores de la API siguen RFC~7807~\cite{rfc7807}
(\emph{Problem Details for HTTP APIs}), formato est\'andar para comunicar
fallos HTTP de forma estructurada y legible por m\'aquina.

\section{Seguridad de aplicaciones web: OWASP Top~10}

El OWASP Top~10~\cite{owasp-top10} es la clasificaci\'on de referencia de
la industria para los riesgos de seguridad m\'as cr\'iticos en aplicaciones
web (control de acceso roto, fallas criptogr\'aficas, inyecci\'on,
configuraci\'on de seguridad, fallas de identificaci\'on/autenticaci\'on,
fallas de registro/monitoreo, entre otros). BIOPET usa esta clasificaci\'on
como marco de auditor\'ia manual (cap\'itulo~\ref{cap:seguridad}),
complementada con an\'alisis din\'amico (OWASP ZAP) y est\'atico (SpotBugs
+ Find Security Bugs), siguiendo una l\'ogica de defensa en profundidad:
ning\'un m\'etodo \'unico de auditor\'ia se considera suficiente por s\'i
solo.

\section{Acceso a datos: JPA y procedimientos almacenados}

Spring Data JPA (\emph{Java Persistence API}) provee un mapeo
objeto-relacional (ORM) que traduce operaciones sobre entidades Java a
sentencias SQL generadas autom\'aticamente, adecuado para operaciones CRUD
simples sobre una \'unica tabla. Para operaciones que involucran l\'ogica
de negocio, validaci\'on cruzada entre tablas o agregaciones complejas,
BIOPET usa procedimientos y funciones almacenados de PostgreSQL, invocados
desde el backend mediante \texttt{@Procedure}/\texttt{@NamedStoredProcedureQuery}
de Spring Data JPA (secci\'on~\ref{sec:acceso-jpa-formal}). Esta
estrategia h\'ibrida sigue la misma l\'ogica documentada por Zmaranda et
al.~\cite{zmaranda2020crud}: dejar que cada mecanismo de acceso a datos
resuelva el tipo de operaci\'on para el que fue dise\~nado, en vez de
forzar toda la l\'ogica a trav\'es de un \'unico mecanismo. Esta decisi\'on
es consistente con evidencia emp\'irica comparativa adicional: SQL nativo
supera en tiempo de ejecuci\'on a las consultas generadas por ORM tanto en
C\#~\cite{nowicki2022comparative} como, v\'ia JDBC directo, en
Java~\cite{zuchnik2021comparative}, y la implementaci\'on de la capa de
persistencia tiene un impacto medible sobre el rendimiento general de la
aplicaci\'on~\cite{siebyla2020impact}. La comparaci\'on de rendimiento
entre pilas REST equivalentes (Spring~Boot y .NET~Core), relevante para la
decisi\'on de arquitectura registrada en ADR-002, sigue tambi\'en una
metodolog\'ia de \emph{benchmarking} comparable a la usada
aqu\'i~\cite{dhalla2021performance,godinho2024performance}.

\section{Patr\'on \emph{cache-aside}}

El patr\'on \emph{cache-aside} (tambi\'en llamado \emph{lazy loading})
consiste en que la aplicaci\'on consulta primero la cach\'e; si el dato no
est\'a presente (\emph{cache miss}), lo obtiene del almac\'en de datos
primario y lo escribe en la cach\'e para lecturas futuras; en cada
escritura sobre el dato primario, la entrada de cach\'e correspondiente se
invalida. BIOPET implementa este patr\'on de forma declarativa con las
anotaciones \texttt{@Cacheable} (lectura) y \texttt{@CacheEvict}
(invalidaci\'on) de Spring sobre Redis/Valkey
(secci\'on~\ref{sec:rendimiento-final}, listado~\ref{lst:cache-aside}),
respaldado por trabajos previos sobre la efectividad de esta estrategia en
acceso a datos relacionales~\cite{zulfa2020caching}. La configuraci\'on
del tiempo de vida (TTL) de cada entrada de cach\'e, externalizada por
variable de entorno en BIOPET (\texttt{CACHE\_TTL\_MS}), sigue la misma
recomendaci\'on de an\'alisis a gran escala de cl\'usteres de cach\'e en
memoria: la configuraci\'on del TTL es un factor clave del comportamiento
observado~\cite{yang2021large}. El despliegue contenerizado de BIOPET
(cap\'itulo~\ref{cap:despliegue}) tiene, como cualquier despliegue en
contenedores, un \emph{overhead} medible seg\'un su
configuraci\'on~\cite{ghatrehsamani2020art}.

\section{S\'intesis}

Los conceptos anteriores no se aplican de forma aislada: la ingenier\'ia
de requisitos (29148/INCOSE) alimenta la matriz de trazabilidad
(cap\'itulo~\ref{cap:trazabilidad-final}); el modelo C4 documenta la
arquitectura que implementa esos requisitos; ISO/IEC~25010 organiza la
evidencia emp\'irica que verifica si esa arquitectura cumple sus atributos
de calidad; y REST/JWT/OWASP/JPA/\emph{cache-aside} son las decisiones de
implementaci\'on concretas evaluadas por esa evidencia. Este marco
com\'un es lo que permite, en el cap\'itulo~\ref{cap:discusion}, discutir
los resultados de BIOPET en un mismo vocabulario en vez de como hallazgos
inconexos.

\chapter{Trabajos relacionados}
\label{cap:trabajos-relacionados}

Este cap\'itulo integra la revisi\'on de trabajos relacionados realizada por
Zaida (bloque requisitos/usabilidad/dominio veterinario), Jaime (bloque
arquitectura/seguridad/rendimiento/calidad) y Fred (bloque arquitectura de
datos/rendimiento), consolidada y cerrada con un protocolo inspirado en
PRISMA~2020~\cite{page2021prisma} en \texttt{docs/checklists/prisma2020.md}.
PRISMA fue dise\~nado para revisiones sistem\'aticas de intervenciones
sanitarias; aqu\'i se adapta expl\'icitamente como marco de \textbf{rigor y
transparencia} para una revisi\'on narrativa de trabajos relacionados de
ingenier\'ia de software, no como un metaan\'alisis cl\'inico --- varios
\'items orientados a estad\'istica de efectos cl\'inicos se marcan
\emph{No aplica} en el checklist original en vez de forzarlos.

\section{Criterios de elegibilidad y fuentes}

\textbf{Inclusi\'on:} estudio con DOI o enlace verificable; publicado en
fuente acad\'emica o en una revista/\emph{proceedings} con revisi\'on por
pares (o norma/documento oficial de industria sin DOI, como ISO u OWASP);
relevante a al menos uno de los tres bloques tem\'aticos (requisitos,
usabilidad o dominio veterinario --- Zaida; arquitectura, seguridad,
rendimiento o calidad --- Jaime; arquitectura de datos/rendimiento ---
Fred). \textbf{Exclusi\'on:} blogs comerciales sin revisi\'on por pares,
p\'aginas de venta de software, trabajos sin autor o fecha verificable.
\textbf{Fuentes consultadas:} IEEE Xplore, ACM Digital Library, Scopus y
Google Scholar (bloque Zaida, b\'usqueda 2026-08-17); auditor\'ia de la
bibliograf\'ia ya existente m\'as verificaci\'on dirigida contra Crossref y
p\'aginas oficiales de editorial/conferencia/norma (bloque Jaime, misma
fecha); verificaci\'on dirigida de DOI reales, uno a uno, con un filtro
temporal expl\'icito de 2020--2026 (bloque Fred, 2026-08-17). El proceso de
selecci\'on fue independiente por bloque (cada autor filtr\'o su propio
tema), sin doble revisor cruzado por estudio --- limitaci\'on declarada
expl\'icitamente, no oculta.

De un total de 15 candidatos revisados por Zaida (12 excluidos por no
cumplir el criterio de elegibilidad), una verificaci\'on dirigida de 4
candidatos por Jaime, y 13 candidatos por DOI verificados uno a uno por
Fred, el conjunto final incluido en la \textbf{s\'intesis de trabajos
relacionados de este cap\'itulo} suma \textbf{10 estudios} (Zaida~3 +
Jaime~4 + Fred~3), cumpliendo la meta original de $\geq$8--9. Esos
\textbf{10} son espec\'ificamente los estudios incluidos en esta
s\'intesis narrativa --- no el total de referencias retenidas por los
tres bloques en todo el informe.

\textbf{Trazabilidad completa del bloque Fred (corregida 2026-09-04, tras
auditor\'ia cruzada de Jaime y verificaci\'on de Fred).} La etapa de
\emph{identificaci\'on} parti\'o de 13 candidatos por DOI (rango
2020--2026). La etapa de \emph{verificaci\'on} (17-08-2026) descart\'o 2
(Yang et al., 2018, ICSE --- fuera del rango temporal; un art\'iculo de
TURCOMAT --- inaccesible, HTTP~403), dejando \textbf{11 referencias
elegibles/retenidas}, las 11 con DOI verificado (ver
\texttt{docs/investigacion/handoff-fred-referencias.bib}). Sobre esas 11,
se aplic\'o una \textbf{clasificaci\'on editorial para este informe ---
no un paso de exclusi\'on PRISMA}: \textbf{3} son trabajos relacionados
de dominio/arquitectura y se incorporan a la s\'intesis de este cap\'itulo
(Bucao et al., Llaneta et al., Zmaranda et al. ---
secci\'on~\ref{sec:sintesis-trabajos-relacionados}); las otras
\textbf{8 no fueron excluidas en ning\'un momento}: son referencias de
apoyo metodol\'ogico/rendimiento citadas en otros cap\'itulos del informe
(arquitectura de acceso a datos, cach\'e, despliegue contenerizado ---
\texttt{nowicki2022comparative}, \texttt{zuchnik2021comparative},
\texttt{siebyla2020impact}, \texttt{dhalla2021performance},
\texttt{godinho2024performance}, \texttt{zulfa2020caching},
\texttt{yang2021large}, \texttt{ghatrehsamani2020art}), fuera del alcance
de la s\'intesis narrativa de trabajos relacionados pero citadas con
funci\'on propia en el informe. Los tres bloques est\'an integrados en el
checklist PRISMA (\texttt{docs/checklists/prisma2020.md}) y en el
diagrama de flujo que se muestra en la figura~\ref{fig:prisma-flow}.

% EVIDENCIA-COMPARTIDA-PRISMA
% Ruta esperada: figuras/zaida/06-prisma-flow-diagram.png
\evidencia{figuras/zaida/06-prisma-flow-diagram.png}%
{Diagrama de flujo PRISMA de identificaci\'on y selecci\'on de estudios (regenerado 2026-09-04, trazabilidad del bloque Fred corregida y verificada por Fred). Bloque Zaida completamente cuantificado (15 identificados $\to$ 12 excluidos $\to$ 3 incluidos); bloque Fred completamente cuantificado (identificaci\'on: 13 candidatos $\to$ verificaci\'on, 17-08-2026: 2 descartados $\to$ 11 elegibles/retenidos $\to$ clasificaci\'on editorial para este informe, no un paso de exclusi\'on PRISMA: 3 trabajos relacionados + 8 de apoyo metodol\'ogico/rendimiento, estos \'ultimos citados en otros cap\'itulos, no en la s\'intesis de n=10); bloque Jaime con el conteo final de incluidos (4), sin registro de candidatos excluidos (\'unico gap que persiste, se\~nalado con l\'inea punteada).}%
{fig:prisma-flow}

\section{Cadena de búsqueda, fuentes y ventana temporal por bloque}
\label{sec:cadena-busqueda-trabajos-relacionados}

Los tres bloques usaron estrategias de búsqueda distintas entre sí ---
diferencia declarada explícitamente aquí, no oculta detrás del checklist.
A continuación, la cadena exacta (cuando existió), las bases/fuentes y la
ventana temporal aplicada por cada uno, tal como quedaron registradas el
2026-08-17 en \texttt{docs/informe/borradores/zaida/estudios-primarios-zaida.md},
\texttt{docs/informe/borradores/jaime/trabajos-relacionados.md} y
\texttt{docs/investigacion/handoff-fred-trabajos-relacionados.md}
respectivamente (fuente primaria de cada bloque; el detalle completo,
incluida la tabla de exclusiones, sigue disponible en
\texttt{docs/checklists/prisma2020.md}, sección 4, ítems 6--7).

\textbf{Bloque Zaida (requisitos, usabilidad, dominio veterinario).}
Bases consultadas: IEEE Xplore, ACM Digital Library, Scopus, Google
Scholar. Fecha de búsqueda: 2026-08-17. \textbf{Sin ventana temporal
declarada} (no se aplicó filtro de año como criterio formal). Tres
cadenas de búsqueda booleanas exactas, una por sub-tema:

\begin{quote}
\texttt{"requirements elicitation" AND "systematic literature review"} \\
\texttt{"system usability scale" AND web application} \\
\texttt{veterinary management system web application case study}
\end{quote}

Resultado: 15 candidatos revisados, 12 excluidos por no cumplir el
criterio de elegibilidad de la sección anterior, 3 incluidos.

\textbf{Bloque Jaime (arquitectura, seguridad, rendimiento, calidad).}
\textbf{No usó una cadena booleana única.} Se auditó primero la
bibliografía ya existente en \texttt{docs/informe/referencias.bib} (15
entradas, todas normas/documentación oficial, ninguna un trabajo
académico de \emph{related work}) y, al no encontrar ahí trabajos
académicos pertinentes, se hizo verificación dirigida por tema
(metodología ágil de gestión veterinaria, pruebas de rendimiento web,
vulnerabilidades JWT, aplicación práctica de ISO/IEC~25010) contra
Crossref (\texttt{api.crossref.org}) y páginas oficiales de
editorial/conferencia. Fecha: 2026-08-17. \textbf{Sin ventana temporal
declarada.} Esta es una diferencia metodológica real frente al bloque de
Zaida, no un vacío de documentación: se trató de verificación de
referencias ya identificadas por relevancia temática directa, no de una
búsqueda exploratoria amplia que requiriera una cadena. Resultado: 4
estudios incluidos; el conteo de candidatos evaluados y excluidos antes
de esos 4 no fue registrado por Jaime (limitación declarada en la
sección~\ref{sec:limitaciones-trabajos-relacionados} y en el diagrama de
flujo, figura~\ref{fig:prisma-flow}).

\textbf{Bloque Fred (arquitectura de datos, rendimiento).} Tampoco usó
una cadena booleana: verificación dirigida de DOI reales, uno a uno,
abriendo cada \texttt{https://doi.org/<DOI>} y confirmando título,
autores y año contra Crossref/IEEE/ACM/IOP. Fecha: 2026-08-17.
\textbf{Única ventana temporal explícita de los tres bloques: 2020--2026}
(criterio propio de Fred, ``máx. 5--6 años''). Resultado: 13 candidatos
por DOI, 2 descartados (Yang et al., ICSE 2018, DOI
\texttt{10.1145/3180155.3180194} --- fuera del rango 2020--2026; un
artículo de TURCOMAT --- inaccesible, HTTP~403), 3 incluidos.

\section{Estudios incluidos}

\subsection{Pacheco, Garc\'ia \& Reyes (2018) --- T\'ecnicas de elicitaci\'on de requisitos}
Revisi\'on sistem\'atica sobre 140 estudios (1993--2015) que identifica las
entrevistas estructuradas como la t\'ecnica de elicitaci\'on de requisitos
con mayor evidencia emp\'irica de madurez y efectividad~\cite{pacheco2018requirements}.
El SRS de BIOPET (\texttt{docs/requisitos/SRS.md}) documenta requisitos ya
elicitados (historias de usuario, casos de uso, patr\'on ISO/IEC/IEEE~29148)
pero no declara expl\'icitamente qu\'e t\'ecnica se us\'o para llegar a
ellos; a diferencia de Pacheco et al., que comparan t\'ecnicas entre s\'i,
BIOPET aplica una t\'ecnica ya elegida y la documenta con trazabilidad
formal completa (SRS~$\to$~HU~$\to$~CU~$\to$~matriz), sin reportar esa
comparaci\'on.

\subsection{Bangor, Kortum \& Miller (2008) --- Benchmark emp\'irico del SUS}
Analiza aproximadamente diez a\~nos de datos SUS sobre cientos de
productos, estableciendo un puntaje medio de referencia de la industria de
68/100 (``sobre el promedio'') y 80{,}3 como umbral de
``excelente''~\cite{bangor2008sus}. Este estudio es la fuente del
benchmark contra el que se interpreta el puntaje SUS propio de BIOPET
(secci\'on~\ref{sec:sus-final}): BIOPET no construye un benchmark nuevo,
usa el ya existente para interpretar una medici\'on real sobre un sistema
concreto.

\subsection{Cede\~no Ochoa, Catuto Murillo \& Rodas-Silva (2021) --- Aplicaciones web para cl\'inicas veterinarias en Ecuador}
Caso de estudio sobre el desarrollo de una aplicaci\'on web para gesti\'on
de una cl\'inica veterinaria ecuatoriana, documentando la mejora
cualitativa de procesos administrativos al migrar de manejo manual a una
plataforma centralizada~\cite{cedeno2021veterinaria}. Es el trabajo con
contexto pa\'is m\'as cercano encontrado y verificado. A diferencia de
BIOPET, no reporta una evaluaci\'on emp\'irica multim\'etrica del sistema
resultante (sin cobertura de pruebas, rendimiento bajo carga, seguridad
auditada ni usabilidad medida con instrumento estandarizado).

\subsection{Rodr\'iguez, Llerena, Guevara, Baren \& Castro (2024) --- OSCRUM}
Estudio de caso sobre el desarrollo de SysVet, un sistema web de c\'odigo
abierto para gesti\'on veterinaria, usando la metodolog\'ia \'agil
OSCRUM~\cite{rodriguez2024oscrum}. Comparte con BIOPET el inter\'es en la
reproducibilidad del artefacto (c\'odigo abierto, proceso documentado),
pero documenta el \emph{proceso} de desarrollo \'agil, no una evaluaci\'on
emp\'irica multim\'etrica del \emph{producto} resultante.

\subsection{Putri, Hadi \& Ramdani (2017) --- Pruebas de rendimiento de un sistema web acad\'emico}
Caracteriza cuantitativamente el rendimiento de un sistema de admisi\'on
de estudiantes bajo carga, mediante \emph{benchmarking} t\'ecnico contra el
sistema real~\cite{putri2017perftesting}. Es metodol\'ogicamente an\'alogo a
la evaluaci\'on de rendimiento de BIOPET con k6
(secci\'on~\ref{sec:rendimiento-final}): ambos son estudios de
\emph{benchmarking} de un artefacto web real en un entorno acad\'emico, sin
dise\~no experimental controlado. A diferencia de BIOPET, Putri et al.
miden \'unicamente rendimiento, sin seguridad, usabilidad ni cobertura de
pruebas.

\subsection{Yang et al. (2026) --- Vulnerabilidades en implementaciones de JWT (NDSS)}
Auditor\'ia sistem\'atica de 43 implementaciones reales de JSON Web Token
en diez lenguajes, con una herramienta de \emph{fuzzing} dirigido
(JWTeemo), que descubre 31 vulnerabilidades previamente desconocidas (20
con CVE asignado)~\cite{yang2026tokentimebomb}. BIOPET usa JWT como
mecanismo central de autenticaci\'on (cap\'itulo~\ref{cap:seguridad}); este
trabajo es la referencia acad\'emica m\'as actual disponible sobre los
riesgos reales de esa misma tecnolog\'ia. A diferencia de Yang et al., que
auditan \emph{implementaciones/librer\'ias} gen\'ericas, BIOPET no reclama
haber descubierto vulnerabilidades JWT como contribuci\'on cient\'ifica:
aporta evidencia de que una aplicaci\'on concreta que usa JWT fue auditada
con m\'ultiples t\'ecnicas y no present\'o los patrones de mayor severidad
detectables por ellas, sin afirmar estar ``libre'' de la categor\'ia de
vulnerabilidades de \emph{fuzzing} dirigido a la librer\'ia subyacente
(\texttt{jjwt}), que queda expl\'icitamente fuera del alcance de la
auditor\'ia de BIOPET.

\subsection{Estdale \& Georgiadou (2018) --- Aplicaci\'on pr\'actica de ISO/IEC~25010}
Discute la dificultad pr\'actica de aplicar el modelo de calidad
ISO/IEC~25010 a un producto de software real, m\'as all\'a de la cita
te\'orica de la norma~\cite{estdale2018applying25010}. Es paralelo
metodol\'ogico directo al mapeo realizado en
\texttt{docs/arquitectura/ISO-25010.md} (secci\'on~\ref{sec:iso25010}),
donde se mapean las ocho caracter\'isticas de ISO/IEC~25010:2011 contra
evidencia real de BIOPET, declarando honestamente ``no evaluada
directamente'' cuando no existe evidencia. A diferencia del trabajo de
Estdale \& Georgiadou, que discute el modelo en general, BIOPET aporta un
mapeo concreto con rutas de evidencia verificables por cada
subcaracter\'istica.

\subsection{Bucao, Quinones, Abong, Penuela \& Sun (2023) --- VETelgeuse}
Sistema de gesti\'on web y m\'ovil para cl\'inicas veterinarias peque\~nas y
medianas, evaluado con el cuestionario de usabilidad USE sobre 30
usuarios reales~\cite{bucao2023vetelgeuse}. Confirma, junto con Cede\~no
Ochoa et al. y Llaneta et al., que el dominio de BIOPET (gesti\'on
veterinaria web) tiene evidencia acad\'emica real m\'as all\'a de un
\'unico estudio. A diferencia de BIOPET, su validaci\'on emp\'irica se
limita a usabilidad, sin cobertura de pruebas, rendimiento bajo carga ni
auditor\'ia de seguridad.

\subsection{Llaneta, Guelas, Labanan, Mercado \& Sasis (2022) --- TerraVet}
Propone un \emph{framework} web y m\'ovil para conectar due\~nos de
mascotas con cl\'inicas veterinarias, publicado en ICIST~2022 (ACM), con
una evaluaci\'on preliminar de la propuesta~\cite{llaneta2022terravet}. Es
paralelo de dominio y de arquitectura web+m\'ovil a BIOPET, pero no eval\'ua
rendimiento bajo carga ni despliegue en la nube, ambos reportados en este
informe (cap\'itulos~\ref{cap:pruebas-calidad} y \ref{cap:despliegue}).

\subsection{Zmaranda, Pop-Fele, Gyor\"odi, Gyor\"odi \& Pecherle (2020) --- CRUD con ORM vs. SQL directo}
Compara el tiempo de ejecuci\'on de operaciones CRUD usando mapeadores
objeto-relacionales (ORM) de .NET frente a acceso directo a
datos~\cite{zmaranda2020crud}. Respalda directamente la decisi\'on de
arquitectura de BIOPET de usar un \textbf{acceso h\'ibrido}: Spring Data
JPA para operaciones CRUD simples y SQL nativo en procedimientos
almacenados para consultas complejas y validaciones cruzadas
(secci\'on~\ref{sec:acceso-jpa-formal}). A diferencia de BIOPET, el estudio
de Zmaranda et al. se limita a una sola pila (.NET) y no eval\'ua
procedimientos almacenados ni agregaciones equivalentes a
\texttt{fn\_reporte\_dashboard} o \texttt{fn\_historial\_clinico\_mascota}.

\section{S\'intesis comparativa}
\label{sec:sintesis-trabajos-relacionados}

\begin{table}[H]
  \centering
  \tiny
  \caption{S\'intesis comparativa de los diez estudios incluidos frente a BIOPET.}
  \label{tab:sintesis-relacionados}
  \begin{tabular}{@{}p{2.4cm}p{2.2cm}p{3.0cm}p{4.8cm}@{}}
    \toprule
    Estudio & Dominio & Evaluaci\'on reportada & Diferencia principal frente a BIOPET \\
    \midrule
    Pacheco et al. (2018) & Elicitaci\'on de requisitos & Madurez/efectividad comparada de t\'ecnicas & BIOPET no compara t\'ecnicas; aplica una y la traza formalmente \\
    Bangor et al. (2008) & Usabilidad (SUS) & Benchmark de interpretaci\'on (68/80{,}3) & BIOPET usa el benchmark, no construye uno nuevo \\
    Cede\~no Ochoa et al. (2021) & Gesti\'on veterinaria (Ecuador) & Mejora cualitativa de procesos & Sin evaluaci\'on multim\'etrica; BIOPET s\'i \\
    Rodr\'iguez et al. (2024) --- OSCRUM & Gesti\'on veterinaria, c\'odigo abierto & Proceso \'agil documentado & Documenta proceso, no eval\'ua emp\'iricamente el producto \\
    Putri et al. (2017) & Sistema web acad\'emico & Rendimiento bajo carga & Solo rendimiento; sin seguridad, usabilidad ni cobertura \\
    Yang et al. (2026) --- NDSS & Seguridad JWT (multi-lenguaje) & Vulnerabilidades descubiertas & Audita librer\'ias gen\'ericas, no una aplicaci\'on completa \\
    Estdale \& Georgiadou (2018) & Calidad de software (ISO/IEC~25010) & Discusi\'on de evidenciabilidad & Gu\'ia general, no un mapeo concreto enlazado a evidencia \\
    Bucao et al. (2023) --- VETelgeuse & Gesti\'on veterinaria (Filipinas) & Usabilidad (USE, n=30) & Solo usabilidad; sin rendimiento, seguridad ni cobertura \\
    Llaneta et al. (2022) --- TerraVet & Gesti\'on veterinaria (marco web+m\'ovil) & Propuesta de \emph{framework}, evaluaci\'on preliminar & No eval\'ua rendimiento ni despliegue en la nube \\
    Zmaranda et al. (2020) & Acceso a datos (ORM vs. SQL directo) & Tiempos de operaciones CRUD & Limitado a .NET; no cubre procedimientos almacenados ni agregaciones \\
    \bottomrule
  \end{tabular}
\end{table}

Como resume la tabla~\ref{tab:sintesis-relacionados}, tomados
individualmente los diez estudios eval\'uan siempre \textbf{una
sola dimensi\'on a la vez}: elicitaci\'on de requisitos sin evaluar el
producto resultante, un benchmark de usabilidad sin aplicarlo a un sistema
concreto, el dominio veterinario sin evaluaci\'on multim\'etrica (tres
veces: Cede\~no Ochoa, Bucao, Llaneta), un proceso \'agil sin evaluaci\'on
emp\'irica del producto, rendimiento sin seguridad ni usabilidad, seguridad
JWT gen\'erica sin evaluar una aplicaci\'on completa, el modelo de calidad
discutido sin aplicarlo con evidencia enlazada a un caso real, o
rendimiento de acceso a datos limitado a una sola pila tecnol\'ogica.
BIOPET se posiciona frente a estos diez trabajos como una evaluaci\'on
\textbf{multim\'etrica} de un \'unico artefacto: cobertura de pruebas,
rendimiento, usabilidad, calidad web, y seguridad por tres \'angulos
distintos (revisi\'on manual OWASP, an\'alisis din\'amico ZAP, an\'alisis
est\'atico SpotBugs/Find Security Bugs), todas enmarcadas bajo
ISO/IEC~25010 como vocabulario com\'un. Con la incorporaci\'on del bloque
de Fred, el dominio veterinario exacto de BIOPET pasa de tener un solo
estudio de contraste (Cede\~no Ochoa et al.) a tres (sumando Bucao et al. y
Llaneta et al.), aunque ninguno de los tres combina el contexto
ecuatoriano con una evaluaci\'on multim\'etrica --- esa combinaci\'on
espec\'ifica sigue sin evidencia externa directa. No se afirma que BIOPET
sea metodol\'ogicamente superior a ninguno de estos trabajos: no existe una
comparaci\'on emp\'irica directa que lo demuestre; la diferencia se\~nalada
es de \textbf{amplitud de la evidencia reportada} sobre un mismo sistema,
no de calidad relativa del artefacto.

\section{Limitaciones de esta revisi\'on}
\label{sec:limitaciones-trabajos-relacionados}

\begin{itemize}[nosep]
  \item Solo se incluyeron fuentes en espa\~nol e ingl\'es; podr\'ian existir
    estudios relevantes en otros idiomas no cubiertos.
  \item El bloque de Jaime no document\'o el conteo de candidatos evaluados
    y excluidos (solo el resultado final de 4 incluidos), a diferencia de
    los bloques de Zaida (15 candidatos, 12 exclusiones documentadas) y de
    Fred (13 candidatos por DOI, 2 descartados con motivo individual);
    este gap se se\~nala expl\'icitamente con l\'inea punteada en el
    diagrama de flujo regenerado (figura~\ref{fig:prisma-flow}), \'unico
    bloque que a\'un no est\'a completamente cuantificado.
  \item \textbf{Inconsistencia num\'erica del bloque de Fred (detectada en
    auditor\'ia cruzada de Jaime el 2026-09-04, corregida el mismo d\'ia
    con la trazabilidad que Fred confirm\'o).} Una versi\'on anterior de
    este cap\'itulo y del diagrama de flujo reportaba, para el bloque
    Fred, ``13 candidatos $\to$ 2 excluidos $\to$ 3 incluidos'', dejando 8
    candidatos sin categor\'ia expl\'icita. La trazabilidad real es ``13
    identificados $\to$ 2 descartados $\to$ 11 retenidos $\to$ 3 trabajos
    relacionados + 8 de apoyo metodol\'ogico/rendimiento (no excluidos)'',
    ya reflejada en el p\'arrafo de trazabilidad de Fred arriba y en la
    figura~\ref{fig:prisma-flow}.
  \item La selecci\'on no tuvo un segundo revisor independiente por
    estudio: cada bloque fue filtrado por una sola persona sobre su propio
    tema, sin validaci\'on cruzada entre integrantes --- recurso real de un
    equipo de tres personas en un Proyecto Fin de Curso, declarado
    expl\'icitamente en vez de omitirse. El bloque de Fred aplic\'o,
    adem\'as, un filtro temporal expl\'icito (2020--2026) que los otros dos
    bloques no declararon como criterio formal --- diferencia de criterio
    entre bloques declarada, no oculta.
  \item Con el bloque de Fred, el dominio veterinario exacto de BIOPET pasa
    de un estudio de contraste a tres, pero ninguno combina el contexto
    ecuatoriano con una evaluaci\'on multim\'etrica del sistema: esa
    combinaci\'on espec\'ifica sigue subrepresentada en la literatura
    acad\'emica disponible, no por sesgo de b\'usqueda sino por escasez
    real, hallazgo reportado independientemente por los tres bloques de
    revisi\'on.
\end{itemize}

\chapter{Metodolog\'ia}
\label{cap:metodologia}

Este cap\'itulo integra, en formato de informe acad\'emico, el borrador
metodol\'ogico de Jaime (\texttt{docs/informe/borradores/jaime/metodologia-y-amenazas.md}),
apoyado en el checklist ya completado
\texttt{docs/checklists/ralph2021-engineering-research.md}~\cite{ralph2021-empirical-standards}.

\section{Enfoque metodol\'ogico}

BIOPET corresponde, de forma predominante, a una investigaci\'on de
\textbf{ingenier\'ia orientada al dise\~no y evaluaci\'on de un artefacto de
software} (\emph{Engineering Research}, tambi\'en denominada \emph{Design
Science} en la literatura de ingenier\'ia de software emp\'irica). El
equipo dise\~n\'o, construy\'o y evalu\'o emp\'iricamente un sistema real
--- backend Spring~Boot, frontend Angular e infraestructura Docker ---
frente a requisitos expl\'icitos (\texttt{docs/requisitos/SRS.md}), usando
m\'ultiples fuentes de evidencia t\'ecnica: pruebas automatizadas y
cobertura (JaCoCo), rendimiento (k6), usabilidad (SUS), calidad web
autom\'atica (Lighthouse), y seguridad (OWASP, ZAP, SpotBugs/Find Security
Bugs, auditor\'ia de SQL din\'amico).

\subsection*{Por qu\'e \emph{Engineering Research} y no una alternativa}
\begin{itemize}[nosep]
  \item \textbf{No es un experimento controlado.} Ninguna medici\'on
    manipula deliberadamente una variable independiente con asignaci\'on
    aleatoria a grupos de tratamiento/control. Las corridas de k6 miden el
    comportamiento del propio sistema bajo dos condiciones observacionales
    (cach\'e fr\'ia y caliente), no un dise\~no experimental con grupos
    comparados.
  \item \textbf{No es un estudio de caso observacional tradicional.} BIOPET
    es un artefacto construido por el propio equipo dentro de un Proyecto
    Fin de Curso, no un fen\'omeno observado \emph{in situ} en una
    cl\'inica veterinaria real ya operando con el sistema.
  \item \textbf{S\'i encaja con el patr\'on de \emph{Engineering
    Research}}, cuyos atributos esenciales (descripci\'on del artefacto,
    justificaci\'on de su relevancia, evaluaci\'on conceptual y emp\'irica,
    discusi\'on de limitaciones) est\'an todos presentes y verificados en
    el checklist Ralph et al. ya completado.
\end{itemize}

\subsection*{Afirmaciones que este informe evita deliberadamente}
No se afirma que BIOPET sea un experimento controlado (sin grupos de
control ni asignaci\'on aleatoria); no se afirma validaci\'on industrial
(toda la evaluaci\'on se realiz\'o en un entorno acad\'emico, sin usuarios
ni profesionales de una organizaci\'on externa real); no se afirma ninguna
forma de certificaci\'on (ISO/IEC~25010 se usa como marco de
clasificaci\'on de calidad, no como certificaci\'on obtenida); no se afirma
causalidad experimental en ninguna medici\'on; y no se afirma que los
resultados sean universalmente generalizables.

\section{Design Science Research: las seis actividades de Peffers et al.}
\label{sec:dsr-peffers}

\begin{table}[H]
  \centering
  \small
  \caption{Mapeo de BIOPET sobre las seis actividades de Peffers et al.~\cite{peffers2007dsrm}.}
  \label{tab:peffers}
  \begin{tabular}{@{}p{3.2cm}p{1.6cm}p{7.8cm}@{}}
    \toprule
    Actividad & Estado & Evidencia \\
    \midrule
    1. Identificaci\'on del problema y motivaci\'on & Completada & 15 observaciones docentes reales registradas y cerradas (\texttt{docs/observaciones/OBSERVACIONES.md}); SRS \S1 \\
    2. Definici\'on de objetivos de la soluci\'on & Completada & Requisitos REQ-F/REQ-NF verificables, trazados a HU/CU (\texttt{docs/trazabilidad/matriz.csv}) \\
    3. Dise\~no y desarrollo & Completada & Modelo C4 (contexto/contenedores/componentes), 7 ADR, c\'odigo fuente real en \texttt{Backend/src/}, \texttt{frontend/src/} \\
    4. Demostraci\'on & Completada & \texttt{make up} de punta a punta; colecciones Postman; k6 contra el backend real; ZAP contra el contenedor real \\
    5. Evaluaci\'on & Completada & JaCoCo, k6, SUS, Lighthouse, OWASP/ZAP/SpotBugs, todos con datos reales (secci\'on~\ref{sec:diseno-evaluacion}) \\
    6. Comunicaci\'on & Completada (este documento) & Este informe reemplaza al de la Tercera Entrega, con cifras actualizadas de v1.0.0 \\
    \bottomrule
  \end{tabular}
\end{table}

\section{Goal--Question--Metric (GQM)}
\label{sec:gqm}

\textbf{Goal:} analizar el artefacto BIOPET (backend Spring~Boot +
frontend Angular + infraestructura Docker) \textbf{con el prop\'osito de}
evaluar \textbf{con respecto a} su calidad funcional, mantenibilidad,
rendimiento, usabilidad, calidad web y seguridad, \textbf{desde el punto de
vista} del equipo de desarrollo y de evaluadores acad\'emicos,
\textbf{en el contexto de} un Proyecto Fin de Curso ejecutado en un
entorno de desarrollo local con Docker y desplegado en producci\'on sobre
el plan gratuito de Render~\cite{vansolingen2002gqm}.

\begin{table}[H]
  \centering
  \tiny
  \caption{Preguntas y m\'etricas GQM, con el valor de evidencia vigente al cierre de la Entrega Final.}
  \label{tab:gqm}
  \begin{tabular}{@{}p{2.0cm}p{4.6cm}p{4.6cm}@{}}
    \toprule
    Dimensi\'on & Pregunta & Valor de evidencia (fuente) \\
    \midrule
    Calidad funcional & \textquestiondown Los requisitos declarados tienen respaldo trazable? & Matriz consistente frente al SRS (\texttt{scripts/validate-traceability.sh}) \\
    Calidad funcional & \textquestiondown La suite de pruebas pasa consistentemente? & 205 pruebas, 0 fallos, 0 errores (\texttt{Backend/target/surefire-reports/}) \\
    Mantenibilidad & \textquestiondown Qu\'e proporci\'on del c\'odigo se ejercita? & LINE 91{,}80\,\% / BRANCH 79{,}39\,\%, umbral 70\,\% (\texttt{docs/mediciones/jacoco/METRICS.md}) \\
    Mantenibilidad & \textquestiondown El proceso de build es reproducible y automatizado? & CI con 6 jobs verdes (\texttt{.github/workflows/ci.yml}) \\
    Rendimiento & \textquestiondown Latencia y error del endpoint cacheado bajo carga? & p95 6{,}4--21{,}1~ms (caliente/fr\'io), error 0{,}0\,\% en 10 corridas (\texttt{docs/mediciones/perf/REPORT.md}) \\
    Usabilidad & \textquestiondown Qu\'e nivel de usabilidad percibida presenta el frontend seg\'un los registros de la evaluaci\'on SUS? & SUS media 74{,}44/100, IC95\,\% $[63{,}33,85{,}56]$, n=18 (\texttt{docs/mediciones/sus/REPORT.md}) \\
    Calidad web & \textquestiondown El frontend cumple umbrales base de calidad web? & Performance/Accessibility/Best Practices/SEO cumplidos en las 12 corridas mobile+desktop del 2026-08-18 (\texttt{docs/mediciones/lighthouse/}) \\
    Seguridad (din\'amica) & \textquestiondown Hay hallazgos de severidad alta por escaneo din\'amico? & 0 alertas altas, 1 informativa, ZAP Baseline (\texttt{docs/mediciones/sec/zap/README.md}) \\
    Seguridad (est\'atica) & \textquestiondown El an\'alisis est\'atico detecta patrones de inyecci\'on SQL? & 66 hallazgos totales, 0 de tipo \texttt{SQL\_*} (\rutalarga{docs/mediciones/sec/static-analysis/README.md}) \\
    Seguridad (SQL din\'amico) & \textquestiondown Los procedimientos almacenados construyen SQL de forma segura? & 0 hallazgos, \texttt{exit 0} (\texttt{scripts/audit-sql-dynamic.sh}) \\
    \bottomrule
  \end{tabular}
\end{table}

\section{Dise\~no de evaluaci\'on}
\label{sec:diseno-evaluacion}

Para cada fuente de evidencia emp\'irica se documenta prop\'osito, unidad
evaluada, procedimiento reproducible, m\'etrica principal y limitaci\'on
conocida (detalle ampliado por dimensi\'on en los
cap\'itulos~\ref{cap:pruebas-calidad} y \ref{cap:seguridad}):

\begin{enumerate}[nosep]
  \item \textbf{Pruebas automatizadas:} \texttt{mvn clean verify}
    (JUnit~5 + MockMvc + 2 pruebas de integraci\'on con Testcontainers
    contra PostgreSQL real). M\'etrica: pruebas/fallos/errores.
  \item \textbf{Cobertura JaCoCo:} regla \texttt{BUNDLE} sobre todo el
    m\'odulo backend; contadores LINE/BRANCH/COMPLEXITY.
  \item \textbf{Rendimiento (k6):} 5 corridas en fr\'io + 5 en caliente
    contra \texttt{GET /api/mascotas} sobre HTTPS/TLS~1.3; IC~95\,\% con
    distribuci\'on t de Student; prueba de Wilcoxon pareada entre
    condiciones fr\'ia/caliente.
  \item \textbf{Usabilidad (SUS):} cuestionario de Brooke (1996) sin
    modificar, aplicado a 18 participantes (P01--P18) que, seg\'un
    declaraci\'on del equipo, son externos al equipo; la disponibilidad
    de consentimiento informado individual firmado para cada
    participante no pudo confirmarse documentalmente durante una
    auditor\'ia posterior (ver secci\'on~\ref{sec:sus-final} y
    \texttt{docs/etica/regularizacion-sus/}).
  \item \textbf{Calidad web (Lighthouse):} \texttt{npx @lhci/cli autorun}
    contra el frontend servido por contenedor real, perfiles m\'ovil
    simulado y escritorio, 3 corridas por ruta y por perfil (12 corridas
    en la re-ejecuci\'on final del 2026-08-18).
  \item \textbf{Seguridad din\'amica (ZAP Baseline):} escaneo no
    autenticado contra el backend real dentro de la red Docker Compose.
  \item \textbf{Seguridad est\'atica (SpotBugs + Find Security Bugs):}
    an\'alisis de bytecode compilado, \texttt{effort=Max},
    \texttt{threshold=Low}.
  \item \textbf{Auditor\'ia de SQL din\'amico:} an\'alisis de patrones
    inseguros (\texttt{EXECUTE}, concatenaci\'on) sobre \texttt{db/procs/*.sql}.
  \item \textbf{Trazabilidad de requisitos:} validaci\'on cruzada
    autom\'atica SRS $\leftrightarrow$ matriz.
\end{enumerate}

Cada m\'etodo declara su propia limitaci\'on en el cap\'itulo
correspondiente y en el cap\'itulo~\ref{cap:amenazas-limitaciones-final}; ninguna
cifra de este informe se present\'o sin una fuente citada.

\subsection{Muestreo SUS: reclutamiento, poblaci\'on y l\'imites de generalizaci\'on}
\label{sec:muestreo-sus-metodologia}

Los 18 participantes (P01--P18) de la evaluaci\'on SUS
(secci\'on~\ref{sec:sus-final}) se reclutaron por \textbf{conveniencia,
dentro del c\'irculo cercano al equipo del proyecto}
(\texttt{docs/mediciones/sus/REPORT.md}, secci\'on ``Amenazas a la
validez''): no hubo aleatorizaci\'on, plataforma de reclutamiento pagada
ni acceso a usuarios de una cl\'inica veterinaria real. Se declara
expl\'icitamente hasta d\'onde llega esta evidencia, sin completar
detalles de canal de reclutamiento (por ejemplo, redes sociales
espec\'ificas o convocatoria formal) que no quedaron registrados y que el
equipo no puede sostener con evidencia.

\textbf{Poblaci\'on de origen (nivel de generalidad sostenible por la
evidencia):} adultos con acceso a un dispositivo con navegador web y
disposici\'on a colaborar con el equipo, sin relaci\'on documentada con
el dominio veterinario ni con la operaci\'on real de una cl\'inica. La
muestra result\'o demogr\'aficamente heterog\'enea dentro de ese c\'irculo
(\texttt{docs/mediciones/sus/REPORT.md}): edad 19--61 a\~nos (media
34{,}3), sexo balanceado (F=9, M=9), experiencia previa con aplicaciones
web variada (avanzada=6, intermedia=7, b\'asica=3, ninguna=2) y
dispositivo variado (laptop=10, escritorio=4, tablet=2, celular=2) --- lo
que \textbf{atenu\'a, sin eliminar}, la homogeneidad t\'ipica de una
muestra de conveniencia reclutada dentro de un c\'irculo social \'unico.

\textbf{Por qu\'e se us\'o muestreo por conveniencia y qu\'e sesgos
introduce.} Es el m\'etodo real y disponible para un equipo de tres
estudiantes en un Proyecto Fin de Curso, sin presupuesto para una
plataforma de reclutamiento ni acceso institucional a personal o
clientes de una cl\'inica veterinaria real; se declara como limitaci\'on
de recursos, no como decisi\'on metodol\'ogica \'optima. Introduce al
menos dos sesgos concretos: (1) \textbf{sesgo de selecci\'on} hacia
personas con alguna conexi\'on social al equipo, que pueden compartir
rasgos (nivel educativo, familiaridad general con tecnolog\'ia) no
representativos de la poblaci\'on general; y (2) \textbf{posible sesgo de
complacencia/deseabilidad social}, ya declarado en
\texttt{docs/mediciones/sus/REPORT.md}: participantes con relaci\'on
personal con el equipo pueden puntuar de forma m\'as favorable que un
usuario sin esa relaci\'on.

\textbf{A qu\'e poblaci\'on se pueden y no se pueden extrapolar los
resultados.} Los resultados (SUS medio 74{,}44/100, categor\'ia
``Bueno'') son razonablemente informativos sobre la percepci\'on de
usabilidad de \textbf{adultos con acceso a la web y diversidad et\'aria y
de experiencia digital similar a la observada en la muestra}, al
completar una tarea de \emph{onboarding} controlada (inicio de
sesi\'on, alta, edici\'on y baja l\'ogica de una mascota) en un entorno
sin presi\'on operativa real. \textbf{No se pueden extrapolar} a
usuarios reales de una cl\'inica veterinaria operando bajo carga de
trabajo real (veterinarios, auxiliares o personal administrativo
durante una jornada cl\'inica), ni a due\~nos de mascota en el momento de
una urgencia real, contextos que la evaluaci\'on no reprodujo; tampoco se
document\'o la ubicaci\'on geogr\'afica ni la relaci\'on de los
participantes con el sector veterinario, por lo que ese eje de
representatividad tampoco puede argumentarse en ning\'un sentido.

\textbf{Respaldo bibliogr\'afico.} Ralph, Baltes et al.
(2021)~\cite{ralph2021-empirical-standards}, en su est\'andar de
Ingenier\'ia de Software Emp\'irica, tratan el muestreo por conveniencia
como una pr\'actica leg\'itima y frecuente en investigaci\'on de
ingenier\'ia de software \textbf{siempre que la amenaza se reporte de
forma expl\'icita y la generalizaci\'on reclamada se acote al alcance que
la muestra realmente sostiene} --- exactamente el criterio aplicado en
esta secci\'on: no se afirma representatividad de la poblaci\'on general
ni del dominio veterinario, solo del segmento de poblaci\'on
efectivamente alcanzado.

\textbf{L\'imite \'etico declarado (sin evidencia inventada).} Como ya
se documenta en la secci\'on~\ref{sec:sus-final} y en
\texttt{docs/etica/regularizacion-sus/}, no existe evidencia verificable
de consentimiento informado individual firmado por cada participante, y
la constancia de regularizaci\'on firmada por el equipo \textbf{no
sustituye} esos consentimientos. Esta secci\'on no afirma una aprobaci\'on
\'etica previa ni consentimientos que no existan; el reclutamiento por
conveniencia descrito arriba no depende de ---ni corrige--- ese vac\'io
documental, que sigue declarado como limitaci\'on abierta.

\section{Reproducibilidad}

El sistema completo se reconstruye desde una clonaci\'on limpia con un
\'unico comando (\texttt{make up}), con im\'agenes de terceros fijadas por
\emph{digest} \texttt{sha256} para evitar derivas silenciosas entre
reconstrucciones. Los scripts que generan cada pieza de evidencia
(\rutalarga{scripts/archive-jacoco-evidence.sh},
\rutalarga{scripts/run-zap-baseline.sh}, \rutalarga{scripts/audit-sql-dynamic.sh},
\rutalarga{scripts/validate-traceability.sh}, \rutalarga{scripts/perf-analysis.py},
\rutalarga{scripts/analisis-sus.py}, \rutalarga{scripts/run-lighthouse.sh}, entre
otros) est\'an versionados, de modo que la evidencia no es solo consultable
sino regenerable. El detalle completo de reproducibilidad de despliegue se
documenta en el cap\'itulo~\ref{cap:despliegue}.

\section{Amenazas a la validez (resumen metodol\'ogico)}

Las amenazas a la validez de cada medici\'on se documentan en detalle en
el cap\'itulo~\ref{cap:amenazas-limitaciones-final}, clasificadas seg\'un
las cuatro categor\'ias est\'andar (interna, externa, de constructo y de
conclusi\'on). Este cap\'itulo de metodolog\'ia se limita a declarar el
principio que las gobierna: cada amenaza documentada incluye su efecto
posible, la mitigaci\'on ya aplicada (si existe) y el riesgo residual que
persiste, sin minimizar ninguna de ellas.

\chapter{Ingenier\'ia de requisitos}
\label{cap:requisitos}

Este cap\'itulo integra \texttt{docs/requisitos/SRS.md} (Especificaci\'on
de Requisitos de Software, con patr\'on ISO/IEC/IEEE~29148), las historias
de usuario, los casos de uso y \texttt{docs/requisitos/CHANGELOG-REQ.md}.
El SRS distingue expl\'icitamente requisitos \textbf{verificados},
\textbf{implementados} (sin prueba automatizada dedicada) y
\textbf{pendientes}; este cap\'itulo respeta esa misma distinci\'on: ning\'un
requisito que el SRS marca como pendiente se declara aqu\'i como cumplido.

\section{Priorizaci\'on MoSCoW}

Los requisitos se priorizan con la t\'ecnica MoSCoW (Must, Should, Could,
Won't). El n\'ucleo \emph{Must Have} corresponde al compromiso funcional
operativo de esta entrega: autenticaci\'on, autorizaci\'on por rol y por
propiedad, y CRUD de mascotas. Los m\'odulos de citas, consultas y
vacunas, incorporados durante la Unidad~IV, se formalizaron como
\texttt{REQ-F-023} a \texttt{REQ-F-025} (secci\'on~\ref{sec:req-unidad4}).

\section{Requisitos funcionales implementados y verificados}

Los requisitos \texttt{REQ-F-001} a \texttt{REQ-F-012} y \texttt{REQ-F-021}
cubren registro de usuario, login/logout/refresh por cookies,
autorizaci\'on por rol, perfil del usuario autenticado, y el CRUD completo
de mascotas con resumen por especie; todos con estado \textbf{verificado}
en el SRS, respaldados por pruebas automatizadas de controlador/servicio
(cap\'itulo~\ref{cap:pruebas-calidad}).

\subsection{Requisitos incorporados en la Unidad~IV (\texttt{REQ-F-023} a \texttt{REQ-F-025})}
\label{sec:req-unidad4}

Estos tres requisitos formalizan funcionalidades reales incorporadas al
backend durante la Unidad~IV, distintas de Citas y Consultas (que ya
ten\'ian requisito formal hist\'orico propio --- \texttt{REQ-F-015} y
\texttt{REQ-F-013} respectivamente --- y no se duplican aqu\'i):

\begin{itemize}[nosep]
  \item \textbf{REQ-F-023 --- Gesti\'on administrativa de usuarios}
    (\texttt{UsuarioController}/\texttt{UsuarioService}): un
    \texttt{ADMIN} lista, consulta, crea, actualiza y da de baja l\'ogica
    cuentas de cualquier rol, sin poder autoescalar el rol de su propia
    cuenta. Estado: \textbf{verificado} (\texttt{UsuarioControllerTest},
    16 pruebas).
  \item \textbf{REQ-F-024 --- Gesti\'on de vacunas} (\texttt{VacunaController}):
    registro, consulta (global, por mascota o por identificador),
    actualizaci\'on y baja l\'ogica de vacunas, con propiedad restringida
    para \texttt{ROLE\_DUENO}. Estado: \textbf{verificado}
    (\texttt{VacunaControllerTest}, 10 pruebas; colecci\'on Postman
    \texttt{BIOPET-Vacunas.postman\_collection.json}); la operaci\'on
    \texttt{GET /api/vacunas/mascota/\{mascotaId\}} no tiene prueba
    automatizada dedicada.
  \item \textbf{REQ-F-025 --- Consulta de informaci\'on externa de especies}
    (\texttt{ExternalApiController}): consulta de datos de una especie
    (nombre cient\'ifico, h\'abitat, dieta) desde un servicio externo (API
    Ninjas), con cach\'e Redis \emph{cache-aside} para evitar llamadas
    repetidas. Estado: \textbf{implementado} (sin prueba automatizada
    dedicada; no se marca ``verificado'' siguiendo la misma convenci\'on ya
    usada en el resto del SRS); evidencia emp\'irica manual de la
    comparaci\'on cach\'e fr\'ia/caliente en
    \rutalarga{docs/u4/evidencias/fred/redis-cache-comparacion.md}.
\end{itemize}

Rendimiento bajo carga de citas, consultas y vacunas no fue medido
espec\'ificamente con k6 (la medici\'on de rendimiento del
cap\'itulo~\ref{cap:pruebas-calidad}, secci\'on~\ref{sec:rendimiento-final},
se limita al endpoint de mascotas), limitaci\'on declarada en el
cap\'itulo~\ref{cap:amenazas-limitaciones-final}.

Las seis rutinas almacenadas de PostgreSQL (\texttt{fn\_resumen\_mascotas\_por\_especie},
\texttt{fn\_historial\_clinico\_mascota}, \texttt{fn\_reporte\_dashboard},
\texttt{fn\_siguiente\_numero\_ficha}, \texttt{sp\_actualizar\_estado\_citas\_masivas},
\texttt{sp\_registrar\_consulta\_validada}) respaldan estos m\'odulos con
acceso JPA formal (secci\'on~\ref{sec:acceso-jpa-formal}).

\section{Requisitos pendientes, heredados de la Entrega~1A}

Los requisitos \texttt{REQ-F-013} a \texttt{REQ-F-020} y \texttt{REQ-F-022}
permanecen, con honestidad, en estado \textbf{pendiente} o \textbf{parcial}
en el SRS, sin implementaci\'on de backend ni de frontend a la fecha de
este informe salvo lo explicitado:

\begin{table}[H]
  \centering
  \small
  \caption{Requisitos funcionales pendientes o parciales heredados de la Entrega 1A (fuente: \texttt{docs/requisitos/SRS.md}).}
  \label{tab:req-pendientes}
  \begin{tabular}{@{}p{2.0cm}p{6.4cm}p{1.6cm}p{2.0cm}@{}}
    \toprule
    Requisito & Descripci\'on & Prioridad & Estado \\
    \midrule
    REQ-F-013 & Registro/consulta de historial cl\'inico cronol\'ogico como m\'odulo independiente & Should & Pendiente \\
    REQ-F-014 & Registro de medicamentos prescritos & Could & Pendiente \\
    REQ-F-015 & Calendario interactivo de citas (m\'odulo b\'asico de citas ya existe; el calendario interactivo no) & Should & Pendiente \\
    REQ-F-016 & API de recepci\'on de telemetr\'ia IoT & Could & Pendiente \\
    REQ-F-017 & Recomendaciones cl\'inicas generadas por IA externa & Could & Pendiente \\
    REQ-F-018 & Facturaci\'on digital y registro de pagos & Should & Pendiente \\
    REQ-F-019 & Reportes estad\'isticos exportables (PDF/Excel) & Could & Pendiente \\
    REQ-F-020 & Auditor\'ia de operaciones del sistema & Should & Parcial \\
    REQ-F-022 & Notificaciones por correo electr\'onico & Could & Pendiente \\
    \bottomrule
  \end{tabular}
\end{table}

\textbf{Nota de honestidad sobre \texttt{REQ-F-015} y \texttt{REQ-F-013}.}
El backend expone \texttt{CitaController} y \texttt{ConsultaController}
reales, con endpoints funcionales y respaldados por las rutinas
almacenadas de la Unidad~IV (secci\'on~\ref{sec:req-unidad4}); sin embargo,
el SRS mantiene \texttt{REQ-F-013} (historial cl\'inico cronol\'ogico
completo, con consulta longitudinal) y la parte de \emph{calendario
interactivo} de \texttt{REQ-F-015} como pendientes, porque esas
funcionalidades espec\'ificas --- m\'as all\'a del CRUD b\'asico de
consultas y citas ya implementado --- no est\'an completas ni tienen
frontend dedicado. Este informe reporta ambos estados tal como los declara
el SRS vigente, sin resolver la aparente tensi\'on ampliando el alcance
por cuenta propia: se trata como lo que es, una actualizaci\'on parcial del
SRS que a\'un no reclasific\'o expl\'icitamente cada subrequisito heredado
frente al c\'odigo nuevo de la Unidad~IV, y se documenta como limitaci\'on
de consistencia documental en el cap\'itulo~\ref{cap:amenazas-limitaciones-final}.

\texttt{REQ-F-020} (auditor\'ia) se marca \textbf{parcial}, no
\textbf{pendiente}: el registro de eventos de autenticaci\'on
(\texttt{AUTH\_AUDIT}, cap\'itulo~\ref{cap:seguridad}) s\'i existe y est\'a
verificado, pero no cubre auditor\'ia gen\'erica de todas las operaciones
CRUD sobre mascotas ni de los m\'odulos de citas/consultas/vacunas.

\section{Requisitos no funcionales relevantes}

\begin{table}[H]
  \centering
  \small
  \caption{Requisitos no funcionales relevantes citados en este informe (fuente: \texttt{docs/requisitos/SRS.md}, secci\'on 3.2).}
  \label{tab:req-nf-relevantes}
  \begin{tabular}{@{}p{2.0cm}p{6.6cm}p{2.6cm}@{}}
    \toprule
    Requisito & Descripci\'on & Estado \\
    \midrule
    REQ-NF-001 & Rendimiento del listado de mascotas ($\leq$200~ms p95 caliente, $\leq$500~ms fr\'ia) & Verificado \\
    REQ-NF-002 & Comunicaci\'on cifrada obligatoria (HTTPS/TLS) & Verificado (TLS de desarrollo) \\
    REQ-NF-003 & Expiraci\'on configurable de tokens JWT & Verificado \\
    REQ-NF-004 & Claims est\'andar del JWT (RFC~7519) & Verificado \\
    REQ-NF-005 & Interfaz responsiva (320--1440px) / Lighthouse & Verificado (secci\'on~\ref{sec:lighthouse-final}) \\
    REQ-NF-006 & Compatibilidad de navegadores (Chrome/Firefox/Edge) & Pendiente de evidencia archivada \\
    REQ-NF-008 & Cifrado de contrase\~nas (BCrypt) & Verificado \\
    REQ-NF-009 & Registro de eventos de autenticaci\'on (OWASP A09) & Verificado \\
    REQ-NF-010 & Limitaci\'on de intentos de login (OWASP A07) & Verificado \\
    REQ-NF-011 & Arquitectura en capas (\texttt{controller}/\texttt{service}/\texttt{repository}) & Verificado \\
    REQ-NF-012 & Cach\'e del listado de mascotas con TTL externo & Verificado parcialmente (TTL confirmado; sin hit ratio medido) \\
    REQ-NF-013 & Estrategia h\'ibrida de acceso a datos (JPA + procedimientos almacenados) & Verificado (cap\'itulo~\ref{cap:arquitectura}) \\
    \bottomrule
  \end{tabular}
\end{table}

El requisito de interfaz responsiva/calidad web (\texttt{REQ-NF-005}) est\'a
\textbf{verificado}: las cuatro categor\'ias de Lighthouse (Performance,
Accessibility, Best Practices, SEO) cumplen su umbral en la corrida del
2026-08-18, en ambos perfiles mobile y desktop
(secci\'on~\ref{sec:lighthouse-final}) --- superando el estado
``verificado parcialmente'' que declaraba una revisi\'on anterior del SRS,
fechada antes de esa corrida. \texttt{REQ-NF-012} (hit ratio de cach\'e
Redis) permanece verificado solo parcialmente: existe evidencia de TTL y
de existencia de clave bajo carga, pero no una medici\'on expl\'icita de
\texttt{keyspace\_hits}/\texttt{keyspace\_misses} a lo largo del tiempo.
No existe, en el SRS actual, un requisito no funcional formal dedicado
espec\'ificamente a la usabilidad medida con SUS: la medici\'on SUS
(secci\'on~\ref{sec:sus-final}) responde a un criterio de evaluaci\'on de
la Gu\'ia de la Entrega Final, no a un \texttt{REQ-NF} numerado del SRS.

\section{Trazabilidad y estado de la validaci\'on autom\'atica}

\texttt{scripts/validate-traceability.sh} valida de forma cruzada que todo
requisito del SRS tenga fila en \texttt{docs/trazabilidad/matriz.csv} con
historia de usuario y caso de uso asociados. El detalle completo de la
matriz se presenta en el cap\'itulo~\ref{cap:trazabilidad-final}.

\section{Marco INCOSE}

La redacci\'on de cada requisito sigue el patr\'on ``El sistema
deber\'a\ldots'' con entradas, resultado esperado y criterios de
aceptaci\'on verificables, alineado con las caracter\'isticas de calidad de
requisitos de la \emph{INCOSE Guide to Writing Requirements}~\cite{incose2023requirements}
(\texttt{docs/checklists/incose2023-req.md}): en particular
\emph{Unambiguous} y \emph{Verifiable}, aplicadas expl\'icitamente al
cerrar la ambig\"uedad hist\'orica de \texttt{REQ-F-017} (observaci\'on
OBS-04, ya cerrada).

\section{Cambios entre versiones}

\texttt{docs/requisitos/CHANGELOG-REQ.md} documenta la evoluci\'on del SRS
entre la Entrega~1A, la Entrega~1B, la Tercera Entrega y la Entrega Final,
incluyendo la incorporaci\'on de \texttt{REQ-F-022} (recuperaci\'on formal
de RF-07, cierre de OBS-02) y de \texttt{REQ-F-023} a \texttt{REQ-F-025}
(Unidad~IV). Ning\'un identificador de requisito se reutiliz\'o para dos
funcionalidades distintas.

\chapter{Dise\~no y arquitectura}
\label{cap:arquitectura}

Este cap\'itulo integra la documentaci\'on arquitect\'onica versionada en
\texttt{docs/diagrams/} (modelo C4~\cite{c4-model}) y las decisiones de
arquitectura relevantes (\texttt{docs/adr/}), actualizada para reflejar el
estado v1.0.0 (PostgreSQL~18/Valkey~8 en producci\'on, acceso JPA formal a
procedimientos almacenados).

\section{Arquitectura general}

BIOPET sigue una arquitectura de tres niveles: un \emph{frontend} Angular~17
(aplicaci\'on de p\'agina \'unica, componentes \emph{standalone}), un
\emph{backend} Spring~Boot~3.2 sobre Java~21 que expone una API
REST~\cite{spring-boot-docs}, una base de datos relacional
PostgreSQL~\cite{postgresql-docs} (16 en desarrollo local, 18 en
producci\'on) gestionada con Flyway (migraciones \texttt{V1} a
\texttt{V6}), y Redis/Valkey~\cite{redis-docs} (7 en desarrollo, Valkey~8
en producci\'on) usado tanto para la lista negra de tokens JWT revocados
como para la cach\'e de lecturas de mascotas. Todo el sistema se
levanta mediante Docker Compose y un \texttt{Makefile} con un \'unico
comando reproducible (\texttt{make up}).

\section{C4 Nivel 2 --- Contenedores}
\label{sec:c4-nivel-2}

El diagrama de contenedores (\texttt{docs/diagrams/c4-contenedores/})
describe los cuatro contenedores del sistema: el frontend Angular, el
backend Spring Boot, PostgreSQL y Redis/Valkey, junto con los protocolos
reales entre ellos (HTTPS/JSON con cookies entre frontend y backend; JDBC
entre backend y PostgreSQL; el cliente Redis de Spring Data entre backend
y Redis/Valkey).

% EVIDENCIA-COMPARTIDA-C4-L2
% Ruta esperada: figuras/compartidas/c4-contenedores.png
\evidencia{figuras/compartidas/c4-contenedores.png}%
{Diagrama C4 Nivel 2 (contenedores) de BIOPET: frontend, backend, PostgreSQL y Redis/Valkey, con los protocolos reales entre ellos.}%
{fig:c4-nivel-2-final}

\section{C4 Nivel 3 --- Componentes del backend}
\label{sec:c4-nivel-3}

El diagrama de componentes del backend fue elaborado mediante inspecci\'on
directa del c\'odigo fuente (constructores, llamadas directas,
configuraci\'on de Spring), no por convenci\'on de nombres. A los
componentes ya documentados en la Tercera Entrega (controladores de
autenticaci\'on, mascotas y usuario; servicios de seguridad; manejo de
errores; persistencia; infraestructura del servidor) se suman, en v1.0.0,
los controladores de la Unidad~IV: \texttt{CitaController},
\texttt{ConsultaController}, \texttt{VacunaController} y
\texttt{ExternalApiController}, junto con el repositorio formal de acceso a
procedimientos almacenados.

% EVIDENCIA-JAIME-C4-L3
% Ruta esperada: figuras/jaime/c4-componentes-backend.png
\evidencia{figuras/jaime/c4-componentes-backend.png}%
{Diagrama C4 Nivel 3 (componentes) del backend de BIOPET.}%
{fig:c4-nivel-3-final}

\section{Acceso JPA formal a procedimientos almacenados}
\label{sec:acceso-jpa-formal}

Las seis rutinas de \texttt{db/procs/} son objetos PostgreSQL de tipo
\texttt{PROCEDURE}, invocadas desde el backend mediante
\texttt{@Procedure}/\texttt{@NamedStoredProcedureQuery} de Spring Data JPA,
en vez de \texttt{@Query(nativeQuery=true)} plano (usado en la Tercera
Entrega solo para la primera rutina). Las cuatro rutinas con prefijo
\texttt{fn\_*} originalmente eran \texttt{FUNCTION}; se reclasificaron a
\texttt{PROCEDURE} (conservando el nombre por compatibilidad) para permitir
esa invocaci\'on formal.

\begin{table}[H]
  \centering
  \small
  \caption{Cat\'alogo de las seis rutinas almacenadas de BIOPET (fuente: \texttt{docs/basedatos/CATALOGOSP.md}).}
  \label{tab:catalogo-sp}
  \begin{tabular}{@{}p{4.0cm}p{2.2cm}p{6.6cm}@{}}
    \toprule
    Rutina & Categor\'ia & Prop\'osito \\
    \midrule
    \texttt{fn\_resumen\_mascotas\_por\_especie} & Agregado & Resumen de mascotas activas por especie, filtrable por due\~no \\
    \texttt{fn\_historial\_clinico\_mascota} & Multi-tabla & Historial cl\'inico consolidado: datos, due\~no, conteos de consultas/citas, \'ultima consulta, \'ultima y pr\'oxima vacuna \\
    \texttt{fn\_reporte\_dashboard} & Reporte & Indicadores de dashboard para un rango de fechas \\
    \texttt{fn\_siguiente\_numero\_ficha} & C\'odigo secuencial & Siguiente n\'umero de ficha \texttt{PREFIJO-NNNNNN} \\
    \texttt{sp\_actualizar\_estado\_citas\_masivas} & Actualizaci\'on masiva & \texttt{UPDATE} controlado del estado de citas por veterinario/fecha l\'imite \\
    \texttt{sp\_registrar\_consulta\_validada} & Validaci\'on cruzada & Registra una consulta solo si la mascota y el veterinario son v\'alidos y activos \\
    \bottomrule
  \end{tabular}
\end{table}

\textbf{Estrategia de migraci\'on sin romper compatibilidad.} La
reclasificaci\'on no se aplic\'o reescribiendo la migraci\'on Flyway
original (\texttt{V5\_\_procedimientos\_biopet.sql}, que pudo haberse
aplicado ya en bases persistentes: reescribirla habr\'ia producido
\texttt{checksum mismatch}). En su lugar, una migraci\'on nueva y aditiva
(\texttt{V6\_\_formalizar\_procedimientos\_jpa.sql}) hace \texttt{DROP
FUNCTION IF EXISTS} de las cuatro rutinas creadas por \texttt{V5} y las
vuelve a crear como \texttt{PROCEDURE} con su firma final. Flyway aplica
\texttt{V1} a \texttt{V6} en orden sobre un checkout nuevo, y solo
\texttt{V6} sobre una base que ya tenga \texttt{V1--V5} aplicadas.

% EVIDENCIA-FRED-STOREPROCEDURE
% Ruta esperada: figuras/fred/StoreProceduro-AccesoHibrido.png
\evidencia{figuras/fred/StoreProceduro-AccesoHibrido.png}%
{Acceso h\'ibrido a procedimientos almacenados: invocaci\'on formal JPA sobre las rutinas reclasificadas como \texttt{PROCEDURE}.}%
{fig:sp-acceso-hibrido-final}

% EVIDENCIA-FRED-POSTMAN-SP
% Ruta esperada: figuras/fred/Postman-STOREPROCEDURE.png
\evidencia{figuras/fred/Postman-STOREPROCEDURE.png}%
{Invocaci\'on real de una rutina almacenada verificada desde Postman.}%
{fig:postman-sp-final}

\section{Decisiones de arquitectura relevantes (ADR)}
\label{sec:adr-relevantes-final}

\begin{table}[H]
  \centering
  \caption{ADR vigentes del repositorio.}
  \label{tab:adr-relevantes-final}
  \begin{tabular}{@{}p{2cm}p{9.5cm}p{2.5cm}@{}}
    \toprule
    ADR & Decisi\'on & Estado \\
    \midrule
    ADR-002 & Migraci\'on de la pila de backend de ASP.NET Core~8 a Java~21 / Spring Boot~3.2. & Aceptado \\
    ADR-003 & Revocaci\'on de JWT mediante lista negra en Redis, indexada por \texttt{jti}. & Aceptado \\
    ADR-004 & Estrategia de base de datos reproducible: Flyway como fuente de verdad del esquema, m\'as scripts de inicializaci\'on de Docker con privilegios m\'inimos. & Aceptado \\
    ADR-005 & Reproducibilidad del despliegue: \texttt{Makefile} y pinning de im\'agenes por digest \texttt{sha256}. & Aceptado \\
    ADR-006 & Decisiones de autenticaci\'on y seguridad: JWT por cookies, claims, revocaci\'on, rate limiting, cabeceras, TLS. & Aceptado \\
    ADR-007 & Acceso formal JPA a procedimientos almacenados de PostgreSQL (\texttt{@Procedure}/\texttt{@NamedStoredProcedureQuery}), reclasificaci\'on \texttt{fn\_*} a \texttt{PROCEDURE}. & Aceptado \\
    \bottomrule
  \end{tabular}
\end{table}

ADR-001 (selecci\'on inicial de ASP.NET Core~8 en la Entrega~1A) qued\'o
formalmente superado por ADR-002 y se conserva \'unicamente como registro
hist\'orico.

\section{Patrones y decisiones de seguridad relevantes al dise\~no}

El patr\'on \emph{cache-aside} declarativo de Spring
(\texttt{@Cacheable}/\texttt{@CacheEvict} en \texttt{MascotaService})
respalda las lecturas de \texttt{GET /api/mascotas}, con Redis/Valkey como
almac\'en de cach\'e. La autenticaci\'on usa JWT firmado con HMAC-SHA256,
transportado en cookies \texttt{HttpOnly}+\texttt{Secure}+\texttt{SameSite=Strict},
con revocaci\'on real v\'ia lista negra en Redis y rate limiting de login
en memoria (detalle completo en el cap\'itulo~\ref{cap:seguridad}). El
formato uniforme de error sigue RFC~7807 (\texttt{ProblemDetail})~\cite{rfc7807}.

\section{ISO/IEC~25010: marco de calidad}
\label{sec:iso25010}

\texttt{docs/arquitectura/ISO-25010.md} mapea las ocho caracter\'isticas de
calidad de ISO/IEC~25010:2011~\cite{iso25010-2011} (Funcionalidad,
Eficiencia de desempe\~no, Compatibilidad, Usabilidad, Confiabilidad,
Seguridad, Mantenibilidad, Portabilidad) contra evidencia real de BIOPET,
siguiendo el enfoque de aplicaci\'on pr\'actica de Estdale \&
Georgiadou~\cite{estdale2018applying25010} (secci\'on~\ref{cap:trabajos-relacionados}).
Cada subcaracter\'istica se declara \textbf{evaluada} (con ruta de
evidencia citada) o \textbf{no evaluada directamente} cuando no existe
medici\'on espec\'ifica --- por ejemplo, \emph{Interoperability} m\'as
all\'a de la API REST documentada con OpenAPI, o \emph{Adaptability} m\'as
all\'a de la configuraci\'on externalizada por variables de entorno.

\begin{table}[H]
  \centering
  \small
  \caption{Mapeo resumido de ISO/IEC~25010 sobre evidencia de BIOPET.}
  \label{tab:iso25010-resumen}
  \begin{tabular}{@{}p{3.4cm}p{9.4cm}@{}}
    \toprule
    Caracter\'istica & Evidencia principal \\
    \midrule
    Funcionalidad (\emph{Functional Suitability}) & Matriz de trazabilidad, SRS, 205 pruebas automatizadas \\
    Eficiencia de desempe\~no & 10 corridas k6 (\texttt{docs/mediciones/perf/REPORT.md}) \\
    Compatibilidad & Docker Compose multi-servicio, API REST documentada con OpenAPI \\
    Usabilidad & SUS (n=18), atributos ARIA verificados en c\'odigo \\
    Confiabilidad & 0{,}0\,\% tasa de error en las 10 corridas k6, healthcheck de producci\'on \texttt{UP} \\
    Seguridad & OWASP Top~10 manual, ZAP Baseline, SpotBugs/Find Security Bugs (cap\'itulo~\ref{cap:seguridad}) \\
    Mantenibilidad & Cobertura JaCoCo 91{,}80\,\%/79{,}39\,\%, CI con 6 jobs \\
    Portabilidad & Im\'agenes Docker por digest \texttt{sha256}, despliegue reproducible en Render \\
    \bottomrule
  \end{tabular}
\end{table}

\chapter{Implementaci\'on}
\label{cap:implementacion}

Este cap\'itulo describe, a nivel de dise\~no y sin degenerar en un listado
de c\'odigo fuente, las implementaciones principales de BIOPET v1.0.0.

\section{Autenticaci\'on y autorizaci\'on}

El backend emite JWT firmados con HMAC-SHA256 (\texttt{jjwt}~0.12.6),
transportados en cookies \texttt{access\_token}/\texttt{refresh\_token}
(\texttt{HttpOnly}+\texttt{Secure}+\texttt{SameSite=Strict}). Cada token
contiene diez \emph{claims}: los siete est\'andar de RFC~7519~\cite{rfc7519}
(\texttt{iss}, \texttt{sub}, \texttt{aud}, \texttt{iat}, \texttt{nbf},
\texttt{exp}, \texttt{jti}) y tres personalizados (\texttt{email},
\texttt{rol}, \texttt{typ}). La autorizaci\'on combina rol
(\texttt{@PreAuthorize}) y propiedad (verificaci\'on expl\'icita de que un
\texttt{ROLE\_DUENO} solo accede a sus propias mascotas). El logout es
idempotente y revoca el token en Redis mediante su \texttt{jti}. Detalle
completo en el cap\'itulo~\ref{cap:seguridad}.

\section{Gesti\'on de usuarios}

\texttt{UsuarioController} expone el perfil del usuario autenticado
(\texttt{GET /api/usuarios/me}). El registro (\texttt{POST
/api/auth/registro}) crea siempre un usuario con \texttt{ROLE\_DUENO}; los
dem\'as roles (\texttt{ADMIN}, \texttt{VETERINARIO}, \texttt{AUXILIAR}) se
asignan mediante el usuario administrador sembrado o por un administrador
ya autenticado, no de forma autoasignada por el registro p\'ublico.

\section{Gesti\'on de mascotas}

CRUD completo (\texttt{GET}/\texttt{POST}/\texttt{PUT}/\texttt{DELETE}
sobre \texttt{/api/mascotas} y \texttt{/api/mascotas/\{id\}}), con
eliminaci\'on l\'ogica (\texttt{activo=false}), no f\'isica, y un endpoint
de resumen por especie respaldado por \texttt{fn\_resumen\_mascotas\_por\_especie}.

\section{Citas}

\texttt{CitaController} gestiona la programaci\'on de citas veterinarias
(alta, consulta, cambio de estado). La actualizaci\'on masiva del estado de
citas vencidas de un veterinario se resuelve con la rutina
\texttt{sp\_actualizar\_estado\_citas\_masivas} (transacci\'on \'unica en
PostgreSQL, no un bucle de actualizaciones individuales desde el backend).

\section{Consultas}

\texttt{ConsultaController} registra atenciones veterinarias. El registro
usa la rutina \texttt{sp\_registrar\_consulta\_validada}, que valida en la
propia base de datos --- antes de insertar --- que la mascota exista y
est\'e activa, y que el veterinario referenciado tenga rol
\texttt{VETERINARIO} o \texttt{ADMIN} y est\'e activo, lanzando
\texttt{RAISE EXCEPTION} si no, en vez de delegar esa validaci\'on
\'unicamente a la capa de servicio de Java.

\section{Vacunas}

\texttt{VacunaController} registra y consulta vacunas aplicadas a una
mascota, insumo de \texttt{fn\_historial\_clinico\_mascota} para mostrar la
\'ultima y la pr\'oxima vacuna de una mascota en el resumen consolidado.

\section{Integraci\'on con API externa}

\texttt{ExternalApiController} consume un servicio externo desde el
backend. Se documenta como \textbf{implementado}, no \textbf{verificado}:
no cuenta con una prueba automatizada dedicada a la fecha de este informe
(secci\'on~\ref{sec:req-unidad4}).

\section{Cach\'e (patr\'on cache-aside)}

\texttt{MascotaService} usa la abstracci\'on de cach\'e de Spring
(\texttt{@Cacheable} para lecturas, \texttt{@CacheEvict} para
invalidaci\'on tras escritura) contra Redis/Valkey. El TTL de cach\'e se
externaliza por variable de entorno, no est\'a codificado en el fuente.
El listado~\ref{lst:cache-aside} muestra el fragmento real de
\texttt{MascotaService} que implementa este patr\'on: el listado paginado
se cachea con una clave compuesta (correo, p\'agina, tama\~no, orden) para
que cada combinaci\'on de par\'ametros tenga su propia entrada, y toda
operaci\'on de escritura invalida el cach\'e completo del listado.

\begin{lstlisting}[language=Java,caption={Patr\'on \emph{cache-aside} declarativo en \texttt{MascotaService.java} (fragmento real).},label={lst:cache-aside}]
@Cacheable(value = "mascotas",
    key = "#email + '-' + #pageable.pageNumber + '-' +
           #pageable.pageSize + '-' + #pageable.sort.toString()")
@Transactional(readOnly = true)
public Page<MascotaResponse> listar(Pageable pageable, String email) {
    Usuario usuario = usuarioRepository
        .findByEmailAndActivoTrue(email)
        .orElseThrow(() -> new RecursoNoEncontradoException(
            "Usuario no encontrado"));

    if (usuario.getRol() == Rol.ROLE_DUENO) {
        return mascotaRepository
            .findAllByDuenioIdAndActivoTrue(usuario.getId(), pageable)
            .map(this::toResponse);
    }
    return mascotaRepository.findAllByActivoTrue(pageable)
        .map(this::toResponse);
}

@CacheEvict(value = "mascotas", allEntries = true)
@Transactional
public MascotaResponse crear(MascotaRequest request) {
    // ... validacion y persistencia (omitida por brevedad)
}
\end{lstlisting}

\section{Procedimientos almacenados}

Ver cap\'itulo~\ref{cap:arquitectura}, secci\'on~\ref{sec:acceso-jpa-formal},
para el cat\'alogo completo y la estrategia de migraci\'on que preserva
compatibilidad con bases ya migradas. El listado~\ref{lst:procedure-jpa}
muestra el repositorio dedicado a la invocaci\'on formal de las seis
rutinas almacenadas: cada m\'etodo usa \texttt{@Procedure}, sin ning\'un
CRUD gen\'erico expuesto (el repositorio extiende \texttt{Repository}, no
\texttt{JpaRepository}, deliberadamente).

\begin{lstlisting}[language=Java,caption={Acceso JPA formal a procedimientos almacenados en \texttt{ProcedimientoBiopetRepository.java} (fragmento real).},label={lst:procedure-jpa}]
public interface ProcedimientoBiopetRepository
        extends Repository<Mascota, Long> {

    @Procedure(name = "fn_resumen_mascotas_por_especie")
    List<ResumenEspecie> resumenPorEspecie(Long duenioId);

    @Procedure(procedureName = "sp_registrar_consulta_validada",
               outputParameterName = "p_consulta_id")
    Long registrarConsultaValidada(
        @Param("p_mascota_id") Long mascotaId,
        @Param("p_veterinario_id") Long veterinarioId,
        @Param("p_motivo") String motivo,
        @Param("p_diagnostico") String diagnostico,
        @Param("p_tratamiento") String tratamiento,
        @Param("p_observaciones") String observaciones);
}
\end{lstlisting}

\section{Migraciones Flyway}
\label{sec:migraciones-flyway}

El esquema se gestiona \textbf{exclusivamente} con Flyway como \'unica
fuente de verdad, aplicando en orden \texttt{V1\_\_schema\_inicial.sql},
\texttt{V2\_\_citas.sql}, \texttt{V3\_\_consultas.sql},
\texttt{V4\_\_vacunas.sql}, \texttt{V5\_\_procedimientos\_biopet.sql} y
\texttt{V6\_\_formalizar\_procedimientos\_jpa.sql}, verificado desde un
volumen de PostgreSQL completamente vac\'io. La arquitectura de arranque
reproducible en Docker es, en su versi\'on final v1.0.0:

\begin{enumerate}[nosep]
  \item PostgreSQL inicia desde un volumen vac\'io.
  \item \texttt{docker-entrypoint-initdb.d} ejecuta \'unicamente
    \texttt{db/roles-bootstrap.sql}: un \emph{bootstrap} m\'inimo que crea
    el rol \texttt{biopet\_app} \textbf{sin ninguna dependencia de objetos
    de esquema} (sin tablas, secuencias, funciones ni procedimientos),
    porque \texttt{application.yml} configura ese rol como el usuario del
    \emph{pool} principal de Spring (Hikari/JPA) y debe existir antes de
    que el contenedor del backend intente conectarse, independientemente
    del estado de Flyway.
  \item Flyway ejecuta realmente las migraciones \texttt{V1} a \texttt{V6}
    en orden, sobre el esquema todav\'ia vac\'io (sin
    \texttt{baseline-on-migrate} enmascarando un esquema preexistente).
  \item \texttt{afterMigrate.sql} (\emph{callback} de Flyway) otorga a
    \texttt{biopet\_app} los privilegios sobre los objetos de esquema
    reci\'en creados por \texttt{V1--V6}, una vez que esos objetos ya
    existen.
  \item \texttt{DataInitializer} (\texttt{CommandLineRunner
    seedAdmin}) siembra el usuario administrador de forma idempotente
    \textbf{despu\'es} de que Flyway termina, dentro del propio ciclo de
    arranque de Spring Boot, no como script de Postgres.
  \item \texttt{db/schema.sql}, \texttt{db/seed.sql}, \texttt{db/roles.sql}
    y \texttt{db/procs/} \textbf{no se ejecutan autom\'aticamente} en este
    flujo: se conservan en el repositorio como cat\'alogo, evidencia
    hist\'orica y documentaci\'on de referencia (\texttt{docs/basedatos/CATALOGOSP.md}),
    no como parte activa del \emph{bootstrap} de Docker.
\end{enumerate}

Esta arquitectura corrige un dise\~no anterior en el que \texttt{db/schema.sql}
duplicaba \texttt{V1} y, montado en \texttt{docker-entrypoint-initdb.d} junto
con \texttt{spring.flyway.baseline-on-migrate}, enmascaraba la migraci\'on
real detr\'as de un baseline sobre un volumen que aparentaba ser nuevo. La
versi\'on final v1.0.0 elimina esa duplicidad: Flyway es la \'unica v\'ia
por la que se crea el esquema, verificable ejecutando \texttt{make up}
desde una clonaci\'on limpia y observando \texttt{V1} a \texttt{V6}
aplicarse en orden en los logs del backend.

\section{Manejo uniforme de errores (RFC~7807)}

Todos los errores de la API usan el formato \texttt{application/problem+json}
(RFC~7807~\cite{rfc7807}), con \texttt{type}, \texttt{title},
\texttt{status}, \texttt{detail} e \texttt{instance}
(\texttt{GlobalExceptionHandler}, \texttt{ProblemDetailFactory}). El
\texttt{detail} nunca incluye \emph{stack trace} ni el valor bruto de un
par\'ametro inv\'alido. El listado~\ref{lst:problem-detail} muestra la
f\'abrica real que construye cada \texttt{ProblemDetail}, reutilizada por
todos los manejadores de excepci\'on del backend.

\begin{lstlisting}[language=Java,caption={F\'abrica de errores RFC~7807 en \texttt{ProblemDetailFactory.java} (fragmento real).},label={lst:problem-detail}]
public static ProblemDetail build(HttpStatus status, ProblemType type,
        String title, String detail, HttpServletRequest request) {
    ProblemDetail problemDetail =
        ProblemDetail.forStatusAndDetail(status, detail);
    problemDetail.setType(type.uri());
    problemDetail.setTitle(title);
    problemDetail.setInstance(URI.create(request.getRequestURI()));
    return problemDetail;
}
\end{lstlisting}

\section{Frontend}

Aplicaci\'on Angular~17.3 con componentes \emph{standalone} y formularios
reactivos. \rutalarga{core/auth.service.ts}/\rutalarga{auth.guard.ts} gestionan
la sesi\'on por cookies (\texttt{withCredentials: true}, sin JWT en
\texttt{localStorage}); \rutalarga{core/http-error.interceptor.ts} traduce el
\texttt{ProblemDetail} del backend a mensajes legibles y reintenta
\texttt{refresh} ante un 401 en rutas protegidas; \rutalarga{features/login.component.ts}
y \rutalarga{features/mascotas.component.ts} usan atributos ARIA
(\texttt{aria-invalid}, \texttt{aria-describedby}, \texttt{role="alert"},
\texttt{aria-live}) para anunciar errores a lectores de pantalla, en
l\'inea con WCAG~2.1~\cite{wcag21}.

% EVIDENCIA-ZAIDA-01
% Ruta esperada: figuras/zaida/01-frontend-login.png
\evidencia{figuras/zaida/01-frontend-login.png}%
{Pantalla de inicio de sesi\'on del frontend de BIOPET.}%
{fig:frontend-login-final}

% EVIDENCIA-ZAIDA-02
% Ruta esperada: figuras/zaida/02-crud-mascotas.png
\evidencia{figuras/zaida/02-crud-mascotas.png}%
{CRUD de mascotas en el frontend: listado, alta y edici\'on.}%
{fig:crud-mascotas-final}

\section{Persistencia}

Spring Data JPA sobre PostgreSQL, con consultas derivadas para
\texttt{UsuarioRepository} y \texttt{MascotaRepository}, y acceso formal
(\texttt{@Procedure}/\texttt{@NamedStoredProcedureQuery}) para las seis
rutinas almacenadas descritas en el cap\'itulo~\ref{cap:arquitectura}.
Hibernate/JPA se conecta en tiempo de ejecuci\'on con la cuenta de
privilegios m\'inimos \texttt{biopet\_app}; Flyway usa la cuenta propietaria
\texttt{biopet\_user}, separando la conexi\'on que aplica migraciones de la
que sirve tr\'afico de la aplicaci\'on.

\chapter{Pruebas y calidad}
\label{cap:pruebas-calidad}

\section{Pruebas automatizadas del backend}

La suite completa del backend, ejecutada con \texttt{mvn clean verify}
sobre el commit exacto del tag hist\'orico e inmutable \texttt{v1.0.0}
(\texttt{0d5cd52}) y verificada directamente contra
\texttt{Backend/target/surefire-reports/} de esa reproducci\'on
(\texttt{docs/mediciones/sec/reproduccion-v1.0.0/}), reporta:

\begin{table}[H]
  \centering
  \caption{Resultado real de la suite de pruebas del backend sobre el commit del tag \texttt{v1.0.0} (evidencia hist\'orica; ver estado del \texttt{HEAD} actual m\'as abajo).}
  \label{tab:pruebas-final}
  \begin{tabular}{@{}ll@{}}
    \toprule
    M\'etrica & Valor \\
    \midrule
    Pruebas ejecutadas & 205 \\
    Fallos & 0 \\
    Errores & 0 \\
    Omitidas & 0 \\
    Resultado del build & \texttt{BUILD SUCCESS} \\
    \bottomrule
  \end{tabular}
\end{table}

La figura~\ref{fig:maven-verify-final} resume este resultado, regenerada
autom\'aticamente desde ese mismo log crudo versionado mediante
\texttt{docs/informe/scripts/generar-figuras-evidencia.py}, sin capturas
manuales.

% EVIDENCIA-JAIME-01 (regenerada automaticamente; sustituye a la figura
% historica figuras/jaime/01-maven-verify.png, que se conserva en el
% repositorio como artefacto historico pero ya no se referencia aqui)
% Ruta esperada: figuras/jaime/07-maven-verify-final.png
\evidencia{figuras/jaime/07-maven-verify-final.png}%
{Resultado real de \texttt{mvn clean verify} sobre el commit del tag \texttt{v1.0.0}: 205 pruebas, 0 fallos, 0 errores, 0 omitidas, \texttt{BUILD SUCCESS} (regenerada por \texttt{docs/informe/scripts/generar-figuras-evidencia.py} desde el log crudo versionado).}%
{fig:maven-verify-final}

Estas 205 pruebas corresponden al estado hist\'orico e inmutable del tag
\texttt{v1.0.0}. Ejecutar \texttt{mvn clean verify} sobre el
\texttt{HEAD} actual de la rama de recalificaci\'on -- que incorpora
pruebas nuevas sobre esa misma base, sin sustituirla -- produce
\textbf{217} pruebas (0 fallos, 0 errores, 0 omitidas),
\texttt{BUILD SUCCESS}; ambas cifras son correctas dentro de su alcance
temporal respectivo. Detalle de la diferencia (205~$\rightarrow$~217) y
reproducci\'on de ambos estados en
\texttt{docs/mediciones/TEST-COUNT-PROVENANCE.md} y en la tabla de
\texttt{README.md} (secci\'on ``Estado t\'ecnico verificado'').

\section{Cobertura JaCoCo}
\label{sec:jacoco-final}

JaCoCo~\cite{jacoco-docs} mide y verifica autom\'aticamente la cobertura
con una regla \texttt{BUNDLE} sobre todo el m\'odulo backend. El umbral
autom\'atico configurado en \texttt{Backend/pom.xml} se elev\'o de 60\,\%
(Tercera Entrega) a \textbf{70\,\%} para v1.0.0.

\begin{table}[H]
  \centering
  \caption{Cobertura medida por \texttt{jacoco:check} sobre el commit del tag \texttt{v1.0.0} (evidencia hist\'orica; fuente: \texttt{docs/mediciones/jacoco/METRICS.md}). Ver estado del \texttt{HEAD} actual m\'as abajo.}
  \label{tab:jacoco-final}
  \begin{tabular}{@{}lrrrr@{}}
    \toprule
    Contador & Cubierto & No cubierto & Cobertura & Umbral \\
    \midrule
    LINE & 885 & 79 & 91{,}80\,\% & 70\,\% \\
    BRANCH & 181 & 47 & 79{,}39\,\% & 70\,\% \\
    COMPLEXITY & 299 & 77 & 79{,}52\,\% & 60\,\% \\
    \bottomrule
  \end{tabular}
\end{table}

La figura~\ref{fig:jacoco-resumen-final} muestra LINE y BRANCH frente al
umbral de 70\,\% (COMPLEXITY, con umbral distinto de 60\,\%, no se
incluye), regenerada autom\'aticamente desde \texttt{jacoco.csv} (45
clases) mediante el mismo script.

% EVIDENCIA-JAIME-02 (regenerada automaticamente; sustituye a la figura
% historica figuras/jaime/02-jacoco-resumen.png, que se conserva en el
% repositorio como artefacto historico pero ya no se referencia aqui)
% Ruta esperada: figuras/jaime/08-jacoco-resumen-final.png
\evidencia{figuras/jaime/08-jacoco-resumen-final.png}%
{Cobertura JaCoCo real sobre las 45 clases del backend: LINE 91{,}80\,\% y BRANCH 79{,}39\,\%, ambas por encima del umbral de calidad de 70\,\% (regenerada por \texttt{docs/informe/scripts/generar-figuras-evidencia.py} desde \texttt{jacoco.csv}).}%
{fig:jacoco-resumen-final}

Los tres contadores superan con margen su umbral configurado. La cobertura
mide ejecuci\'on de pruebas, no ausencia de defectos: una l\'inea cubierta
por una prueba no implica que est\'e libre de errores l\'ogicos no
contemplados por la aserci\'on de esa prueba.

Esta cobertura (LINE 91{,}80\,\%, BRANCH 79{,}39\,\%) corresponde al
estado hist\'orico e inmutable del tag \texttt{v1.0.0}. Recalculada sobre
el \texttt{HEAD} actual de la rama de recalificaci\'on, \texttt{jacoco.csv}
arroja LINE~\textbf{90{,}22\,\%} (876/971) y BRANCH~\textbf{76{,}27\,\%}
(180/236); ambas siguen superando con margen el umbral de 70\,\% del
\texttt{pom.xml}. Las diferencias corresponden al estado actual del
c\'odigo y de su suite de pruebas respecto del corte hist\'orico
\texttt{v1.0.0} (detalle en \texttt{README.md}, secci\'on ``Estado
t\'ecnico verificado'').

\section{Rendimiento (k6)}
\label{sec:rendimiento-final}

Herramienta de carga: k6~\cite{k6-docs}. La metodolog\'ia sigue el mismo
patr\'on de \emph{benchmarking} t\'ecnico contra un artefacto web real
descrito por Hasselbring~\cite{hasselbring2021benchmarking} como est\'andar
emp\'irico de ingenier\'ia de software, y usa los percentiles p95/p99 como
marco de referencia para interpretar el rendimiento de aplicaciones web
de negocio, en la l\'inea propuesta por Rempel~\cite{rempel2015webperf}.
Para v1.0.0 se ampli\'o el n\'umero de corridas de 3+3 (Tercera Entrega) a
\textbf{5 en fr\'io + 5 en caliente}, contra \texttt{GET /api/mascotas}
sobre HTTPS/TLS~1.3 real, 50~VUs, \emph{ramp-up} de 5~s y carga sostenida
de 30~s por corrida.

\begin{table}[H]
  \centering
  \tiny
  \caption{Resultado real de las diez corridas de k6 (fuente: \texttt{docs/mediciones/perf/REPORT.md}).}
  \label{tab:k6-final}
  \begin{tabular}{@{}lrrrrrrr@{}}
    \toprule
    Corrida & n & Media (ms) & Mediana (ms) & p95 (ms) & p99 (ms) & Error (\%) & Throughput (req/s) \\
    \midrule
    caliente-01 & 3196 & 10{,}54 & 8{,}68 & 20{,}15 & 30{,}68 & 0{,}0 & 91{,}17 \\
    caliente-02 & 3226 & 7{,}19 & 6{,}63 & 9{,}70 & 17{,}18 & 0{,}0 & 92{,}09 \\
    caliente-03 & 3226 & 6{,}60 & 6{,}01 & 10{,}86 & 16{,}13 & 0{,}0 & 92{,}13 \\
    caliente-04 & 3236 & 5{,}21 & 4{,}92 & 6{,}83 & 10{,}12 & 0{,}0 & 92{,}40 \\
    caliente-05 & 3236 & 5{,}40 & 4{,}96 & 7{,}81 & 12{,}59 & 0{,}0 & 92{,}41 \\
    fr\'io-01 & 3206 & 10{,}51 & 8{,}47 & 21{,}10 & 33{,}14 & 0{,}0 & 91{,}49 \\
    fr\'io-02 & 3211 & 9{,}37 & 7{,}60 & 17{,}23 & 25{,}79 & 0{,}0 & 91{,}62 \\
    fr\'io-03 & 3204 & 10{,}80 & 8{,}43 & 22{,}13 & 33{,}40 & 0{,}0 & 91{,}36 \\
    fr\'io-04 & 3206 & 10{,}71 & 8{,}50 & 20{,}16 & 32{,}65 & 0{,}0 & 91{,}45 \\
    fr\'io-05 & 3199 & 10{,}61 & 7{,}88 & 20{,}71 & 40{,}87 & 0{,}0 & 91{,}51 \\
    \bottomrule
  \end{tabular}
\end{table}

La tasa de error fue \textbf{0{,}0\,\%} en las diez corridas, sin
excepci\'on. Intervalos de confianza al 95\,\% (distribuci\'on t de
Student) reportados por corrida en \texttt{docs/mediciones/perf/REPORT.md};
por ejemplo, caliente-01: $[10{,}30, 10{,}78]$~ms. Adem\'as del IC~95\,\%,
se aplic\'o la prueba de Wilcoxon pareada entre condiciones fr\'ia y
caliente (por \'indice de llegada), con tama\~no de efecto $r$ de
Rosenthal: el efecto es peque\~no en la corrida~01 ($r=-0{,}07$) y grande en
las corridas~03--04 ($r=-0{,}65$ a $-0{,}84$), consistente con una cach\'e
que reduce la latencia de forma creciente conforme el sistema se estabiliza
entre corridas consecutivas.

% EVIDENCIA-FRED-06
% Ruta esperada: figuras/fred/06-performance-report.png
\evidencia{figuras/fred/06-performance-report.png}%
{Vista del reporte agregado de rendimiento.}%
{fig:performance-report-final}

% EVIDENCIA-FRED-CALIENTE
% Ruta esperada: figuras/fred/Rendimiento-PruebaCaliente.png
\evidencia{figuras/fred/Rendimiento-PruebaCaliente.png}%
{Salida de consola de una corrida de k6 en caliente.}%
{fig:rendimiento-caliente-final}

% EVIDENCIA-FRED-FRIO
% Ruta esperada: figuras/fred/Rendimiento-PruebaFria.png
\evidencia{figuras/fred/Rendimiento-PruebaFria.png}%
{Salida de consola de una corrida de k6 en fr\'io.}%
{fig:rendimiento-frio-final}

\section{Usabilidad (SUS)}
\label{sec:sus-final}

El instrumento System Usability Scale de Brooke (1996)~\cite{sus-brooke1996},
alineado con los conceptos de usabilidad de
ISO~9241-11:2018~\cite{iso9241-11} (eficacia, eficiencia y satisfacci\'on
en un contexto de uso), sin modificar, se aplic\'o a \textbf{18
registros} (P01--P18, ampliado desde los 10 de la Tercera Entrega), con
una tarea de referencia com\'un (login, alta, edici\'on, eliminaci\'on
l\'ogica, logout). Los 18 registros y sus respuestas Q1--Q10 son
evidencia t\'ecnicamente verificable en el repositorio, matem\'aticamente
reproducible (\texttt{scripts/analisis-sus.py}); que corresponden a
participantes reales es una \textbf{declaraci\'on del equipo}. El
protocolo de \texttt{docs/etica/ETHICS.md} exige consentimiento informado
firmado por participante, pero una auditor\'ia documental posterior a
esta entrega constat\'o que esos formularios individuales no est\'an
actualmente disponibles como evidencia verificable; esta situaci\'on se
document\'o mediante una constancia de regularizaci\'on firmada por los
integrantes del proyecto
(\texttt{docs/etica/regularizacion-sus/CONSTANCIA-REGULARIZACION-SUS-BIOPET-2026-08-31.pdf}),
que \textbf{no sustituye} consentimientos individuales.

\begin{table}[H]
  \centering
  \caption{Resultados agregados del instrumento SUS (n = 18).}
  \label{tab:sus-final}
  \begin{tabular}{lr}
    \toprule
    M\'etrica & Valor \\
    \midrule
    Media (SUS Score) & 74{,}44 / 100 \\
    Desviaci\'on t\'ipica (muestral, $n-1$) & 22{,}35 \\
    Intervalo de confianza 95\,\% & $[63{,}33,\ 85{,}56]$ (margen $\pm 11{,}12$) \\
    Mediana (p50) & 82{,}50 \\
    M\'inimo / M\'aximo & 22{,}50 / 97{,}50 \\
    Clasificaci\'on cualitativa (Bangor et al.) & Bueno \\
    \bottomrule
  \end{tabular}
\end{table}

La figura~\ref{fig:sus-resultados-final} presenta la distribuci\'on de
los 18 puntajes SUS reales (diagrama de caja con los puntos individuales
superpuestos), regenerada autom\'aticamente desde \texttt{sus-raw.csv}
mediante el mismo script, sin captura manual.

% EVIDENCIA-JAIME-06 (regenerada automaticamente por Jaime a partir del
% hallazgo del docente sobre esta figura; sustituye, solo en esta
% referencia del informe, a la figura historica
% figuras/zaida/05-sus-resultados.png -- que se conserva SIN modificar en
% el repositorio como artefacto historico de Zaida, no se sobrescribe ni
% se borra)
% Ruta esperada: figuras/jaime/06-sus-resultados-final.png
\evidencia{figuras/jaime/06-sus-resultados-final.png}%
{Distribuci\'on de los 18 puntajes SUS reales: media 74{,}44, DT 22{,}35, mediana 82{,}50, IC 95\,\% $[63{,}33,\ 85{,}56]$ (regenerada por \texttt{docs/informe/scripts/generar-figuras-evidencia.py} desde \texttt{sus-raw.csv}).}%
{fig:sus-resultados-final}

\textbf{Interpretaci\'on.} La media obtenida (74{,}44) se ubica en la
categor\'ia \textbf{Bueno} de la escala de adjetivos SUS, con un IC~95\,\%
que incluye el umbral de referencia de 68 puntos considerado ``sobre el
promedio'' de la industria seg\'un el benchmark emp\'irico de Bangor,
Kortum \& Miller~\cite{bangor2008sus}. El participante con menor puntaje
(P08, 22{,}5) declar\'o experiencia web ``ninguna'', consistente con el
efecto conocido de curva de aprendizaje inicial en usuarios sin experiencia
digital previa. La ampliaci\'on de la muestra de $n=10$ a $n=18$ redujo el
margen del intervalo de confianza (de un margen mayor en la Tercera
Entrega a $\pm 11{,}12$ aqu\'i), pero \textbf{un tama\~no muestral mayor no
implica, por s\'i solo, representatividad externa}: los participantes
siguen siendo reclutados por conveniencia (secci\'on~\ref{cap:amenazas-limitaciones-final}).

\section{Calidad web autom\'atica (Lighthouse)}
\label{sec:lighthouse-final}

La configuraci\'on de auditor\'ia declara umbrales Performance~$\geq$~80,
Accessibility~$\geq$~90, Best Practices~$\geq$~90 y SEO~$\geq$~90~\cite{lighthouse-docs}.
El equipo ejecut\'o tres corridas sucesivas: la inicial del 2026-08-01
(solo perfil m\'ovil, contra \texttt{localhost}, SEO~82/100, umbral no
cumplido); una corrida completa el 2026-08-18
(\texttt{lighthouserc.js}/\texttt{lighthouserc.desktop.js}, tambi\'en
contra \texttt{localhost:4200}, ambos perfiles, SEO~100/100); y la
\textbf{corrida definitiva contra el despliegue p\'ublico de Render}
el \textbf{2026-09-03}
(\texttt{lighthouserc.render.js}/\texttt{lighthouserc.render.desktop.js},
\texttt{docs/mediciones/lighthouse/lhci-20260903-2102-render-*.json},
metadatos en \texttt{lhci-20260903-2102-render.meta.txt}; metodolog\'ia y
medias reproducibles en \texttt{docs/mediciones/lighthouse/RENDER-REPORT.md}),
la primera de las tres auditando \texttt{requestedUrl} p\'ublico
(\texttt{https://biopet-frontend.onrender.com}) en vez de \texttt{localhost}.
Cubre ambos perfiles (m\'ovil y escritorio), 3~corridas por ruta y por
perfil (12~corridas totales) sobre las rutas \texttt{/login} y
\texttt{/mascotas} (ninguna es la portada). El script
\texttt{scripts/run-lighthouse-render.sh} verifica, antes de auditar, que
el despliegue responda \texttt{200} en ambas rutas (sin despliegue en
curso), y valida el campo \texttt{requestedUrl} de cada JSON generado
contra el dominio p\'ublico antes de archivarlo.

\begin{table}[H]
  \centering
  \caption{Resultado real de la corrida Lighthouse del 2026-09-03 contra el despliegue p\'ublico de Render (12 corridas, mobile+desktop, medias calculadas por \texttt{compute-render-averages.mjs}; fuente: \texttt{docs/mediciones/lighthouse/lhci-20260903-2102-render-*.json}).}
  \label{tab:lighthouse-final}
  \begin{tabular}{@{}llrrl@{}}
    \toprule
    Categor\'ia & Perfil & Ruta & Media (0--100) & Umbral \\
    \midrule
    Performance & Escritorio & /login & 99{,}7 & $\geq$ 80 --- Cumplido \\
    Performance & Escritorio & /mascotas & 100{,}0 & $\geq$ 80 --- Cumplido \\
    Performance & M\'ovil & /login & 100{,}0 & $\geq$ 80 --- Cumplido \\
    Performance & M\'ovil & /mascotas & 99{,}0 & $\geq$ 80 --- Cumplido \\
    Accessibility & Ambos & Ambas & 91{,}0 & $\geq$ 90 --- Cumplido \\
    Best Practices & Ambos & /login & 100{,}0 & $\geq$ 90 --- Cumplido \\
    Best Practices & Ambos & /mascotas & 96{,}0 & $\geq$ 90 --- Cumplido \\
    SEO & Ambos & Ambas & \textbf{100{,}0} & $\geq$ 90 --- \textbf{Cumplido} \\
    \bottomrule
  \end{tabular}
\end{table}

% EVIDENCIA-ZAIDA-04
% Ruta esperada: figuras/zaida/04-lighthouse.png
\evidencia{figuras/zaida/04-lighthouse.png}%
{Reporte de Lighthouse CI del frontend de BIOPET contra el despliegue p\'ublico de Render (2026-09-03), por perfil y ruta.}%
{fig:lighthouse-final}

\textbf{Estado real para v1.0.0.} Las \textbf{cuatro} categor\'ias cumplen
su umbral en las 12 corridas del 2026-09-03, en ambos perfiles y ambas
rutas, contra el dominio p\'ublico real (\texttt{requestedUrl} y
\texttt{finalUrl} de los 12 JSON coinciden con
\texttt{https://biopet-frontend.onrender.com/login} o
\texttt{.../mascotas}, sin redirecci\'on). Performance se mantiene en el
rango 99{,}0--100{,}0 en los cuatro grupos perfil/ruta, mejorando el rango
m\'ovil 90--94 de la corrida local del 2026-08-18. Accessibility (91{,}0),
Best Practices (96{,}0--100{,}0) y SEO (100{,}0) se mantienen id\'enticos a
la corrida local previa. Las corridas del 2026-08-01 y 2026-08-18 quedan
como evidencia hist\'orica del proceso de mejora (SEO corregido de 82 a
100) y se conservan sin modificar en
\texttt{docs/mediciones/lighthouse/}; la corrida del 2026-09-03 contra
Render es la que se reporta como definitiva por ser la \'unica auditada
contra el sistema realmente evaluable en producci\'on.

\section{Diferenciaci\'on de los tipos de medici\'on}

Este informe distingue expl\'icitamente: rendimiento (k6, agregado por
corrida, IC~95\,\% con distribuci\'on t), cach\'e (verificaci\'on aislada
por clave con \texttt{redis-cli}), cobertura (JaCoCo, LINE/BRANCH/COMPLEXITY
sobre el \texttt{BUNDLE} completo), usabilidad (SUS, ejecutada, $n=18$),
calidad web (Lighthouse, ejecutada en ambos perfiles mobile y desktop, 12
corridas, cuatro categor\'ias cumplidas), y seguridad (revisi\'on manual OWASP +
an\'alisis din\'amico ZAP + an\'alisis est\'atico SpotBugs/Find Security
Bugs, cap\'itulo~\ref{cap:seguridad}), para evitar mezclar cifras de
naturaleza distinta.

\chapter{Seguridad}
\label{cap:seguridad}

Este cap\'itulo integra la auditor\'ia de seguridad de BIOPET desde tres
\'angulos complementarios: revisi\'on manual sobre OWASP Top~10~\cite{owasp-top10},
an\'alisis din\'amico (OWASP ZAP Baseline) y an\'alisis est\'atico
(SpotBugs + Find Security Bugs), m\'as la auditor\'ia dedicada de SQL
din\'amico en procedimientos almacenados.

\section{Autenticaci\'on por cookies y revocaci\'on}

JWT firmado con HMAC-SHA256, transportado en cookies
\texttt{HttpOnly}+\texttt{Secure}+\texttt{SameSite=Strict}
(\texttt{access\_token}, \texttt{refresh\_token}), con revocaci\'on real
v\'ia lista negra en Redis indexada por \texttt{jti} (ADR-003), rate
limiting de login en memoria (5 fallos consecutivos $\to$ 401; el sexto
$\to$ 429 con cabecera \texttt{Retry-After}) y auditor\'ia estructurada
(\texttt{AUTH\_AUDIT}) que solo acepta \texttt{ip} y \texttt{subject} como
par\'ametros, sin v\'ia de c\'odigo para que una contrase\~na, un JWT
completo o un encabezado \texttt{Authorization} llegue al logger.

\section{Autorizaci\'on por rol y por propiedad}

Autorizaci\'on combinada: por rol (\texttt{@PreAuthorize}) y por propiedad
(\texttt{ROLE\_DUENO} solo accede a sus propias mascotas). 401 para
autenticaci\'on ausente/inv\'alida, 403 para autenticaci\'on v\'alida sin
autorizaci\'on, ambos como \texttt{ProblemDetail} sin \emph{stack trace}
(figura~\ref{fig:postman-401-final}).

% EVIDENCIA-JAIME-05
% Ruta esperada: figuras/jaime/05-postman-401.png
\evidencia{figuras/jaime/05-postman-401.png}%
{Respuesta 401 (\texttt{ProblemDetail}) capturada desde la colecci\'on Postman de BIOPET.}%
{fig:postman-401-final}

\section{Cabeceras de seguridad y TLS}

\texttt{SecurityConfig} (Spring Security~\cite{spring-security-docs}) fija
cinco cabeceras en todas las respuestas: \texttt{X-Frame-Options: DENY},
\texttt{X-Content-Type-Options: nosniff},
\texttt{Content-Security-Policy: default-src 'self'; frame-ancestors 'none'; object-src 'none'},
\texttt{Referrer-Policy: no-referrer}, y \texttt{Strict-Transport-Security}
solo sobre HTTPS. TLS~1.3 verificado en vivo
(\texttt{TLS\_AES\_256\_GCM\_SHA384}, figura~\ref{fig:tls-openssl-final})
en el entorno de desarrollo; en
producci\'on, HTTPS es provisto autom\'aticamente por Render sobre el
subdominio \texttt{*.onrender.com} (cap\'itulo~\ref{cap:despliegue}).

% EVIDENCIA-JAIME-03
% Ruta esperada: figuras/jaime/03-tls-openssl.png
\evidencia{figuras/jaime/03-tls-openssl.png}%
{Salida de \texttt{openssl s\_client -tls1\_3}: TLSv1.3 y \texttt{TLS\_AES\_256\_GCM\_SHA384} negociados (entorno de desarrollo).}%
{fig:tls-openssl-final}

% EVIDENCIA-JAIME-04
% Ruta esperada: figuras/jaime/04-security-headers.png
\evidencia{figuras/jaime/04-security-headers.png}%
{Cabeceras de seguridad reales capturadas con \texttt{curl.exe -i}.}%
{fig:security-headers-final}

\section{Auditor\'ia OWASP Top~10 (resumen)}

La tabla~\ref{tab:owasp-final} resume el resultado por control de la
revisi\'on manual sobre los nueve controles OWASP Top~10 con evidencia
real y verificable.

\begin{table}[H]
  \centering
  \caption{Resultado de la auditor\'ia manual OWASP por control (fuente: \rutalarga{docs/mediciones/sec/REPORT.md}).}
  \label{tab:owasp-final}
  \begin{tabular}{@{}ll@{}}
    \toprule
    Control OWASP & Resultado \\
    \midrule
    A01 --- Control de acceso roto & PASS \\
    A02 --- Fallas criptogr\'aficas (TLS) & PASS \\
    A03 --- Inyecci\'on & PASS \\
    A04 --- Dise\~no inseguro & Evaluado (cierre 2026-08-09) \\
    A05 --- Configuraci\'on de seguridad & PASS \\
    A06 --- Componentes vulnerables/desactualizados & Evaluado, con hallazgos documentados \\
    A07 --- Identificaci\'on y autenticaci\'on & PASS \\
    A09 --- Registro y monitoreo & PASS \\
    XSS (transversal) & Evaluado (cierre 2026-08-09) \\
    \bottomrule
  \end{tabular}
\end{table}

A08 y A10 no fueron objeto de auditor\'ia dedicada y no se documentan aqu\'i
para no inventar evidencia inexistente.

\section{Protecci\'on estructural contra inyecci\'on (A03)}

Todos los m\'etodos de repositorio son consultas derivadas de Spring Data o
invocaciones formales de procedimientos almacenados
(\texttt{@Procedure}/\texttt{@NamedStoredProcedureQuery}), con par\'ametros
enlazados y fuertemente tipados. Los payloads de inyecci\'on probados
(\texttt{1 OR 1=1}, \texttt{1; DROP TABLE usuarios; --}, \texttt{\%' UNION
SELECT NULL --}) son rechazados en la capa de conversi\'on de Spring MVC
antes de alcanzar el repositorio, verificado en \texttt{SqlInjectionSecurityTest}.

\section{Auditor\'ia est\'atica: SpotBugs + Find Security Bugs}
\label{sec:spotbugs}

Ejecutada con \texttt{spotbugs-maven-plugin}~4.10.3.0 + Find Security
Bugs~1.14.0, \texttt{effort=Max}, \texttt{threshold=Low} (m\'axima
exhaustividad), sobre todas las clases compiladas del backend; la
tabla~\ref{tab:spotbugs-final} resume los hallazgos por severidad.

\begin{table}[H]
  \centering
  \caption{Resumen del an\'alisis est\'atico (fuente: \rutalarga{docs/mediciones/sec/static-analysis/README.md}).}
  \label{tab:spotbugs-final}
  \begin{tabular}{@{}lr@{}}
    \toprule
    M\'etrica & Valor \\
    \midrule
    Hallazgos totales & 66 \\
    Severidad alta (\texttt{priority=1}) & 1 (\texttt{SPRING\_CSRF\_PROTECTION\_DISABLED}) \\
    Severidad media (\texttt{priority=2}) & 13 \\
    Severidad baja (\texttt{priority=3}) & 52 \\
    Hallazgos relacionados con SQL (\texttt{SQL\_*}) & \textbf{0} \\
    \bottomrule
  \end{tabular}
\end{table}

\subsection*{Tratamiento del hallazgo CSRF (\'unico de severidad alta)}
\label{sec:csrf}

\texttt{SecurityConfig.java} deshabilita expl\'icitamente la protecci\'on
CSRF integrada de Spring Security, lo que SpotBugs/Find Security Bugs
reporta correctamente como \texttt{SPRING\_CSRF\_PROTECTION\_DISABLED}
(\texttt{priority=1}). Este informe \textbf{no trata este hallazgo como un
falso positivo}: la herramienta detect\'o exactamente lo que hay en el
c\'odigo. Se trata, en cambio, como un \textbf{riesgo mitigado por
arquitectura}: la protecci\'on CSRF tradicional basada en tokens existe
para defender flujos de autenticaci\'on por \emph{sesi\'on con cookie
ambiental} en formularios HTML cl\'asicos, donde el navegador env\'ia la
cookie de sesi\'on autom\'aticamente en cualquier solicitud
\emph{cross-site}, incluida una iniciada por un sitio malicioso. BIOPET no
sigue ese patr\'on: el atributo \texttt{SameSite=Strict} en las cookies
\texttt{access\_token}/\texttt{refresh\_token} (ADR-006) impide que el
navegador adjunte esas cookies en solicitudes iniciadas desde un origen
distinto, que es precisamente el vector que CSRF explota. La combinaci\'on
\texttt{HttpOnly}+\texttt{Secure}+\texttt{SameSite=Strict} sustituye
funcionalmente al token CSRF cl\'asico para este esquema de autenticaci\'on
espec\'ifico, sin que eso implique que el riesgo no exist\'ia --- exist\'ia,
y se document\'o expl\'icitamente en vez de suprimirse. El gate de CI
\texttt{security-static} (secci\'on~\ref{sec:ci-jobs}) falla \'unicamente
ante hallazgos \texttt{SQL\_*}; el hallazgo CSRF se archiva siempre como
parte del reporte completo (sin filtros ni \texttt{@SuppressFBWarnings}),
para que quede visible en cada corrida en vez de ocultarse.

\section{Auditor\'ia din\'amica: OWASP ZAP}
\label{sec:zap}

Esta secci\'on documenta \textbf{dos} corridas de OWASP ZAP, claramente
diferenciadas: (A) el \emph{baseline} hist\'orico sin autenticaci\'on
(evidencia original de esta fase, conservada sin alterar), y (B) un
escaneo pasivo \textbf{autenticado}, incorporado en la recalificaci\'on
para cerrar la limitaci\'on que (A) ya declaraba expl\'icitamente. Ninguna
de las dos corridas es un \emph{active scan}: ambas son \emph{baseline}/pasivas,
sin intentos de explotaci\'on.

\subsection{A. Baseline hist\'orico (sin autenticaci\'on)}
\label{sec:zap-baseline-historico}

Ejecutada con \texttt{zap-baseline.py} (imagen oficial
\texttt{ghcr.io/zaproxy/zaproxy:stable}, versi\'on~2.17.0) contra el
backend real, dentro de la red Docker Compose (\texttt{http://backend:8080},
\textbf{sin TLS}, \textbf{sin autenticaci\'on}), el
\texttt{2026-08-18T02:53:21Z} (fuente:
\texttt{docs/mediciones/sec/zap/RUN-METADATA.txt}).

\begin{table}[H]
  \centering
  \caption{Resumen ZAP Baseline hist\'orico (fuente: \texttt{docs/mediciones/sec/zap/README.md}).}
  \label{tab:zap-final}
  \begin{tabular}{@{}lr@{}}
    \toprule
    Severidad & Alertas \\
    \midrule
    Alta & \textbf{0} \\
    Media & 0 \\
    Baja & 0 \\
    Informativa & 1 (\texttt{Non-Storable Content}, 3 instancias) \\
    \bottomrule
  \end{tabular}
\end{table}

C\'odigo de salida de \texttt{zap-baseline.py}: \texttt{0}. La
tabla~\ref{tab:zap-final} resume la severidad de las alertas; en t\'erminos
de cobertura real, el \emph{spider} no autenticado solo alcanz\'o
\textbf{3~URI} (\texttt{/}, \texttt{/robots.txt}, \texttt{/sitemap.xml} ---
las tres instancias de \texttt{Non-Storable Content} de la tabla) sobre
\textbf{2~endpoints \'unicos} seg\'un \texttt{insight.endpoint.total} del
propio JSON de ZAP (m\'etrica de ZAP que deduplica por endpoint, distinta
del conteo de URI), con \textbf{100\,\%} de las respuestas en \texttt{4xx}
(\texttt{insight.code.4xx}). El objetivo de la Entrega Final (cero
hallazgos de severidad alta) se cumple, pero esta corrida \textbf{no
constituye evidencia representativa de la superficie protegida}: el
spider de ZAP no fue autenticado y, por dise\~no, todo lo que no sea
\texttt{/api/auth/**} responde \texttt{401} (limitaci\'on declarada
expl\'icitamente en \texttt{docs/mediciones/sec/zap/README.md}, no
oculta).

\subsection{B. Recalificaci\'on --- escaneo pasivo autenticado (local, TLS)}
\label{sec:zap-authenticated-local}

Para cerrar la limitaci\'on de (A), se ejecut\'o un segundo escaneo con
OWASP ZAP~2.17.0 (misma imagen), esta vez \textbf{autenticado} y sobre
\textbf{HTTPS/TLS} real, contra el mismo backend pero levantado con el
perfil TLS local del propio repositorio
(\texttt{docker-compose.tls.yml}, \texttt{https://backend:8443}), el
\texttt{2026-09-03T21:38:48Z} (fuente:
\texttt{docs/mediciones/sec/zap/RUN-METADATA-AUTH-LOCAL.txt}). Sigue
siendo un entorno \textbf{local}, del propio repositorio acad\'emico ---
\textbf{no} se escane\'o el despliegue de Render ni ning\'un otro
repositorio (BIOPET-V2 es un proyecto separado, fuera de alcance).

La autenticaci\'on us\'o el m\'etodo \texttt{json} del \emph{Automation
Framework} de ZAP contra \texttt{POST~/api/auth/login}, con gesti\'on de
sesi\'on por \textbf{cookie} (m\'etodo nativo de ZAP, coincide con el
mecanismo real de \texttt{JwtCookieService}) y verificaci\'on peri\'odica
contra \texttt{GET~/api/usuarios/me}. El login y las siete rutas GET
protegidas objetivo (\texttt{/api/usuarios/me}, \texttt{/api/mascotas},
\texttt{/api/mascotas/resumen-especies}, \texttt{/api/citas},
\texttt{/api/consultas}, \texttt{/api/vacunas}, \texttt{/api/usuarios})
devolvieron exactamente \texttt{200}, verificado con la aserci\'on nativa
\texttt{responseCode: 200} de ZAP AF en cada petici\'on, con \textbf{0
diferencias} registradas.

Seg\'un los \emph{insights} agregados del propio JSON de ZAP,
\texttt{insight.endpoint.total = 9} (endpoints \'unicos, deduplicados);
el resumen de consola de \texttt{outputSummary} reporta, por separado,
``Total of 11 URLs'' --- un contador \textbf{distinto}, sin deduplicar
(incluye la repetici\'on de \texttt{POST~/api/auth/login}, invocado dos
veces por dise\~no del ciclo de verificaci\'on). \textbf{Estas dos cifras
no deben sumarse ni tratarse como la misma m\'etrica.} De igual modo,
\texttt{insight.code.2xx = 93} e \texttt{insight.code.4xx = 6} son
estad\'isticas agregadas internas de ZAP, \textbf{no} un desglose por
URL; ninguno de los reportes retenidos (HTML, XML ni JSON) expone el
denominador exacto de ese porcentaje, por lo que este informe
\textbf{no afirma} ``93\,\% de 11 peticiones'' ni ninguna fracci\'on
literal equivalente. Lo que s\'i qued\'o verificado por partida doble es
que ese \texttt{4xx} agregado \textbf{no corresponde a ninguna de las
siete rutas objetivo ni al login} (ambos confirmados en \texttt{200} por
la aserci\'on \texttt{responseCode} con 0 diferencias) y que una
reconciliaci\'on adicional de solo lectura, exportando el historial HTTP
real de una corrida equivalente, mostr\'o \textbf{0 respuestas 4xx} entre
las nueve peticiones expl\'icitas del plan. La identidad exacta de la
petici\'on individual que origin\'o ese \texttt{4xx} agregado no pudo
reconstruirse a partir de los artefactos retenidos --- se documenta aqu\'i
como \textbf{l\'imite de trazabilidad} de las plantillas de reporte de
ZAP, no como un hecho supuesto.

Alertas de esta corrida: \textbf{0} altas, \textbf{0} medias, \textbf{0}
bajas, \textbf{2 tipos informativos} (\texttt{Authentication Request
Identified} y \texttt{Session Management Response Identified}, 3
instancias en total, ambas sobre \texttt{POST~/api/auth/login}) --- son
ZAP reconociendo correctamente el propio flujo de autenticaci\'on, \textbf{no}
vulnerabilidades.

La tabla~\ref{tab:zap-comparativo} resume la comparaci\'on entre ambas
corridas.

\begin{table}[H]
  \centering
  \caption{Comparaci\'on entre el baseline hist\'orico sin autenticar y el escaneo pasivo autenticado local (TLS).}
  \label{tab:zap-comparativo}
  \begin{tabular}{@{}lll@{}}
    \toprule
    Criterio & A. Baseline hist\'orico & B. Autenticado local (TLS) \\
    \midrule
    Target & \texttt{http://backend:8080} & \texttt{https://backend:8443} \\
    TLS & No & S\'i (autofirmado, local) \\
    Autenticaci\'on & No & S\'i (\texttt{json}, cookie) \\
    Rutas protegidas en \texttt{2xx} & 0 & 7/7 \\
    Endpoints \'unicos & 2 & 9 \\
    Alta & 0 & 0 \\
    Media & 0 & 0 \\
    Baja & 0 & 0 \\
    Limitaci\'on principal & Sin cobertura autenticada & Solo GET, solo local \\
    \bottomrule
  \end{tabular}
\end{table}

\textbf{Alcance expl\'icito de esta corrida (B) --- lo que NO se hizo:}
no se ejecut\'o ning\'un \emph{active scan}; no se atacaron rutas
\texttt{POST}/\texttt{PUT}/\texttt{DELETE} (solo las siete rutas
\texttt{GET} listadas); no se cubri\'o el 100\,\% de los endpoints de la
API (quedan fuera, deliberadamente, todas las rutas de escritura y
\texttt{/api/externa/**}, un proxy a una API externa real); no se
escane\'o el despliegue de Render ni producci\'on; no se escane\'o
BIOPET-V2. Esta corrida \textbf{no valida la seguridad completa del
sistema}; mejora sustancialmente la representatividad del DAST porque
verifica autenticaci\'on real y rutas protegidas reales, sobre un
subconjunto acotado y deliberado de la superficie de la API. Extender
este mismo mecanismo de autenticaci\'on a un escaneo m\'as amplio, o
contra un entorno p\'ublico/de producci\'on, es una extensi\'on posible
para trabajo futuro que requerir\'ia adem\'as control expl\'icito de
credenciales y de la carga generada contra un servicio compartido.

Evidencia completa, trazable y verificable por SHA-256 en
\texttt{docs/mediciones/sec/zap/} (\texttt{zap-authenticated-local-report.json},
\texttt{.html}, \texttt{.xml}, \texttt{RUN-METADATA-AUTH-LOCAL.txt},
\texttt{zap-authenticated-local.yaml}, \texttt{SHA256SUMS-AUTH-LOCAL.txt})
y en \texttt{scripts/run-zap-authenticated-local.sh} (script reproducible,
sin credenciales hardcodeadas). Verificaci\'on:
\texttt{sha256sum -c SHA256SUMS-AUTH-LOCAL.txt} $\to$ \textbf{7/7 OK}.

\section{Auditor\'ia de SQL din\'amico en procedimientos almacenados}

\texttt{scripts/audit-sql-dynamic.sh} analiza los seis archivos de
\texttt{db/procs/*.sql} en busca de construcci\'on insegura de SQL
din\'amico (\texttt{EXECUTE} con concatenaci\'on, patrones ajenos a
PostgreSQL). Resultado: \textbf{0 hallazgos}, \texttt{exit 0}.

\section{Postman}

La colecci\'on \rutalarga{docs/postman/BIOPET.postman_collection.json}
(40+ \emph{requests}) y \rutalarga{docs/postman/BIOPET-Vacunas.postman_collection.json}
cubren estado del servicio, autenticaci\'on completa (incluida la
secuencia de rate limiting), mascotas, citas, consultas, vacunas y un
cat\'alogo de errores controlados (422, 400, 404, 401, 403, 429). Ninguna
depende de un JWT o cookie guardado manualmente en variable: se apoyan
exclusivamente en el \emph{cookie jar} autom\'atico de Postman.

\section{Limitaciones reales de este cap\'itulo}

\begin{itemize}[nosep]
  \item El certificado TLS del entorno de desarrollo es autofirmado y
    exclusivamente acad\'emico; en producci\'on, HTTPS lo provee Render de
    forma autom\'atica (no es un certificado gestionado por el equipo).
  \item El rate limiting de login es en memoria y por instancia: no es
    distribuido entre m\'ultiples r\'eplicas del backend.
  \item Los logs \texttt{AUTH\_AUDIT} se escriben \'unicamente en el log
    local del proceso; no existe integraci\'on con un SIEM centralizado.
  \item El \emph{spider} del baseline hist\'orico de ZAP
    (secci\'on~\ref{sec:zap-baseline-historico}) no fue autenticado; esa
    limitaci\'on espec\'ifica qued\'o parcialmente cerrada por el escaneo
    pasivo autenticado local (secci\'on~\ref{sec:zap-authenticated-local}),
    que verific\'o en \texttt{200} siete rutas \texttt{GET} protegidas
    reales de los m\'odulos de mascotas, citas, consultas, vacunas y
    usuarios --- pero sigue sin cubrir las rutas de escritura
    (\texttt{POST}/\texttt{PUT}/\texttt{DELETE}) de esos mismos m\'odulos,
    ni el despliegue de Render, ni un \emph{active scan}.
  \item El an\'alisis est\'atico no sustituye una revisi\'on de seguridad
    din\'amica dirigida a la librer\'ia JWT subyacente (\texttt{jjwt}); ese
    tipo de an\'alisis (\emph{fuzzing} dirigido, ver
    Yang et al.~\cite{yang2026tokentimebomb}, cap\'itulo~\ref{cap:trabajos-relacionados})
    queda fuera del alcance de esta auditor\'ia.
  \item ``0 alertas de severidad alta''/``0 hallazgos \texttt{SQL\_*}'' no
    implica ausencia total de vulnerabilidades: implica ausencia de los
    patrones que las reglas configuradas de ZAP y SpotBugs pueden detectar.
\end{itemize}

\chapter{Despliegue y reproducibilidad}
\label{cap:despliegue}

\section{Despliegue real en producci\'on}

BIOPET est\'a desplegado en \textbf{Render} (\texttt{render.com}), definido
declarativamente mediante \texttt{render.yaml} (patr\'on \emph{Blueprint}),
con HTTPS autom\'atico sobre subdominios \texttt{*.onrender.com}.

\begin{table}[H]
  \centering
  \caption{Servicios reales de producci\'on (fuente: \texttt{docs/despliegue/DEPLOYMENT.md}).}
  \label{tab:produccion-final}
  \begin{tabular}{@{}ll@{}}
    \toprule
    Servicio & URL / valor \\
    \midrule
    Frontend & \texttt{https://biopet-frontend.onrender.com} \\
    Backend & \texttt{https://biopet-backend-dh5e.onrender.com} \\
    \emph{Healthcheck} & \texttt{GET /actuator/health} $\to$ \texttt{\{"status":"UP","groups":["liveness","readiness"]\}} \\
    Base de datos & PostgreSQL~18 (gestionado por Render) \\
    Cach\'e & Valkey~8 (compatible con Redis, gestionado por Render) \\
    HTTPS & Activo, autom\'atico, sin configuraci\'on TLS manual en la app \\
    \bottomrule
  \end{tabular}
\end{table}

El backend distingue expl\'icitamente la variable \texttt{CORS\_ALLOWED\_ORIGINS}
(fijada al origen real del frontend) de la URL p\'ublica del propio backend,
evitando una configuraci\'on de CORS permisiva por omisi\'on. Las
credenciales de base de datos de producci\'on se inyectan por variables de
entorno gestionadas en el panel de Render, nunca versionadas en el
repositorio.

\section{Continuous Integration}
\label{sec:ci-jobs}

\texttt{.github/workflows/ci.yml} define seis \emph{jobs}, todos en verde a
la fecha de cierre de esta entrega:

\begin{table}[H]
  \centering
  \small
  \caption{Jobs del pipeline de CI.}
  \label{tab:ci-jobs}
  \begin{tabular}{@{}p{2.8cm}p{9.4cm}@{}}
    \toprule
    Job & Prop\'osito \\
    \midrule
    \texttt{backend-test} & \texttt{mvn clean verify}: 205 pruebas + \texttt{jacoco:check} (umbral 70\,\%) \\
    \texttt{frontend-build} & Build de producci\'on del frontend Angular \\
    \texttt{traceability} & \texttt{scripts/validate-traceability.sh}: consistencia SRS $\leftrightarrow$ matriz \\
    \texttt{sql-audit} & \texttt{scripts/audit-sql-dynamic.sh}: SQL din\'amico inseguro en \texttt{db/procs/*.sql} \\
    \texttt{security-static} & SpotBugs + Find Security Bugs; falla solo ante hallazgos \texttt{SQL\_*} (secci\'on~\ref{sec:csrf}) \\
    \texttt{zap-baseline} & \texttt{scripts/run-zap-baseline.sh}: escaneo din\'amico contra el backend real en Docker \\
    \bottomrule
  \end{tabular}
\end{table}

\texttt{ghcr-publish.yml} publica la imagen del backend en GHCR ante un
disparo manual o un \emph{push} de tag \texttt{v*}. \texttt{Makefile}
expone un objetivo \texttt{make all} que encadena las verificaciones
locales equivalentes a los \emph{jobs} de CI, para reproducibilidad fuera
de GitHub Actions.

\section{GHCR (imagen de contenedor del backend)}

La imagen del backend est\'a publicada realmente en GitHub Container
Registry:

\begin{table}[H]
  \centering
  \caption{Publicaci\'on GHCR (fuente: \texttt{docs/publicacion/PAQUETE-V1.0.0.md}).}
  \label{tab:ghcr-final}
  \begin{tabular}{@{}ll@{}}
    \toprule
    Campo & Valor \\
    \midrule
    Imagen & \texttt{ghcr.io/grinjoww/entregafinal-biopet-backend} \\
    Tag publicado & \texttt{sha-fe2f033} \\
    Digest (\texttt{sha256}, inmutable) & \texttt{sha256:ef1e857a95a307a1\allowbreak15ebe01599a41506eab824808b70a3c8e317dcc55bef5163} \\
    Commit asociado & \texttt{fe2f033138e3ec1fad07bf1038e10b8bb140449f} \\
    \bottomrule
  \end{tabular}
\end{table}

\begin{verbatim}
docker pull ghcr.io/grinjoww/entregafinal-biopet-backend@sha256:ef1e857a95a307a1
15ebe01599a41506eab824808b70a3c8e317dcc55bef5163
\end{verbatim}

El tag Git \texttt{v1.0.0} ya fue creado y publicado (apunta al commit
\texttt{0d5cd52}, estado hist\'orico de cierre de la Entrega Final; ver
secci\'on~\ref{sec:pendientes-honestos-publicacion}). Las etiquetas
\texttt{1.0.0} y \texttt{latest} en GHCR se publican autom\'aticamente
por el mismo workflow ante ese \emph{push} de tag; \texttt{sha-fe2f033}
sigue siendo, adem\'as, una etiqueta t\'ecnica real y verificable ligada
al commit espec\'ifico de la imagen publicada.

\section{Archivado permanente en Zenodo}

El software y el conjunto de datos de evidencias emp\'iricas de BIOPET
est\'an archivados en Zenodo con DOI reales, propagados a
\texttt{CITATION.cff}, \texttt{README.md} y \texttt{docs/checklists/fair.md}:

\begin{table}[H]
  \centering
  \caption{DOI publicados en Zenodo.}
  \label{tab:zenodo-final}
  \begin{tabular}{@{}ll@{}}
    \toprule
    Elemento & DOI \\
    \midrule
    Software & \texttt{\doisoftware} \\
    \emph{Dataset} de evidencias emp\'iricas & \texttt{\doidataset} \\
    \bottomrule
  \end{tabular}
\end{table}

Esto cierra la observaci\'on OBS-10 (archivado en Zenodo pendiente,
Tercera Entrega), registrada como CERRADA en
\texttt{docs/observaciones/OBSERVACIONES.md}.

\section{Estado de publicaci\'on del tag hist\'orico}
\label{sec:pendientes-honestos-publicacion}

El tag Git \texttt{v1.0.0} \textbf{ya existe y ya fue publicado}: apunta
al commit \texttt{0d5cd52}
(verificable con \texttt{git rev-list -n 1 v1.0.0}), que registra el
estado hist\'orico de cierre de la Entrega Final. Ese tag es inmutable:
las correcciones de retroalimentaci\'on y de recalificaci\'on posteriores
a su publicaci\'on (incluida la redacci\'on de este propio p\'arrafo) se
realizan sobre commits subsiguientes de la rama de trabajo, sin mover,
borrar ni recrear \texttt{v1.0.0}. A la fecha de cierre de este documento:

\begin{itemize}[nosep]
  \item El tag Git \texttt{v1.0.0} existe en el repositorio y apunta al
    commit \texttt{0d5cd52}.
  \item Al publicarse ese tag (\texttt{git push --tags}), el workflow
    \texttt{ghcr-publish.yml} publica autom\'aticamente las etiquetas
    \texttt{1.0.0} y \texttt{latest} en GHCR, sin intervenci\'on manual
    adicional.
  \item \texttt{CITATION.cff} declara \texttt{version: 1.0.0} junto con
    \texttt{date-released}, fijado con la fecha real de creaci\'on de ese
    tag (\texttt{git for-each-ref refs/tags/v1.0.0 --format=\%(creatordate:iso8601)}),
    en vez de quedar deliberadamente vac\'io como en versiones anteriores
    de este documento.
  \item La visibilidad p\'ublica del paquete GHCR (para que terceros puedan
    hacer \texttt{docker pull} sin autenticarse) qued\'o documentada como
    pendiente de confirmaci\'on expl\'icita en
    \texttt{docs/publicacion/PAQUETE-V1.0.0.md}.
\end{itemize}

Ninguno de estos pendientes afecta la validez de la evidencia emp\'irica ya
reportada en los cap\'itulos~\ref{cap:pruebas-calidad} y
\ref{cap:seguridad}: todas las mediciones se ejecutaron contra el sistema
real, ya sea en el entorno de desarrollo local o en el despliegue de
producci\'on ya activo.

% EVIDENCIA-FRED-BD-REPRODUCIBLE
% Ruta esperada: figuras/fred/Bd-reproducible-Flyaway.png
\evidencia{figuras/fred/Bd-reproducible-Flyaway.png}%
{Migraciones Flyway (V1--V6) aplicadas en orden desde una base vac\'ia.}%
{fig:bd-reproducible-final}

% EVIDENCIA-FRED-DOCKER-HEALTHY
% Ruta esperada: figuras/fred/Docker-compose-servicios-healthy.png
\evidencia{figuras/fred/Docker-compose-servicios-healthy.png}%
{Servicios de \texttt{docker compose} en estado \texttt{healthy} tras \texttt{make up}.}%
{fig:docker-healthy-final}

% EVIDENCIA-FRED-DOCKER-DIGEST
% Ruta esperada: figuras/fred/Docker-compose-diges-sha.png
\evidencia{figuras/fred/Docker-compose-diges-sha.png}%
{Im\'agenes de terceros fijadas por \emph{digest} \texttt{sha256} en \texttt{docker-compose.yml}.}%
{fig:docker-digest-final}

\section{Reproducibilidad de punta a punta}

\begin{itemize}[nosep]
  \item \textbf{Arranque completo:} \texttt{make up} (o \texttt{docker
    compose up --build -d}), sin pasos manuales adicionales.
  \item \textbf{Im\'agenes de terceros fijadas por \emph{digest}
    \texttt{sha256}} (no solo \emph{tag}), para evitar derivas silenciosas
    entre reconstrucciones.
  \item \textbf{Base de datos reproducible desde vac\'ia:} Flyway aplica
    \texttt{V1} a \texttt{V6} en orden en un checkout nuevo, verificado
    expl\'icitamente para esta entrega.
  \item \textbf{Configuraci\'on externalizada:} \texttt{.env.example}
    documenta todas las variables de entorno necesarias, sin secretos
    reales.
  \item \textbf{Migraciones verificadas desde base vac\'ia:} evidencia real
    de \texttt{V1} a \texttt{V6} aplicadas en orden por Flyway
    (figura~\ref{fig:bd-reproducible-final}).
  \item \textbf{Scripts de evidencia versionados y regenerables:} ver
    listado completo en el cap\'itulo~\ref{cap:metodologia},
    secci\'on~\ref{sec:diseno-evaluacion}.
\end{itemize}

\chapter{Matriz de trazabilidad}
\label{cap:trazabilidad-final}

La matriz cruda vive en \texttt{docs/trazabilidad/matriz.csv} (id de
requisito, tipo, prioridad MoSCoW, historia de usuario, caso de uso,
m\'odulo, endpoint, prueba automatizada, evidencia emp\'irica y estado),
validada autom\'aticamente por \texttt{scripts/validate-traceability.sh}
frente al SRS. Esta secci\'on selecciona las filas m\'as representativas de
las tres \'areas de trabajo del equipo, con el estado reportado tal como lo
declara el SRS/matriz vigente --- ning\'un estado se marca
\emph{verificado} sin una prueba automatizada o un documento de evidencia
citado.

\begin{landscape}
\small
\begin{longtable}{@{}
  >{\raggedright\arraybackslash}p{2.0cm}
  >{\raggedright\arraybackslash}p{1.8cm}
  >{\raggedright\arraybackslash}p{5.4cm}
  >{\raggedright\arraybackslash}p{4.6cm}
  >{\raggedright\arraybackslash}p{6.4cm}
  >{\raggedright\arraybackslash}p{2.2cm}
@{}}
\caption{Matriz de trazabilidad (selecci\'on representativa, v1.0.0). Matriz completa: \texttt{docs/trazabilidad/matriz.csv}.}
\label{tab:trazabilidad-final} \\
\toprule
Requisito & Resp. & Implementaci\'on & Componente & Prueba / evidencia & Estado \\
\midrule
\endfirsthead
\multicolumn{6}{c}{\tablename\ \thetable\ -- continuaci\'on} \\
\toprule
Requisito & Resp. & Implementaci\'on & Componente & Prueba / evidencia & Estado \\
\midrule
\endhead
\bottomrule
\endfoot

REQ-F-001 & Jaime & Registro de usuario (\texttt{ROLE\_DUENO}) & \rutalarga{AuthController/AuthService} & \texttt{AuthControllerTest}; Postman & Verificado \\
REQ-F-003 & Jaime & Login por cookies + JWT & \rutalarga{AuthController/JwtService} & \texttt{AuthControllerTest}, \texttt{JwtServiceTest}; ADR-006 & Verificado \\
REQ-F-005 & Jaime & Logout con revocaci\'on Redis & \texttt{TokenBlacklistService} & \texttt{AuthControllerTest}; \rutalarga{A07-authentication.md} & Verificado \\
REQ-F-006 & Jaime & Autorizaci\'on por rol (\texttt{@PreAuthorize}) & \rutalarga{SecurityConfig/MascotaController} & \texttt{MascotaControllerTest}; \rutalarga{A01-access-control.md} & Verificado \\
REQ-F-007 & Zaida/Jaime & Perfil de usuario autenticado & \texttt{UsuarioController} & Postman (\rutalarga{docs/postman/}) & Verificado \\
REQ-F-008--012 & Equipo & CRUD completo de mascotas & \rutalarga{MascotaController/MascotaService} & \texttt{MascotaControllerTest}; Postman & Verificado \\
REQ-F-013 & Equipo & Historial cl\'inico cronol\'ogico (m\'odulo independiente) & \texttt{ConsultaController} (CRUD b\'asico existe; historial longitudinal no) & \rutalarga{docs/requisitos/SRS.md} & Pendiente \\
REQ-F-015 & Equipo & Citas veterinarias (CRUD b\'asico existe; calendario interactivo no) & \texttt{CitaController} & \rutalarga{docs/requisitos/SRS.md} & Pendiente \\
REQ-F-020 & Jaime & Auditor\'ia de operaciones (parcial: solo autenticaci\'on) & \texttt{AuthenticationAuditService} & \texttt{AuthenticationAuditServiceTest}; \rutalarga{A09-logging.md} & Parcial \\
REQ-F-021 & Fred & Resumen de mascotas por especie & \texttt{fn\_resumen\_mascotas\_por\_especie} & \texttt{ResumenEspeciesIntegrationTest} (Testcontainers) & Verificado \\
REQ-F-023 & Equipo (Unidad IV) & Gesti\'on administrativa de usuarios & \rutalarga{UsuarioController/UsuarioService} & \texttt{UsuarioControllerTest} (16 pruebas) & Verificado \\
REQ-F-024 & Equipo (Unidad IV) & Gesti\'on de vacunas & \rutalarga{VacunaController/VacunaService} & \texttt{VacunaControllerTest} (10 pruebas); Postman & Verificado \\
REQ-F-025 & Equipo (Unidad IV) & Consulta de informaci\'on externa de especies (cach\'e Redis) & \rutalarga{ExternalApiController/ExternalApiService} & Sin prueba automatizada; evidencia manual cach\'e fr\'ia/caliente & Implementado \\
REQ-NF-001 & Fred & Rendimiento del listado de mascotas ($\leq$200/500~ms p95) & \texttt{MascotaService} (\texttt{@Cacheable}) & k6 (10 corridas), error 0{,}0\,\% & Verificado \\
REQ-NF-002 & Jaime & Comunicaci\'on cifrada obligatoria (HTTPS/TLS) & \texttt{TomcatDualConnectorConfig} & Verificaci\'on en vivo, TLSv1.3 (desarrollo) & Verificado \\
REQ-NF-005 & Zaida & Interfaz responsiva / calidad web (Lighthouse) & Frontend Angular & Lighthouse (12 corridas, mobile+desktop), 4/4 cumplen & Verificado \\
REQ-NF-006 & Jaime & Formato uniforme de error (\texttt{ProblemDetail}) & \texttt{GlobalExceptionHandler} & Postman, RFC~7807 & Verificado \\
REQ-NF-008 & Jaime & Cifrado de contrase\~nas (BCrypt) & \texttt{SecurityConfig} (\texttt{PasswordEncoder}) & \texttt{AuthControllerTest}; \rutalarga{db/seed.sql} & Verificado \\
REQ-NF-009 & Jaime & Auditor\'ia estructurada \texttt{AUTH\_AUDIT} & \texttt{AuthenticationAuditService} & \texttt{AuthenticationAuditServiceTest}; \rutalarga{A09-logging.md} & Verificado \\
REQ-NF-010 & Jaime & Rate limiting de login & \texttt{LoginRateLimiterService} & \texttt{AuthControllerTest}; \rutalarga{A07-authentication.md} & Verificado \\
REQ-NF-012 & Fred & Cach\'e del listado de mascotas con TTL externo & Redis/Valkey, \texttt{application.yml} & TTL confirmado (\rutalarga{docs/mediciones/redis/}); sin hit ratio & Parcial \\
REQ-NF-013 & Fred/Jaime & Estrategia h\'ibrida de acceso a datos (JPA + procedimientos almacenados) & \rutalarga{db/procs/}, repositorios JPA & \texttt{ResumenEspeciesIntegrationTest}; ADR-007 & Verificado \\
\end{longtable}
\normalsize
\end{landscape}

Cobertura de contribuciones reales por integrante en esta matriz: Jaime
(seguridad, autenticaci\'on, autorizaci\'on, cobertura JaCoCo,
an\'alisis est\'atico/din\'amico de seguridad), Fred (persistencia
reproducible, procedimientos almacenados, cach\'e, rendimiento k6, GHCR,
despliegue), y Zaida (frontend, accesibilidad, requisitos, SUS,
Lighthouse, PRISMA). Ning\'un requisito se marca \emph{verificado} sin una
prueba automatizada o un documento de evidencia citado en su columna.

\chapter{Resultados}
\label{cap:resultados-final}

Este cap\'itulo integra, sin repetir el detalle metodol\'ogico ya
presentado, qu\'e se construy\'o, qu\'e se valid\'o, qu\'e m\'etricas se
alcanzaron y qu\'e evidencia respalda cada resultado.

\section{Qu\'e se construy\'o}

Un sistema web de gesti\'on veterinaria con autenticaci\'on y autorizaci\'on
por rol/propiedad, CRUD de mascotas, y los m\'odulos de citas, consultas y
vacunas, respaldados por seis rutinas almacenadas de PostgreSQL con acceso
JPA formal, desplegado realmente en producci\'on (Render) y publicado como
imagen de contenedor inmutable en GHCR.

\section{Qu\'e se valid\'o y con qu\'e resultado}

\begin{table}[H]
  \centering
  \small
  \caption{S\'intesis de resultados cuantitativos de BIOPET v1.0.0.}
  \label{tab:sintesis-resultados}
  \begin{tabular}{@{}p{3.4cm}p{5.4cm}p{4.0cm}@{}}
    \toprule
    Dimensi\'on & Resultado & Fuente \\
    \midrule
    Pruebas automatizadas & 205 pruebas, 0 fallos, 0 errores & \texttt{surefire-reports/} \\
    Cobertura (JaCoCo) & LINE 91{,}80\,\%, BRANCH 79{,}39\,\%, umbral 70\,\% & \rutalarga{docs/mediciones/jacoco/METRICS.md} \\
    Rendimiento (k6) & 10 corridas (5 fr\'io + 5 caliente), error 0{,}0\,\% & \texttt{docs/mediciones/perf/REPORT.md} \\
    Usabilidad (SUS) & Media 74{,}44/100, IC95\,\% $[63{,}33,85{,}56]$, n=18 & \texttt{docs/mediciones/sus/REPORT.md} \\
    Calidad web (Lighthouse) & 4/4 categor\'ias cumplen umbral (SEO 100/100, mobile+desktop) & \texttt{docs/mediciones/lighthouse/} \\
    Seguridad OWASP (manual) & 6 controles auditados, todos \texttt{PASS} & \texttt{docs/mediciones/sec/REPORT.md} \\
    Seguridad din\'amica (ZAP) & 0 alertas altas, 1 informativa & \rutalarga{docs/mediciones/sec/zap/README.md} \\
    Seguridad est\'atica (SpotBugs) & 66 hallazgos, 0 de tipo \texttt{SQL\_*}, 1 alto (CSRF, mitigado por arquitectura) & \rutalarga{docs/mediciones/sec/static-analysis/README.md} \\
    Auditor\'ia SQL din\'amico & 0 hallazgos en \texttt{db/procs/} & \texttt{scripts/audit-sql-dynamic.sh} \\
    CI & 6 jobs verdes & \texttt{.github/workflows/ci.yml} \\
    Despliegue & Frontend y backend en Render, HTTPS activo, \texttt{/actuator/health} \texttt{UP} & \texttt{docs/despliegue/DEPLOYMENT.md} \\
    Empaquetado & GHCR con digest inmutable; DOI Zenodo (software y dataset) & \rutalarga{docs/publicacion/PAQUETE-V1.0.0.md} \\
    \bottomrule
  \end{tabular}
\end{table}

\section{Qu\'e requisitos fueron satisfechos}

El n\'ucleo \emph{Must Have} (autenticaci\'on, autorizaci\'on, CRUD de
mascotas) est\'a completamente verificado. Los m\'odulos de Unidad~IV
(vacunas, integraci\'on de API externa) est\'an verificados o
implementados seg\'un el caso; citas y consultas tienen CRUD b\'asico
funcional pero sus extensiones originales (calendario interactivo,
historial cl\'inico longitudinal completo) permanecen pendientes, tal como
lo declara el SRS (cap\'itulo~\ref{cap:requisitos}). Los requisitos
\emph{Should}/\emph{Could} de facturaci\'on, reportes exportables,
integraci\'on IoT y recomendaciones por IA permanecen fuera del alcance
implementado.

\section{Evidencia que respalda cada resultado}

Cada cifra de la tabla~\ref{tab:sintesis-resultados} tiene una ruta de
archivo real y verificable en el repositorio como fuente, sin excepci\'on;
los cap\'itulos~\ref{cap:pruebas-calidad}, \ref{cap:seguridad} y
\ref{cap:despliegue} desarrollan el detalle completo de c\'omo se
recolect\'o cada una.

\section{Consistencia num\'erica frente a entregas previas}

Este informe corrige expl\'icitamente, respecto a la Tercera Entrega,
las cifras que cambiaron con evidencia real: 109~$\to$~\textbf{205} pruebas
del backend; umbral de cobertura 60\,\%~$\to$~\textbf{70\,\%} (LINE
95{,}87\,\%~$\to$~91{,}80\,\%, BRANCH 76{,}87\,\%~$\to$~79{,}39\,\%); SUS
sin ejecutar~$\to$~\textbf{ejecutado} ($n=18$, media 74{,}44); Lighthouse
sin ejecutar~$\to$~\textbf{ejecutado, mobile+desktop} (SEO
82/100~$\to$~100/100 tras la corrida del 2026-08-18); DOI pendiente~$\to$~\textbf{DOI publicado}; GHCR
pendiente~$\to$~\textbf{imagen publicada}; producci\'on pendiente~$\to$~\textbf{desplegada
y con \emph{healthcheck} \texttt{UP}}. Los documentos hist\'oricos de la
Tercera Entrega (\texttt{informe-entrega-3.tex}/\texttt{.pdf}) se conservan
sin modificar, como registro de esa evaluaci\'on espec\'ifica.

\chapter{Discusi\'on}
\label{cap:discusion}

\section{Cobertura y calidad de c\'odigo}

Elevar el umbral autom\'atico de JaCoCo de 60\,\% a 70\,\% y mantener LINE
en 91{,}80\,\% (muy por encima del umbral) sugiere que la suite de pruebas
creci\'o en proporci\'on al c\'odigo nuevo de la Unidad~IV, no solo en
n\'umero absoluto (109~$\to$~205 pruebas): el crecimiento de BRANCH
(76{,}87\,\%~$\to$~79{,}39\,\%) fue menor que el de LINE, lo que es
consistente con la naturaleza t\'ipica de c\'odigo de validaci\'on
condicional (por ejemplo, en \texttt{sp\_registrar\_consulta\_validada} y
sus equivalentes en Java) que a\~nade ramas dif\'iciles de cubrir en su
totalidad. La cobertura sigue midiendo ejecuci\'on, no correcci\'on: el
hecho de que 0 hallazgos \texttt{SQL\_*} coexistan con 205 pruebas en verde
es evidencia complementaria, no una prueba formal de ausencia de defectos.

\section{Rendimiento}

Las diez corridas de k6 mantienen el mismo patr\'on cualitativo observado
en la Tercera Entrega (latencia mediana de un solo d\'igito de
milisegundos en condiciones de cach\'e caliente, throughput cercano a
92~req/s), con tasa de error 0{,}0\,\% sostenida al duplicar el n\'umero de
corridas (3+3~$\to$~5+5). El an\'alisis de Wilcoxon con tama\~no de efecto
de Rosenthal muestra que el efecto de la cach\'e no es uniforme entre
corridas ($r$ desde $-0{,}07$ hasta $-0{,}84$), lo que sugiere que factores
del entorno de ejecuci\'on (por ejemplo, el estado de la JVM o del pool de
conexiones al inicio de cada corrida) influyen en la magnitud observable
del beneficio de cach\'e, sin que eso invalide la conclusi\'on cualitativa
de que la cach\'e reduce la latencia. Frente a Putri, Hadi \& Ramdani
(2017)~\cite{putri2017perftesting} (cap\'itulo~\ref{cap:trabajos-relacionados}),
BIOPET reporta el mismo tipo de evidencia de \emph{benchmarking} pero
complementada con otras cinco dimensiones de calidad, no en aislamiento.

\section{Seguridad}

El resultado combinado de las tres t\'ecnicas de auditor\'ia (manual,
din\'amica, est\'atica) es consistente entre s\'i: ning\'una detect\'o
patrones de inyecci\'on SQL explotable, y el \'unico hallazgo de severidad
alta (CSRF deshabilitado) tiene una mitigaci\'on arquitect\'onica expl\'icita
y documentada (secci\'on~\ref{sec:csrf}), no ignorada. Comparado con el
trabajo de Yang et al. sobre vulnerabilidades JWT~\cite{yang2026tokentimebomb},
BIOPET no reclama haber verificado la ausencia de las categor\'ias de
vulnerabilidad que ese estudio descubre mediante \emph{fuzzing} dirigido a
la librer\'ia \texttt{jjwt}: esa clase de an\'alisis queda expl\'icitamente
fuera del alcance de las tres t\'ecnicas usadas aqu\'i.

\section{Usabilidad}

Con $n=18$ y media 74{,}44/100, BIOPET se ubica sobre el umbral de
referencia de 68 puntos de Bangor, Kortum \& Miller~\cite{bangor2008sus},
en la categor\'ia ``Bueno''. El intervalo de confianza se estrech\'o
respecto a la muestra de $n=10$ de la Tercera Entrega, pero como se declara
en el cap\'itulo~\ref{cap:pruebas-calidad}, ampliar la muestra no resuelve
por s\'i solo el sesgo de reclutamiento por conveniencia. El participante
de menor puntaje (sin experiencia web previa) refuerza, de forma
consistente con la literatura de usabilidad, la relevancia de una fase de
orientaci\'on breve para usuarios de perfil similar en trabajo futuro.

\section{Calidad}

La combinaci\'on de cobertura alta, cero hallazgos SQL y auditor\'ia de
seguridad multim\'etodo en verde sugiere un nivel de calidad de c\'odigo
consistente con las buenas pr\'acticas esperadas en un proyecto de este
tipo, enmarcado bajo ISO/IEC~25010 (secci\'on~\ref{sec:iso25010}). El
hallazgo de SEO~82/100 (por debajo del umbral 90) detectado en la primera
corrida de Lighthouse (2026-08-01) se document\'o como \'area de mejora
concreta en vez de minimizarse, y qued\'o corregido en la corrida definitiva
del 2026-08-18 (SEO~100/100 en las 12 corridas, ambos perfiles): a la fecha
de cierre de este informe, las cuatro categor\'ias de Lighthouse cumplen su
umbral.

\section{Despliegue y reproducibilidad}

El despliegue real en Render, con \emph{healthcheck} en estado \texttt{UP}
y HTTPS activo, junto con la imagen inmutable en GHCR y el archivado
permanente en Zenodo, cierra la brecha entre ``el sistema funciona en mi
m\'aquina'' y ``el sistema es reproducible y accesible por terceros'' que
persist\'ia en la Tercera Entrega. La reproducibilidad de punta a punta
(\texttt{make up}, im\'agenes fijadas por \emph{digest}, migraciones desde
base vac\'ia verificadas) es, junto con la amplitud de evidencia
multim\'etrica, uno de los dos pilares que distinguen a BIOPET frente a
los trabajos relacionados revisados (cap\'itulo~\ref{cap:trabajos-relacionados}),
ninguno de los cuales documenta reproducibilidad de infraestructura con el
mismo nivel de detalle.

\section{Comparaci\'on con trabajos relacionados: l\'imites de la afirmaci\'on}

Ninguna comparaci\'on de este cap\'itulo implica superioridad absoluta de
BIOPET sobre Rodr\'iguez et al.~\cite{rodriguez2024oscrum}, Cede\~no Ochoa
et al.~\cite{cedeno2021veterinaria}, Putri et al.~\cite{putri2017perftesting}
o Yang et al.~\cite{yang2026tokentimebomb}: no existe una comparaci\'on
emp\'irica directa (mismos datos, mismo entorno) entre BIOPET y ninguno de
esos sistemas. La \'unica afirmaci\'on sostenible, con la evidencia
disponible, es que BIOPET reporta un conjunto \textbf{m\'as amplio} de
dimensiones de calidad medidas sobre un mismo artefacto que cualquiera de
los trabajos revisados individualmente, no que su calidad relativa en cada
dimensi\'on individual sea superior.

\section{Respuesta expl\'icita a las preguntas de investigaci\'on}
\label{sec:respuestas-rq}

\textbf{RQ1 --- \textquestiondown En qu\'e medida BIOPET satisface los
requisitos definidos para la Entrega Final?} El n\'ucleo \emph{Must Have}
(autenticaci\'on, autorizaci\'on por rol/propiedad, CRUD de mascotas) est\'a
completamente verificado, y los tres requisitos formalizados en la
Unidad~IV (\texttt{REQ-F-023} a \texttt{REQ-F-025}: gesti\'on
administrativa de usuarios, vacunas e integraci\'on de API externa) est\'an
verificados o implementados seg\'un el caso (cap\'itulo~\ref{cap:requisitos}). Un
subconjunto expl\'icito de requisitos \emph{Should}/\emph{Could} heredados
de la Entrega~1A (historial cl\'inico longitudinal completo, calendario
interactivo de citas, IoT, facturaci\'on digital, reportes exportables,
recomendaciones por IA) permanece \textbf{pendiente}, sin que este informe
lo declare cumplido. Respuesta: BIOPET satisface el n\'ucleo de requisitos
comprometido para esta entrega y una parte de los requisitos ampliados de
Unidad~IV; no satisface la totalidad del alcance funcional original de la
Entrega~1A.

\textbf{RQ2 --- \textquestiondown Qu\'e resultados presenta BIOPET en las
evaluaciones emp\'iricas ejecutadas?} 205 pruebas automatizadas (0
fallos/errores); cobertura JaCoCo 91{,}80\,\%~LINE/79{,}39\,\%~BRANCH sobre
un umbral autom\'atico de 70\,\%; diez corridas de rendimiento k6 con
0{,}0\,\% de tasa de error; usabilidad SUS con $n=18$, media 74{,}44/100
(categor\'ia ``Bueno''); calidad web Lighthouse con las cuatro categor\'ias
cumplidas en la corrida del 2026-08-18 (SEO~100/100); y seguridad auditada
por tres m\'etodos independientes (OWASP manual, ZAP din\'amico, SpotBugs
est\'atico), con 0 alertas de severidad alta y 0 hallazgos de inyecci\'on
SQL (cap\'itulos~\ref{cap:pruebas-calidad} y \ref{cap:seguridad}). Respuesta:
en las seis dimensiones evaluadas, BIOPET cumple los umbrales configurados
por el propio equipo, con las limitaciones de alcance documentadas en el
cap\'itulo~\ref{cap:amenazas-limitaciones-final} (entorno acad\'emico,
tama\~no muestral, mecanismo de sesi\'on de la corrida Lighthouse del
2026-08-18 sin documentar en prosa).

\textbf{RQ3 --- \textquestiondown El despliegue, CI, GHCR y archivado
permiten reproducir y verificar el artefacto?} S\'i, con evidencia real:
seis \emph{jobs} de CI en verde, despliegue activo en Render con
\emph{healthcheck} \texttt{UP}, imagen de contenedor publicada en GHCR con
referencia inmutable por \emph{digest}, y archivado permanente en Zenodo
con DOI real para el software y para el conjunto de datos de evidencias
(cap\'itulo~\ref{cap:despliegue}). La reproducibilidad tiene l\'imites
declarados: el tag Git \texttt{v1.0.0} y el GitHub Release a\'un no existen,
y el plan gratuito de PostgreSQL en Render no permite verificar hoy una
operaci\'on sostenida m\'as all\'a de 30 d\'ias
(secci\'on~\ref{sec:limitacion-30-dias}). Respuesta: el artefacto es
reproducible de punta a punta en el entorno de desarrollo y verificable en
producci\'on dentro de esas condiciones expl\'icitas.

\chapter{Amenazas a la validez y limitaciones}
\label{cap:amenazas-limitaciones-final}

Este cap\'itulo integra y actualiza, con la evidencia real de v1.0.0, el
borrador de Jaime (\texttt{docs/informe/borradores/jaime/metodologia-y-amenazas.md}),
sin minimizar ninguna amenaza.

\section{Validez interna}

\begin{itemize}[nosep]
  \item Todas las mediciones de rendimiento y seguridad locales se
    ejecutaron en una sola m\'aquina de desarrollo, sin control de
    interferencia de otros procesos del sistema operativo.
  \item El certificado TLS del entorno de desarrollo es autofirmado; en
    producci\'on, HTTPS lo gestiona Render autom\'aticamente, fuera del
    control directo del equipo.
  \item El hit ratio de cach\'e se puede medir con dos m\'etodos que
    producen cifras distintas (global vs. aislado por clave); se document\'o
    expl\'icitamente cu\'al se us\'o y por qu\'e.
  \item El an\'alisis Wilcoxon/Rosenthal muestra tama\~nos de efecto no
    uniformes entre corridas ($r$ de $-0{,}07$ a $-0{,}84$), lo que indica
    variabilidad de entorno no completamente controlada entre corridas
    consecutivas.
\end{itemize}

\section{Validez externa}

\begin{itemize}[nosep]
  \item Todas las mediciones locales se realizaron en un entorno
    acad\'emico; el \'unico entorno equivalente a producci\'on es el
    despliegue real en el plan gratuito de Render, que tiene sus propias
    limitaciones (secci\'on~\ref{sec:limitacion-30-dias}).
  \item La muestra SUS ($n=18$, reclutada por conveniencia) no es
    necesariamente representativa de usuarios finales de una cl\'inica
    veterinaria real; un tama\~no muestral mayor reduce el margen del
    intervalo de confianza pero no cambia el m\'etodo de reclutamiento.
  \item Lighthouse depende del entorno de ejecuci\'on (hardware, red,
    versi\'on del navegador); sus puntajes no son un valor absoluto
    reproducible en cualquier m\'aquina.
  \item Los usuarios de prueba del backend (cuentas acad\'emicas, admin
    sembrado) no son usuarios reales de una cl\'inica veterinaria.
\end{itemize}

\section{Validez de constructo}

\begin{itemize}[nosep]
  \item JaCoCo mide cobertura de ejecuci\'on, no ausencia de defectos.
  \item Lighthouse no representa toda la experiencia de usuario real; se
    complementa con SUS como medici\'on independiente, pero ninguna
    sustituye una evaluaci\'on de UX cualitativa dedicada.
  \item SUS mide percepci\'on de usabilidad, no eficacia cl\'inica ni
    operacional real en un contexto veterinario.
  \item ZAP y SpotBugs no prueban la ausencia total de vulnerabilidades,
    solo la ausencia de los patrones que sus reglas configuradas detectan.
  \item Las pruebas de integraci\'on con Testcontainers verifican
    comportamiento real de PostgreSQL para las rutinas nativas; la mayor\'ia
    del resto de la suite usa H2 en memoria, cuyo comportamiento no es
    id\'entico al de PostgreSQL real.
\end{itemize}

\section{Validez de conclusi\'on}

\begin{itemize}[nosep]
  \item Cada corrida de k6 ($\sim$3\,200 peticiones) corresponde a una
    ventana de tiempo corta ($\sim$30--35~s); no se midieron tendencias en
    periodos prolongados.
  \item El tama\~no muestral de SUS ($n=18$) sigue siendo insuficiente para
    an\'alisis de subgrupos demogr\'aficos con potencia estad\'istica
    adecuada.
  \item Ninguna medici\'on de este informe implica dise\~no experimental
    con grupos de control; ninguna mejora observada entre versiones se
    atribuye causalmente a un cambio espec\'ifico.
  \item Las 205 pruebas y la cobertura JaCoCo se ejecutan en un ambiente
    de prueba mixto (H2 + Testcontainers/PostgreSQL para las rutinas
    nativas), no en un \'unico motor de base de datos consistente en toda
    la suite.
\end{itemize}

\section{Limitaci\'on temporal expl\'icita: operaci\'on sostenida no verificable hoy}
\label{sec:limitacion-30-dias}

El plan gratuito de PostgreSQL en Render \textbf{elimina los datos a los 30
d\'ias} de creado el recurso, documentado expl\'icitamente en
\texttt{docs/despliegue/DEPLOYMENT.md} y en la pol\'itica de retenci\'on de
respaldos de \texttt{docs/despliegue/BACKUP.md} (retenci\'on de 30 d\'ias,
purgado autom\'atico de respaldos con m\'as de 31 d\'ias). Esto significa
que, a la fecha de cierre de este informe, \textbf{no es posible verificar
emp\'iricamente una condici\'on de operaci\'on sostenida del sistema en
producci\'on durante 30 d\'ias o m\'as}: el despliegue actual es real y
verificable (\emph{healthcheck} \texttt{UP}, HTTPS activo), pero su
horizonte temporal de evidencia continua est\'a acotado por la propia
pol\'itica del plan gratuito usado. Este informe documenta esta condici\'on
como una \textbf{limitaci\'on temporal genuina}, no como un requisito
cumplido ni como uno incumplido por negligencia: verificarla exigir\'ia
migrar a un plan pago de base de datos persistente (o instrumentar un
respaldo/restauraci\'on autom\'atico verificado en producci\'on real durante
ese periodo), decisi\'on que excede el alcance de esta entrega acad\'emica.

\section{Otras limitaciones del estudio}

\begin{itemize}[nosep]
  \item \textbf{Contexto exclusivamente acad\'emico:} ninguna medici\'on
    proviene de un despliegue operativo con usuarios reales de una
    cl\'inica veterinaria.
  \item \textbf{Sin evaluaci\'on industrial ni con desarrolladores
    profesionales externos.}
  \item \textbf{Sin comparaci\'on emp\'irica directa con sistemas de
    gesti\'on veterinaria alternativos:} toda comparaci\'on es contra
    umbrales t\'ecnicos propios o contra la literatura, no contra artefactos
    competidores evaluados en las mismas condiciones.
  \item \textbf{Lighthouse: mecanismo de sesi\'on no documentado en
    prosa.} La corrida real del 2026-08-18 (mobile+desktop, 12 corridas,
    4/4 categor\'ias cumplidas, incluido SEO~100/100) audit\'o
    \texttt{/mascotas} directamente (\texttt{finalUrl} sin redirecci\'on a
    \texttt{/login} en los seis reportes de esa ruta, a diferencia de la
    corrida de agosto); sin embargo, el equipo no document\'o
    expl\'icitamente en \texttt{docs/mediciones/lighthouse/README.md} qu\'e
    mecanismo permiti\'o evitar la redirecci\'on del \emph{guard} de
    autenticaci\'on. Este informe reporta el dato num\'erico real
    (verificado en los JSON crudos) sin afirmar un mecanismo de sesi\'on no
    documentado por el equipo. La narrativa completa del
    \texttt{README.md} de esa carpeta a\'un no fue actualizada para
    describir la corrida del 2026-08-18 en prosa.
  \item \textbf{Rate limiting de login no distribuido:} en memoria, por
    instancia \'unica del backend.
  \item \textbf{Sin integraci\'on con SIEM} para los logs de auditor\'ia.
  \item \textbf{Requisitos heredados de la Entrega~1A a\'un pendientes:}
    historial cl\'inico longitudinal completo, calendario interactivo de
    citas, IoT, facturaci\'on digital, reportes exportables y
    recomendaciones por IA (cap\'itulo~\ref{cap:requisitos}).
  \item \textbf{Bloque Jaime de PRISMA sin conteo de exclusiones:} los 10
    estudios incluidos (meta de $\geq$8--9 cumplida) est\'an documentados y
    el diagrama de flujo ya fue regenerado con los tres bloques
    cuantificados; el \'unico gap que persiste es que Jaime no document\'o
    su conteo de candidatos evaluados/excluidos, se\~nalado con l\'inea
    punteada en el propio diagrama (cap\'itulo~\ref{cap:trabajos-relacionados}).
  \item \textbf{Tag \texttt{v1.0.0} y GitHub Release todav\'ia no creados}
    (secci\'on~\ref{sec:pendientes-honestos-publicacion}), fuera del
    alcance expl\'icito de esta tarea de redacci\'on del informe.
\end{itemize}

Ninguna de estas limitaciones se presenta como resuelta; todas condicionan
expl\'icitamente el alcance de las conclusiones del
cap\'itulo~\ref{cap:conclusiones-final}.

\chapter{Conclusiones}
\label{cap:conclusiones-final}

\section{Principales logros}

\begin{enumerate}[nosep]
  \item Un sistema web funcional de gesti\'on veterinaria, con
    autenticaci\'on y autorizaci\'on por rol/propiedad, CRUD de mascotas, y
    los m\'odulos de citas, consultas y vacunas de la Unidad~IV,
    respaldados por seis rutinas almacenadas de PostgreSQL con acceso JPA
    formal.
  \item Evidencia emp\'irica multim\'etrica real: 205 pruebas
    automatizadas (0 fallos), cobertura JaCoCo 91{,}80\,\%/79{,}39\,\% sobre
    un umbral autom\'atico del 70\,\%, diez corridas de rendimiento con
    0{,}0\,\% de tasa de error, usabilidad SUS medida con 18 participantes
    reales, y auditor\'ia de seguridad por tres m\'etodos independientes.
  \item Despliegue real en producci\'on (Render, HTTPS activo,
    \emph{healthcheck} \texttt{UP}), imagen de contenedor inmutable
    publicada en GHCR, y archivado permanente en Zenodo con DOI reales
    para el software y el conjunto de datos de evidencias.
  \item Cierre completo (15/15) de las observaciones docentes formales
    acumuladas a lo largo de tres entregas previas.
\end{enumerate}

\section{Evidencia}

Cada afirmaci\'on cuantitativa de este informe cita su fuente real en el
repositorio (archivo, script o reporte espec\'ifico), reproducible por un
tercero con los comandos documentados en el
cap\'itulo~\ref{cap:despliegue}. Ninguna cifra se estim\'o ni se
extrapol\'o de una entrega anterior sin verificaci\'on directa contra la
evidencia vigente.

\section{Valor del proyecto}

M\'as all\'a de la funcionalidad implementada, el valor distintivo de
BIOPET, frente al panorama de trabajos relacionados revisado
(cap\'itulo~\ref{cap:trabajos-relacionados}), es la \textbf{amplitud
verificable de su evidencia emp\'irica} sobre un mismo artefacto: ning\'un
trabajo comparado reporta simult\'aneamente cobertura de pruebas,
rendimiento, usabilidad, calidad web y seguridad por tres \'angulos
distintos, todo enmarcado bajo ISO/IEC~25010 y con reproducibilidad de
infraestructura de punta a punta. Esto no implica superioridad sobre esos
trabajos: implica que la pr\'actica de documentaci\'on y evidencia aplicada
en BIOPET permite razonar sobre relaciones entre dimensiones de calidad
(por ejemplo, entre cobertura de pruebas y ausencia de hallazgos de
seguridad est\'atica) que un estudio de una sola dimensi\'on no puede
mostrar.

\section{Estado de release}

BIOPET \texttt{v1.0.0} representa el cierre funcional y de evidencia de la
Entrega Final, pero \textbf{no debe interpretarse como un sistema listo
para producci\'on industrial sin condiciones}: el despliegue actual usa el
plan gratuito de Render (con purgado de datos a los 30 d\'ias, sin
verificaci\'on posible hoy de operaci\'on sostenida m\'as all\'a de ese
horizonte), el rate limiting de login no es distribuido, no existe
integraci\'on con un SIEM, y varios requisitos \emph{Should}/\emph{Could}
heredados de la Entrega~1A permanecen pendientes. El tag Git \texttt{v1.0.0}
y el GitHub Release correspondiente todav\'ia no se han creado a la fecha
de cierre de este informe (fuera del alcance expl\'icito de esta tarea de
redacci\'on).

\section{Trabajo futuro}

\begin{itemize}[nosep]
  \item Completar el historial cl\'inico longitudinal, el calendario
    interactivo de citas, la facturaci\'on digital y los reportes
    exportables (\texttt{REQ-F-013}, \texttt{REQ-F-015}, \texttt{REQ-F-018},
    \texttt{REQ-F-019}).
  \item Redactar la narrativa completa del \texttt{README.md} de
    \texttt{docs/mediciones/lighthouse/} para la corrida mobile+desktop del
    2026-08-18 (ya archivada con JSON crudos y las cuatro categor\'ias
    cumplidas), documentando expl\'icitamente el mecanismo de sesi\'on que
    permiti\'o auditar \texttt{/mascotas} directamente en esa corrida.
  \item Medir rendimiento bajo carga espec\'ificamente para los endpoints
    de citas, consultas y vacunas, no solo para el listado de mascotas.
  \item Migrar a un plan de base de datos persistente (o instrumentar
    verificaci\'on de respaldo/restauraci\'on) para poder evaluar
    emp\'iricamente la operaci\'on sostenida m\'as all\'a de 30 d\'ias.
  \item Evaluar un mecanismo de rate limiting distribuido si el sistema
    llegara a desplegarse con m\'ultiples r\'eplicas del backend.
  \item Integrar los logs \texttt{AUTH\_AUDIT} con un sistema de
    recolecci\'on y alertamiento centralizado (SIEM).
  \item Documentar el conteo de candidatos evaluados y excluidos del
    bloque de Jaime en la revisi\'on PRISMA, para homogeneizar el nivel de
    detalle entre los tres bloques (cap\'itulo~\ref{cap:trabajos-relacionados}).
  \item Crear el tag Git \texttt{v1.0.0} y el GitHub Release
    correspondiente, lo que activar\'a autom\'aticamente la publicaci\'on
    de las etiquetas \texttt{1.0.0}/\texttt{latest} en GHCR.
\end{itemize}

\chapter{Declaraciones}
\label{cap:declaraciones}

Este cap\'itulo re\'une las nueve declaraciones obligatorias del informe,
cada una con secci\'on propia: contribuciones de autor\'ia, \'etica,
disponibilidad de datos, disponibilidad de c\'odigo, disponibilidad de
materiales, financiamiento, conflicto de intereses, uso de asistentes
autom\'aticos y agradecimientos.

\section{Contribuciones de autor\'ia (CRediT)}

La matriz completa de roles CRediT~\cite{credit-taxonomy}, verificada
individuo por individuo contra commits, ramas de trabajo y archivos
concretos del repositorio (\texttt{git shortlog -sne --all}), vive en
\texttt{CONTRIBUTORS.md}. La tabla~\ref{tab:credit-final} la resume:

\begin{table}[H]
  \centering
  \small
  \caption{Resumen de contribuciones CRediT por integrante (detalle completo en \texttt{CONTRIBUTORS.md}).}
  \label{tab:credit-final}
  \begin{tabular}{@{}p{3.6cm}p{4.4cm}p{5.4cm}@{}}
    \toprule
    Integrante & \'Areas principales & Evidencia \\
    \midrule
    Mariscal Cabrera, Jaime Josu\'e & Seguridad, autenticaci\'on/autorizaci\'on, cobertura JaCoCo, an\'alisis est\'atico/din\'amico de seguridad, ADR-006/007, CI & \rutalarga{Backend/src/main/java/com/biopet/security/}, \rutalarga{docs/mediciones/sec/}, \rutalarga{.github/workflows/} \\
    Beltr\'an Montiel, Fred Adri\'an & Procedimientos almacenados, acceso a datos, cach\'e, rendimiento k6, despliegue/GHCR & \rutalarga{db/procs/}, \rutalarga{docs/mediciones/perf/}, \rutalarga{docs/despliegue/} \\
    Taipe Mora, Zaida Melissa & Frontend Angular, accesibilidad, requisitos (SRS), SUS, Lighthouse, PRISMA & \rutalarga{frontend/src/app/}, \rutalarga{docs/requisitos/}, \rutalarga{docs/mediciones/sus/} \\
    \bottomrule
  \end{tabular}
\end{table}

Los tres autores de BIOPET, con afiliaci\'on institucional \'unica, son:
Fred Adri\'an Beltr\'an Montiel (\texttt{fbeltranm@uteq.edu.ec}), Jaime
Josu\'e Mariscal Cabrera (\texttt{jmariscalc@uteq.edu.ec}) y Zaida Melissa
Taipe Mora (\texttt{ztaipem@uteq.edu.ec}), todos de la Universidad
T\'ecnica Estatal de Quevedo. Otros colaboradores hist\'oricos que
participaron en fases anteriores del programa acad\'emico no son autores de
esta entrega y no se listan como tales en \texttt{CITATION.cff} ni en este
informe.

\section{\'Etica}

BIOPET utiliza tres categor\'ias de datos: datos sint\'eticos de
desarrollo/prueba; datos emp\'iricos de evaluaci\'on generados por
herramientas automatizadas contra el propio sistema
(\texttt{docs/etica/ETHICS.md}); y los 18 registros anonimizados
P01--P18 de la evaluaci\'on de usabilidad SUS, que seg\'un declaraci\'on
del equipo corresponden a participantes reales (detalle a continuaci\'on).
El sistema no
almacena datos cl\'inicos reales de mascotas ni ha sido operado con datos
de clientes reales de una cl\'inica veterinaria. La prueba de usabilidad
SUS se aplic\'o sobre 18 registros (P01--P18), matem\'aticamente
reproducibles desde sus respuestas Q1--Q10
(\texttt{scripts/analisis-sus.py}); seg\'un declaraci\'on del equipo,
corresponden a participantes reales. El protocolo declarado
(\texttt{docs/etica/ETHICS.md}) exig\'ia consentimiento informado firmado
por cada participante (plantilla en
\texttt{docs/etica/consentimientos/plantilla.md}), pero una auditor\'ia
documental posterior a esta entrega constat\'o que esos formularios
individuales firmados no est\'an actualmente disponibles como evidencia
verificable. Esta situaci\'on se document\'o, sin fabricar evidencia
retrospectiva, mediante una constancia de regularizaci\'on firmada por
los tres integrantes del proyecto
(\texttt{docs/etica/regularizacion-sus/CONSTANCIA-REGULARIZACION-SUS-BIOPET-2026-08-31.pdf}),
que \textbf{no sustituye} consentimientos individuales ni incorpora
firmas de participantes. Los resultados agregados se identifican
\'unicamente por c\'odigo de participante (P01--P18). Ning\'un documento
de este informe ni sus figuras
reproducen contrase\~nas con valor, JWT completos, valores de cookies,
encabezados \texttt{Authorization} con valor real, claves privadas ni el
identificador \texttt{jti} de un token real.

\section{Disponibilidad de datos}

El conjunto de datos de evidencias emp\'iricas (rendimiento, seguridad,
usabilidad, calidad web, cobertura) est\'a archivado permanentemente en
Zenodo con DOI \texttt{\doidataset}, y disponible tambi\'en en
\rutalarga{docs/mediciones/} del repositorio, con procedencia documentada en
\rutalarga{docs/mediciones/DATA-PROVENANCE.md} y diccionario de variables en
\rutalarga{docs/mediciones/DATA-DICTIONARY.md}, siguiendo los principios
FAIR (\rutalarga{docs/checklists/fair.md}).

\section{Disponibilidad de c\'odigo}

El c\'odigo fuente completo (backend, frontend, migraciones, scripts) est\'a
disponible en \texttt{https://github.com/Grinjoww/ENTREGA-FINAL-BIOPET}
bajo licencia MIT, y archivado permanentemente en Zenodo con DOI
\texttt{\doisoftware}. La imagen de contenedor del backend est\'a publicada
en GHCR con referencia inmutable por \emph{digest}
(cap\'itulo~\ref{cap:despliegue}).

\textbf{Reproducibilidad.} Todos los resultados cuantitativos de este
informe son reproducibles con los comandos documentados en el
cap\'itulo~\ref{cap:despliegue} y en \texttt{docs/informe/README.md},
contra el mismo commit citado en cada reporte de evidencia
(\texttt{docs/mediciones/*/README.md} y \texttt{*.md} de reporte
individuales).

\section{Disponibilidad de materiales}

Adem\'as del c\'odigo y los datos de medici\'on (secciones anteriores),
este trabajo produjo materiales complementarios --- protocolos,
instrumentos de evaluaci\'on y configuraciones de auditor\'ia --- que
tambi\'en quedan versionados en el mismo repositorio y archivo Zenodo
citados arriba (DOI \texttt{\doisoftware}), sin un identificador aparte
porque forman parte del mismo commit archivado:

\begin{itemize}
  \item \textbf{Instrumento y protocolo SUS}: cuestionario de Brooke
    (1996)~\cite{sus-brooke1996} sin modificar, plantilla de
    consentimiento informado (\rutalarga{docs/etica/consentimientos/plantilla.md})
    y protocolo \'etico (\rutalarga{docs/etica/ETHICS.md}).
  \item \textbf{Configuraciones de auditor\'ia de calidad web}:
    \texttt{lighthouserc.js}, \texttt{lighthouserc.desktop.js},
    \texttt{lighthouserc.render.js} y \texttt{lighthouserc.render.desktop.js}
    (perfiles m\'ovil/escritorio, contra \texttt{localhost} y contra el
    despliegue p\'ublico de Render).
  \item \textbf{Modelo de arquitectura C4}: fuente Structurizr DSL
    versionada en \rutalarga{docs/diagrams/workspace.dsl}.
  \item \textbf{Decisiones de arquitectura (ADR)} en \rutalarga{docs/adr/}
    y especificaci\'on de despliegue reproducible
    (\texttt{docker-compose.yml}, \texttt{render.yaml}, \texttt{Makefile}).
  \item \textbf{Requisitos}: SRS versionado en \rutalarga{docs/requisitos/SRS.md}.
\end{itemize}

Ninguno de estos materiales incluye credenciales, secretos ni datos
personales identificables de participantes (ver secci\'on de \'Etica).

\section{Financiamiento}

Este proyecto es exclusivamente acad\'emico, desarrollado como Entrega
Final de la asignatura \asignatura{} en la Universidad T\'ecnica Estatal
de Quevedo. No recibi\'o financiamiento externo de ninguna agencia,
empresa ni entidad de fomento a la investigaci\'on (\texttt{CONTRIBUTORS.md},
categor\'ia CRediT \emph{Funding acquisition}: no aplica para ning\'un
integrante).

\section{Conflicto de intereses}

Ninguno de los tres autores declara conflicto de intereses en relaci\'on
con este trabajo.

\section{Uso de asistentes autom\'aticos}

El uso de herramientas de Inteligencia Artificial durante el desarrollo de 
este proyecto se limitó exclusivamente a funciones de apoyo y orientaci\'on acad\'emica. 
Estas herramientas fueron utilizadas como gu\'ia para la organizaci\'on y 
distribuci\'on de tareas entre los integrantes del equipo, as\'i como para la 
revisi\'on de avances y la identificaci\'on de posibles mejoras durante el proceso
 de desarrollo. La Inteligencia Artificial no sustituy\'o el trabajo de los 
 integrantes ni la toma de decisiones del equipo, siendo los autores responsables 
 del dise\~no, implementaci\'on, validaci\'on y resultados finales del proyecto.


\section{Agradecimientos}

El equipo agradece al docente de la asignatura \asignatura{},
\docente{} (\correodocente{}), por la gu\'ia acad\'emica y la
retroalimentaci\'on de evaluaci\'on que orientaron las correcciones de
esta entrega. Se agradece tambi\'en a la \universidad{} y a la
\facultad{} por el acceso a la infraestructura acad\'emica utilizada
durante el desarrollo del proyecto. No se recibi\'o apoyo, patrocinio ni
colaboraci\'on de ninguna persona, instituci\'on o empresa externa al
equipo de tres autores y al docente aqu\'i mencionados.

\chapter{Anexos}
\label{cap:anexos-final}

\section{Glosario de siglas}

\small
\begin{longtable}{@{}ll@{}}
\caption{Siglas usadas en este informe.}
\label{tab:glosario-final} \\
    \toprule
    Sigla & Significado \\
    \midrule
\endfirsthead
    \multicolumn{2}{c}{\tablename\ \thetable\ -- continuaci\'on} \\
    \toprule
    Sigla & Significado \\
    \midrule
\endhead
    \bottomrule
\endfoot
    ADR & Architecture Decision Record \\
    API & Application Programming Interface \\
    ARIA & Accessible Rich Internet Applications \\
    CI & Continuous Integration \\
    CRUD & Create, Read, Update, Delete \\
    CSP & Content Security Policy \\
    CVE & Common Vulnerabilities and Exposures \\
    DOI & Digital Object Identifier \\
    DSR & Design Science Research \\
    FAIR & Findable, Accessible, Interoperable, Reusable \\
    GHCR & GitHub Container Registry \\
    GQM & Goal--Question--Metric \\
    HSTS & HTTP Strict Transport Security \\
    IC & Intervalo de confianza \\
    INCOSE & International Council on Systems Engineering \\
    ISO/IEC & International Organization for Standardization / International Electrotechnical Commission \\
    JWT & JSON Web Token \\
    MoSCoW & Must, Should, Could, Won't \\
    OWASP & Open Worldwide Application Security Project \\
    PRISMA & Preferred Reporting Items for Systematic Reviews and Meta-Analyses \\
    RBAC & Role-Based Access Control \\
    RFC & Request for Comments \\
    SAST & Static Application Security Testing \\
    SEO & Search Engine Optimization \\
    SIEM & Security Information and Event Management \\
    SP & Stored Procedure \\
    SRS & Software Requirements Specification \\
    SUS & System Usability Scale \\
    TLS & Transport Layer Security \\
    UI & User Interface \\
    VU & Usuario virtual (\emph{Virtual User}, k6) \\
    WCAG & Web Content Accessibility Guidelines \\
    ZAP & Zed Attack Proxy \\
\end{longtable}
\normalsize

\section{Endpoints principales de la API}

\begin{table}[H]
  \centering
  \small
  \caption{Endpoints reales expuestos por el backend v1.0.0.}
  \label{tab:endpoints-final}
  \begin{tabular}{@{}llp{7.0cm}@{}}
    \toprule
    M\'etodo & Ruta & Autenticaci\'on / rol \\
    \midrule
    POST & /api/auth/registro & P\'ublico (crea siempre \texttt{ROLE\_DUENO}) \\
    POST & /api/auth/login & P\'ublico \\
    POST & /api/auth/refresh & Cookie \texttt{refresh\_token} \\
    POST & /api/auth/logout & Cookie (idempotente sin sesi\'on) \\
    GET & /api/usuarios/me & Cualquier rol autenticado \\
    GET/POST/PUT/DELETE & /api/mascotas[/\{id\}] & Rol y propiedad (secci\'on \ref{cap:seguridad}) \\
    GET & /api/mascotas/resumen-especies & Cualquier rol autenticado \\
    \ /* & /api/citas/* & \texttt{CitaController}, rol y propiedad \\
    \ /* & /api/consultas/* & \texttt{ConsultaController}, \texttt{VETERINARIO}/\texttt{ADMIN} \\
    \ /* & /api/vacunas/* & \texttt{VacunaController}, rol y propiedad \\
    \ /* & /api/externa/* & \texttt{ExternalApiController} \\
    GET & /actuator/health & P\'ublico \\
    GET & /api/docs, /swagger-ui.html & P\'ublico (OpenAPI 3.0, Springdoc) \\
    \bottomrule
  \end{tabular}
\end{table}

\section{Documentos de referencia en el repositorio}

\begin{itemize}[nosep]
  \item ADR: \texttt{docs/adr/ADR-002} a \texttt{ADR-007}.
  \item Diagramas C4: \texttt{docs/diagrams/} (contenedores, componentes del backend).
  \item Requisitos: \texttt{docs/requisitos/} (SRS, historias, casos de uso, \texttt{CHANGELOG-REQ.md}).
  \item Trazabilidad: \texttt{docs/trazabilidad/matriz.csv}.
  \item Pruebas y cobertura: \texttt{docs/mediciones/jacoco/}.
  \item Rendimiento: \texttt{docs/mediciones/perf/}.
  \item Usabilidad: \texttt{docs/mediciones/sus/}.
  \item Calidad web: \texttt{docs/mediciones/lighthouse/}.
  \item Seguridad: \texttt{docs/mediciones/sec/} (OWASP manual, ZAP, an\'alisis est\'atico).
  \item Procedencia y diccionario de datos: \rutalarga{docs/mediciones/DATA-PROVENANCE.md}, \rutalarga{docs/mediciones/DATA-DICTIONARY.md}.
  \item Base de datos: \texttt{docs/basedatos/CATALOGOSP.md}.
  \item Despliegue: \texttt{docs/despliegue/} (\texttt{DEPLOYMENT.md}, \texttt{RUNBOOK.md}, \texttt{BACKUP.md}).
  \item Publicaci\'on: \texttt{docs/publicacion/PAQUETE-V1.0.0.md}.
  \item \'Etica: \texttt{docs/etica/ETHICS.md}.
  \item Checklists de rigor metodol\'ogico: \rutalarga{docs/checklists/} (\texttt{ralph2021-engineering-research.md}, \texttt{prisma2020.md}, \texttt{incose2023-req.md}, \texttt{fair.md}).
  \item Bit\'acora de observaciones: \texttt{docs/observaciones/OBSERVACIONES.md}.
  \item Contribuciones CRediT: \texttt{CONTRIBUTORS.md}.
  \item Metadatos de citaci\'on: \texttt{CITATION.cff}.
\end{itemize}

\section{Clases de prueba automatizadas citadas en este informe}

\small
\begin{itemize}[nosep]
  \item \texttt{AuthControllerTest}, \texttt{JwtCookieAuthenticationTest},
    \texttt{JwtServiceTest}, \texttt{LoginRateLimiterServiceTest},
    \texttt{AuthenticationAuditServiceTest} (seguridad y autenticaci\'on).
  \item \texttt{MascotaControllerTest} (autorizaci\'on por rol y propiedad).
  \item \texttt{SqlInjectionSecurityTest}, \texttt{SecurityHeadersTest}
    (OWASP A03 y A05).
  \item \texttt{ResumenEspeciesIntegrationTest} (Testcontainers, PostgreSQL real).
\end{itemize}
\normalsize

\section{Anexo A --- Bit\'acora de observaciones docentes}
\label{sec:anexo-observaciones}

Las 15 observaciones formales recibidas a lo largo de las Entregas~1A,
1B y Tercera Entrega est\'an \textbf{cerradas en su totalidad (15/15)},
con evidencia de cierre verificable en
\texttt{docs/observaciones/OBSERVACIONES.md}.

\begin{longtable}{@{}p{1.6cm}p{1.6cm}p{8.4cm}p{2.6cm}@{}}
\caption{Observaciones docentes acumuladas y su estado de cierre.}
\label{tab:anexo-observaciones} \\
\toprule
C\'odigo & Entrega & Observaci\'on (resumen) & Estado \\
\midrule
\endfirsthead
\multicolumn{4}{c}{\tablename\ \thetable\ -- continuaci\'on} \\
\toprule
C\'odigo & Entrega & Observaci\'on (resumen) & Estado \\
\midrule
\endhead
\bottomrule
\endfoot
OBS-01 & 1A & Repositorio no proporcionado o inaccesible & Cerrada \\
OBS-02 & 1A & Ausencia de RF-07 en la lista consolidada de requisitos & Cerrada \\
OBS-03 & 1A & RF-WEB remapeados sin matriz de trazabilidad expl\'icita & Cerrada \\
OBS-04 & 1A & Ambig\"uedad leve en RF-10 (recomendaciones informativas) & Cerrada \\
OBS-05 & 1B & DER entregado como \texttt{.dot}, no como exportaci\'on PNG de pgAdmin & Cerrada \\
OBS-06 & 1B & Colecci\'on Postman no versionada & Cerrada \\
OBS-07 & 1B & Workflow CI fuera de \texttt{.github/workflows/} & Cerrada \\
OBS-08 & 1B & Tag \texttt{v0.1.0-entrega-1b} exigido no creado & Cerrada \\
OBS-09 & 3 & Tag \texttt{v0.9.0-rc} exigido no creado & Cerrada \\
OBS-10 & 3 & Software no archivado en Zenodo, DOI pendiente & Cerrada \\
OBS-11 & 3 & Evidencia/reporte Lighthouse faltante & Cerrada \\
OBS-12 & 3 & Calidad del sistema no enmarcada en ISO/IEC~25010 & Cerrada \\
OBS-13 & 3 & \texttt{CONTRIBUTORS.md} sin roles CRediT individuales & Cerrada \\
OBS-14 & 3 & Afirmaci\'on falsa sobre ausencia de historial Git & Cerrada \\
OBS-15 & 3 & Correo institucional de Jaime en commits & Cerrada \\
\end{longtable}

\section{Anexo B --- Estrategia de b\'usqueda PRISMA}
\label{sec:anexo-prisma-busqueda}

Cadenas y criterios de b\'usqueda documentados por bloque en
\texttt{docs/checklists/prisma2020.md} (\'item~7):

\begin{itemize}[nosep]
  \item \textbf{Zaida} (IEEE Xplore, ACM Digital Library, Scopus, Google
    Scholar; b\'usqueda 2026-08-17):
    \texttt{"requirements elicitation" AND "systematic literature review"}
    \textbullet\ \texttt{"system usability scale" AND web application}
    \textbullet\ \texttt{veterinary management system web application case
    study}.
  \item \textbf{Jaime}: verificaci\'on dirigida por tema (metodolog\'ia
    Engineering Research/DSR/GQM, calidad ISO/IEC~25010, seguridad
    JWT/OWASP, rendimiento web) contra Crossref, sobre referencias ya
    identificadas por relevancia tem\'atica directa, sin cadena booleana
    \'unica documentada.
  \item \textbf{Fred} (2026-08-17): verificaci\'on uno a uno de 13
    candidatos por DOI (rango temporal expl\'icito 2020--2026), 2
    descartados (Yang et al. ICSE~2018 por fuera de rango; art\'iculo de
    TURCOMAT por HTTP~403).
\end{itemize}

\section{Anexo C --- Instrumento SUS aplicado}
\label{sec:anexo-sus}

Las diez preguntas originales de Brooke (1996)~\cite{sus-brooke1996}, sin
modificar la escala de 5 puntos, tal como se aplicaron
(\texttt{docs/mediciones/sus/instrumento-sus.md}):

\begin{enumerate}[nosep]
  \item Creo que usar\'ia este sistema con frecuencia.
  \item Encontr\'e el sistema innecesariamente complejo.
  \item Pens\'e que el sistema era f\'acil de usar.
  \item Creo que necesitar\'ia el apoyo de una persona con conocimientos
    t\'ecnicos para poder usar este sistema.
  \item Encontr\'e que las diversas funciones de este sistema estaban bien
    integradas.
  \item Pens\'e que hab\'ia demasiada inconsistencia en este sistema.
  \item Imagino que la mayor\'ia de las personas aprender\'ian a usar este
    sistema muy r\'apidamente.
  \item Encontr\'e el sistema muy engorroso (poco pr\'actico) de usar.
  \item Me sent\'i muy seguro/a usando el sistema.
  \item Necesit\'e aprender muchas cosas antes de poder usar este sistema.
\end{enumerate}

Escala: 1~=~Totalmente en desacuerdo, \ldots, 5~=~Totalmente de acuerdo.
C\'alculo (Brooke, 1996): \'items impares contribuyen (valor$-1$); \'items
pares contribuyen ($5-$valor); suma$\times 2.5$ = puntaje 0--100
(\texttt{scripts/analisis-sus.py}, funci\'on \texttt{sus\_score()}).

\section{Anexo D --- Evidencia de integraci\'on continua}
\label{sec:anexo-ci}

No se dispone de capturas de pantalla archivadas de corridas de CI en el
repositorio; por eso este anexo documenta la evidencia \textbf{real y
verificable de otra forma} (definici\'on de los \emph{jobs} en el propio
YAML, no una captura), en vez de inventar capturas inexistentes:

\begin{itemize}[nosep]
  \item Definici\'on de los 6 \emph{jobs} verificable en
    \rutalarga{.github/workflows/ci.yml}: \texttt{backend-test},
    \texttt{frontend-build}, \texttt{traceability}, \texttt{sql-audit},
    \texttt{security-static}, \texttt{zap-baseline}.
  \item Reproducci\'on local equivalente: \texttt{make all} (encadena las
    mismas verificaciones que corren en CI).
  \item Resultado esperado de cada \emph{job}, verificado localmente en
    esta fase: \texttt{backend-test} (205 pruebas, \texttt{jacoco:check}
    $\geq$70\,\%); \texttt{sql-audit} (0 hallazgos,
    \rutalarga{scripts/audit-sql-dynamic.sh}); \texttt{security-static}
    (0 hallazgos \texttt{SQL\_*}); \texttt{zap-baseline} (0 alertas altas).
\end{itemize}

\section{Anexo E --- Checklist Ralph et al. (Engineering Research)}
\label{sec:anexo-ralph}

Tabulaci\'on resumida del checklist completo
(\texttt{docs/checklists/ralph2021-engineering-research.md}, 15~\'items):

\begin{longtable}{@{}p{1.2cm}p{9.4cm}p{2.4cm}@{}}
\caption{Resumen del checklist Ralph et al. (2021), est\'andar Engineering Research.}
\label{tab:anexo-ralph} \\
\toprule
\'Item & Descripci\'on (par\'afrasis) & Estado \\
\midrule
\endfirsthead
\multicolumn{3}{c}{\tablename\ \thetable\ -- continuaci\'on} \\
\toprule
\'Item & Descripci\'on (par\'afrasis) & Estado \\
\midrule
\endhead
\bottomrule
\endfoot
E1 & Describe el artefacto propuesto con detalle adecuado & Cumple \\
E2 & Justifica la necesidad/utilidad del artefacto & Cumple \\
E3 & Eval\'ua conceptualmente el artefacto (fortalezas/debilidades/limitaciones) & Cumple \\
E4 & Eval\'ua emp\'iricamente con un m\'etodo reconocido (\emph{benchmarking} + encuesta de usabilidad) & Cumple \\
E5 & Indica claramente la metodolog\'ia emp\'irica usada & Cumple \\
E6 & Discute alternativas de estado del arte & Cumple parcialmente \\
E7 & Compara emp\'iricamente con alternativas o \emph{benchmarks} & Cumple parcialmente \\
E8 & Supuestos expl\'icitos, plausibles y consistentes & Cumple parcialmente \\
E9 & Notaci\'on consistente & Cumple \\
D1 & Materiales suplementarios (c\'odigo, datos) & Cumple \\
D2 & Justifica elementos faltantes del paquete de replicaci\'on & Cumple parcialmente \\
D3 & Discute la base te\'orica del artefacto & Cumple (ver cap\'itulo~\ref{cap:marco-teorico}) \\
D4 & Argumentos de correcci\'on formal (teoremas, demostraciones) & No aplica \\
D5 & Ejemplos ilustrativos del artefacto & Cumple \\
D6 & Eval\'ua en contexto relevante para la industria & No cumple (fuera de alcance de un PFC) \\
\end{longtable}

\section{Anexo F --- Matriz de trazabilidad completa}
\label{sec:anexo-matriz-completa}

Reproducci\'on \'integra de \texttt{docs/trazabilidad/matriz.csv}
(38~requisitos), validada por \texttt{scripts/validate-traceability.sh}.
Columnas de endpoint y evidencia emp\'irica cruda se omiten aqu\'i por
espacio; est\'an disponibles en el CSV original.

\begin{landscape}
\small
\begin{longtable}{@{}p{2.0cm}p{1.3cm}p{1.3cm}p{2.0cm}p{2.0cm}p{4.6cm}p{4.4cm}p{2.0cm}@{}}
\caption{Matriz de trazabilidad completa (fuente: \texttt{docs/trazabilidad/matriz.csv}).}
\label{tab:anexo-matriz-completa} \\
\toprule
Requisito & Tipo & MoSCoW & HU & CU & M\'odulo/c\'odigo & Prueba automatizada & Estado \\
\midrule
\endfirsthead
\multicolumn{8}{c}{\tablename\ \thetable\ -- continuaci\'on} \\
\toprule
Requisito & Tipo & MoSCoW & HU & CU & M\'odulo/c\'odigo & Prueba automatizada & Estado \\
\midrule
\endhead
\bottomrule
\endfoot
REQ-F-001 & funcional & Must & HU-001 & CU-01 & \rutalarga{AuthController/AuthService} & \texttt{AuthControllerTest} & verificado \\
REQ-F-002 & funcional & Must & HU-001 & CU-01 & \texttt{GlobalExceptionHandler} & \texttt{AuthControllerTest} & verificado \\
REQ-F-003 & funcional & Must & HU-002 & CU-02 & \rutalarga{AuthController/JwtService} & \texttt{AuthControllerTest}+\texttt{JwtServiceTest} & verificado \\
REQ-F-004 & funcional & Must & HU-003 & CU-03 & \rutalarga{AuthController/AuthService} & \texttt{AuthControllerTest} & verificado \\
REQ-F-005 & funcional & Must & HU-004 & CU-04 & \rutalarga{AuthController/TokenBlacklistService} & \texttt{AuthControllerTest} & verificado \\
REQ-F-006 & funcional & Must & HU-005 & CU-05 & \rutalarga{SecurityConfig/MascotaController} & \texttt{MascotaControllerTest} & verificado \\
REQ-F-007 & funcional & Should & HU-006 & CU-06 & \rutalarga{UsuarioController/AuthService} & \texttt{UsuarioControllerTest} & verificado \\
REQ-F-008 & funcional & Must & HU-007 & CU-07 & \rutalarga{MascotaController/MascotaService} & \texttt{MascotaControllerTest} & verificado \\
REQ-F-009 & funcional & Must & HU-008 & CU-08 & \rutalarga{MascotaController/MascotaService} & \texttt{MascotaControllerTest} & verificado \\
REQ-F-010 & funcional & Must & HU-009 & CU-09 & \rutalarga{MascotaController/MascotaService} & \texttt{MascotaControllerTest} & verificado \\
REQ-F-011 & funcional & Must & HU-010 & CU-10 & \rutalarga{MascotaController/MascotaService} & \texttt{MascotaControllerTest} & verificado \\
REQ-F-012 & funcional & Must & HU-011 & CU-11 & \rutalarga{MascotaController/MascotaService} & \texttt{MascotaControllerTest} & verificado \\
REQ-F-013 & funcional & Should & HU-012 & CU-12 & \rutalarga{ConsultaController/ConsultaService} & \texttt{ConsultaControllerTest} (parcial) & implementado \\
REQ-F-014 & funcional & Could & HU-013 & CU-13 & --- & --- & pendiente \\
REQ-F-015 & funcional & Should & HU-014 & CU-14 & \rutalarga{CitaController/CitaService} & \texttt{CitaControllerTest} & implementado \\
REQ-F-016 & funcional & Could & HU-015 & CU-15 & --- & --- & pendiente \\
REQ-F-017 & funcional & Could & HU-016 & CU-16 & --- & --- & pendiente \\
REQ-F-018 & funcional & Should & HU-017 & CU-17 & --- & --- & pendiente \\
REQ-F-019 & funcional & Could & HU-018 & CU-18 & --- & --- & pendiente \\
REQ-F-020 & funcional & Should & HU-019 & CU-19 & --- & --- & pendiente \\
REQ-F-021 & funcional & Should & HU-020 & CU-20 & \rutalarga{MascotaController/MascotaRepository} & \texttt{ResumenEspeciesIntegrationTest} & verificado \\
REQ-F-022 & funcional & Could & HU-021 & CU-21 & --- & --- & pendiente \\
REQ-F-023 & funcional & Should & HU-022 & CU-22 & \rutalarga{UsuarioController/UsuarioService} & \texttt{UsuarioControllerTest} & verificado \\
REQ-F-024 & funcional & Should & HU-023 & CU-23 & \rutalarga{VacunaController/VacunaService} & \texttt{VacunaControllerTest} (parcial) & verificado \\
REQ-F-025 & funcional & Could & HU-024 & CU-24 & \rutalarga{ExternalApiController/ExternalApiService} & sin prueba dedicada & implementado \\
REQ-NF-001 & no funcional & Must & --- & --- & Cach\'e Redis (listado mascotas) & k6 v2.1.0 (10 corridas) & verificado \\
REQ-NF-002 & no funcional & Must & --- & --- & \rutalarga{SecurityConfig/TomcatDualConnectorConfig} & \texttt{SecurityHeadersTest} & verificado \\
REQ-NF-003 & no funcional & Must & --- & --- & \texttt{JwtService} & \texttt{JwtServiceTest} & verificado \\
REQ-NF-004 & no funcional & Must & --- & --- & \texttt{JwtService} & \texttt{JwtServiceTest} (6 pruebas) & verificado \\
REQ-NF-005 & no funcional & Must & --- & --- & Frontend Angular & Lighthouse (12 corridas, mobile+desktop) & verificado \\
REQ-NF-006 & no funcional & Should & --- & --- & Frontend Angular & Manual (pendiente) & pendiente \\
REQ-NF-007 & no funcional & Must & --- & --- & \texttt{Makefile}/\texttt{docker-compose.yml} & \emph{healthcheck} Docker Compose & implementado \\
REQ-NF-008 & no funcional & Must & --- & --- & \texttt{SecurityConfig} (BCrypt) & \texttt{AuthControllerTest} & verificado \\
REQ-NF-009 & no funcional & Must & --- & --- & Logging de autenticaci\'on & \texttt{AuthenticationAuditServiceTest} & verificado \\
REQ-NF-010 & no funcional & Must & --- & --- & \texttt{LoginRateLimiterService} & \texttt{LoginRateLimiterServiceTest} & verificado \\
REQ-NF-011 & no funcional & Must & --- & --- & Estructura de paquetes backend & Revisi\'on de c\'odigo & verificado \\
REQ-NF-012 & no funcional & Should & --- & --- & Cach\'e Redis (TTL configurable) & Inspecci\'on + medici\'on Redis & verificado \\
REQ-NF-013 & no funcional & Must & --- & --- & \texttt{MascotaRepository} (\texttt{fn\_resumen\_...}) & \texttt{ResumenEspeciesIntegrationTest} & verificado \\
\end{longtable}
\end{landscape}

\section{Anexo G --- Clasificaci\'on de referencias por nivel de impacto}
\label{sec:anexo-referencias-impacto}

Criterio aplicado: se clasifica como \textbf{alto impacto} \'unicamente
una referencia publicada en una revista o \emph{proceedings} con revisi\'on
por pares, de una editorial acad\'emica reconocida (IEEE, ACM, Springer,
Wiley, Taylor \& Francis, BMJ) con DOI verificado, o en un venue
ampliamente reconocido del \'area (por ejemplo NDSS en seguridad). No se
clasifican como alto impacto las normas t\'ecnicas/RFC/documentaci\'on
oficial de herramientas (que son referencias v\'alidas pero no
acad\'emicas), ni las revistas regionales sin indexaci\'on Scopus/Web of
Science verificable.

\begin{longtable}{@{}p{7.4cm}p{4.8cm}p{2.0cm}@{}}
\caption{Clasificaci\'on de las 41 referencias del informe.}
\label{tab:anexo-referencias-impacto} \\
\toprule
Categor\'ia & Ejemplos & Cantidad \\
\midrule
\endfirsthead
\multicolumn{3}{c}{\tablename\ \thetable\ -- continuaci\'on} \\
\toprule
Categor\'ia & Ejemplos & Cantidad \\
\midrule
\endhead
\bottomrule
\endfoot
Alto impacto (revista/\emph{proceedings} indexado, DOI, editorial reconocida) & BMJ (PRISMA), NDSS, ACM Trans.\ on Storage, IEEE, Springer LNNS/CCIS, IET Software, Int.\ J.\ HCI, JMIS & 16 \\
Acad\'emica verificada, indexaci\'on no confirmada (revista regional/peque\~na con DOI) & Journal of Computer Sciences Institute (x3), Jurnal Teknologi dan Sistem Komputer, Research J.\ of Education Science and Technology, J.\ of Software and Systems Development, J.\ Physics Conf.\ Series, Ecuadorian Science Journal & 7 \\
Norma t\'ecnica / RFC / est\'andar de industria (v\'alida como referencia normativa, no acad\'emica) & OWASP Top~10, RFC~7519/7807, ISO~9241-11, ISO/IEC~25010, WCAG~2.1, INCOSE, CRediT & 7 \\
Documentaci\'on oficial de herramienta ya usada en el proyecto & Spring Security/Boot, JaCoCo, PostgreSQL, Redis, k6, Lighthouse, C4 model & 7 \\
Est\'andar emp\'irico no revisado por pares formalmente (\emph{preprint} arXiv ampliamente adoptado) & Ralph et al. (2021), Empirical Standards & 1 \\
\midrule
\textbf{Total} & & \textbf{38--41*} \\
\end{longtable}

\noindent\footnotesize *La suma var\'ia seg\'un c\'omo se cuenten los
l\'imites entre categor\'ias (p.~ej. \texttt{sus-brooke1996} es un
cap\'itulo de libro Taylor \& Francis cl\'asico y ampliamente citado, se
cuenta en ``alto impacto'' pese a no ser un art\'iculo de revista); el
total de entradas \'unicas en \texttt{referencias.bib} es 41.\normalsize

\textbf{Resultado de la auditor\'ia (bloqueo declarado, no ocultado):} de
las 41 referencias del informe, \textbf{16 cumplen verificablemente el
criterio de alto impacto} definido arriba, por debajo del m\'inimo de 20
exigido por la Gu\'ia Final. No se reclasificaron revistas regionales o
peque\~nas como ``alto impacto'' \'unicamente para alcanzar el n\'umero
--- ver la nota en el cap\'itulo de declaraciones sobre este bloqueo.

\section{Evidencias pendientes de incorporaci\'on como figura}

La lista de capturas pendientes de copiar a \texttt{figuras/} (responsable,
n\'umero, ruta esperada, descripci\'on) se mantiene en
\texttt{docs/informe/README.md}, para no duplicarla en dos lugares del
repositorio. Las figuras ya incorporadas en la Tercera Entrega
(\texttt{figuras/jaime/}, \texttt{figuras/fred/}, \texttt{figuras/zaida/})
se reutilizan directamente en este informe mediante el mismo comando
\texttt{\textbackslash evidencia}.


% ===========================================================================
% BIBLIOGRAFIA (IEEE, BibTeX clasico) --- reutiliza referencias.bib
% ===========================================================================
\bibliographystyle{ieeetr}
\bibliography{referencias}

\end{document}
